\documentclass[11pt,a4paper]{article}
% head height 14pt: the italic Th ligature of a running head stands 13.01pt
\usepackage[margin=35mm,top=32mm,bottom=45mm,headheight=14pt]{geometry}
\usepackage{amsmath,amsthm,mathtools}
\usepackage{fontspec}
\usepackage[mathsf=sym]{unicode-math}
\setmainfont{EBGaramond}[Path=../fonts/, Extension=.otf,
  UprightFont=*-Regular, ItalicFont=*-Italic,
  BoldFont=*-Bold, BoldItalicFont=*-BoldItalic,
  FontFace={sb}{n}{*-SemiBold}, Numbers=Lining,
  SmallCapsFeatures={LetterSpace=4}]
% no sans text face in the family: \textsf (crate names) is code, in the mono
\defaultfontfeatures[Iosevka]{Path=../fonts/, Extension=.ttf,
  UprightFont=*-Regular, ItalicFont=*-Italic, BoldFont=*-Bold,
  FontFace={sb}{n}{*-SemiBold}, Scale=MatchLowercase}
\setsansfont{Iosevka}
\setmonofont{Iosevka}
% stylistic set 3: the calligraphic \mathcal (the default is a roundhand
% script). This copy maps capitals shared with Greek to the Latin, so copy
% and search see Latin (fonts/README-Garamond-Math.txt)
\setmathfont{Garamond-Math-Latin.otf}[Path=../fonts/, StylisticSet=3]
% the math sans is 6 % larger-eyed than the text; scaled to the text x-height
\setmathfont{Garamond-Math-Latin.otf}[Path=../fonts/,
  range=\mathsfup/{latin,Latin,num}, Scale=0.94]
% each \mathsf is one atom, as with a classic math alphabet, so a subscripted
% identifier such as \rivk takes a further subscript
\AtBeginDocument{\let\arbsymsf\mathsf\protected\def\mathsf#1{{\arbsymsf{#1}}}}
\usepackage{polyglossia}
\setmainlanguage[variant=british,ordinalmonthday=false]{english}
\PolyglossiaSetup{english}{frenchspacing=true}
\usepackage[protrusion=true,expansion=false]{microtype}
\usepackage{booktabs,enumitem}
\usepackage{placeins}
\usepackage{tikz}
\usetikzlibrary{arrows.meta,positioning,calc,fit,backgrounds}
\definecolor{tierspend}{RGB}{200,60,60}
\definecolor{tierfull}{RGB}{210,140,40}
\definecolor{tierin}{RGB}{60,130,180}
\definecolor{tierrecv}{RGB}{70,150,90}
\usepackage[psdextra,hidelinks]{hyperref}

\emergencystretch=3em
\hbadness=10000
\hfuzz=2pt

\newtheorem{theorem}{Theorem}[section]
\newtheorem{lemma}[theorem]{Lemma}
\newtheorem{proposition}[theorem]{Proposition}
\newtheorem{corollary}[theorem]{Corollary}
\theoremstyle{definition}
\newtheorem{definition}[theorem]{Definition}
\newtheorem{construction}[theorem]{Construction}
\newtheorem{assumption}[theorem]{Assumption}
\newtheorem{remark}[theorem]{Remark}
\newtheorem{example}[theorem]{Example}

% kind-coded theorem blocks: the rule weight codes the kind, in ink, no fill:
% proved statements 1.6pt, definitions and constructions 0.4pt, assumptions
% dashed 0.8pt; remarks and examples run unframed with an italic head
\usepackage{etoolbox,framed}
\makeatletter
\providecolor{tierin}{RGB}{60,130,180}
\providecolor{tierfull}{RGB}{210,140,40}
\providecolor{tierspend}{RGB}{200,60,60}
\providecolor{muted}{HTML}{5E5E5E}
% a left rule 0.9em from the text; framed spaces it like a list. Inside,
% \topsep is cleared so the statement's own list adds nothing, and
% FrameRestore's \sloppy is dropped so statements break like prose. A frame
% split across pages counts the last line's depth, which framed's
% \boxmaxdepth=0 turns into height
\setlength{\OuterFrameSep}{7pt plus 2pt minus 2pt}
\patchcmd{\fb@put@frame}{\splittopskip\z@}{\splittopskip\z@\splitmaxdepth\z@}{}{\PatchFailed}
% the frame sits flush left, as framed would centre it
\def\arb@flushleft#1{\hbox to\hsize{#1\hfil}}
\patchcmd{\fb@putboxa}{\centerline}{\arb@flushleft}{}{\PatchFailed}
\patchcmd{\fb@put@frame}{\centerline}{\arb@flushleft}{}{\PatchFailed}
\newcommand\arb@kind[2]{%
  \BeforeBeginEnvironment{#2}{\def\FrameCommand{#1}%
    \MakeFramed{\advance\hsize-\width\let\sloppy\relax\FrameRestore\topsep\z@}}%
  \AfterEndEnvironment{#2}{\endMakeFramed}}
\newcommand\arb@dashed[1]{\setbox\z@\hbox{#1}\hbox{\vbox to\ht\z@{\cleaders
  \vbox{\hrule\@width.8pt\@height3pt\vskip2pt}\vfill}\hspace{.9em}\box\z@}}
\arb@kind{\vrule\@width.4pt\hspace{.9em}}{definition}
\arb@kind{\vrule\@width.4pt\hspace{.9em}}{construction}
\arb@kind{\arb@dashed}{assumption}
\arb@kind{\vrule\@width1.6pt\hspace{.9em}}{theorem}
\arb@kind{\vrule\@width1.6pt\hspace{.9em}}{lemma}
\arb@kind{\vrule\@width1.6pt\hspace{.9em}}{proposition}
\arb@kind{\vrule\@width1.6pt\hspace{.9em}}{corollary}
\AtBeginEnvironment{remark}{\thm@headfont{\itshape}}
\AtBeginEnvironment{example}{\thm@headfont{\itshape}}
% a proof head breaks and hyphenates as the text it opens: amsthm sets it as a
% list label, one unbreakable box
\renewenvironment{proof}[1][\proofname]{\par\pushQED{\qed}\normalfont
  \topsep6\p@\@plus6\p@\relax\trivlist\item\relax{\itshape#1\@addpunct{.}}%
  \hskip\labelsep\ignorespaces}{\popQED\endtrivlist\@endpefalse}
% 11/14.5 pt; small (captions, tables, footnotes) 0.875; both size commands
% are robust, so the patches go to their inner macros. A 10pt class (the
% figure render's standalone) keeps its sizes
\ifnum\@ptsize=1
\expandafter\patchcmd\csname normalsize \endcsname{{13.6}}{{14.5}}{}{\PatchFailed}
\expandafter\patchcmd\csname small \endcsname{\@xpt\@xiipt}{{9.625}{12.7}}{}{\PatchFailed}
\fi
% the hyperref package wraps \@footnotetext; the size is set in the wrapped
% original
\patchcmd{\H@@footnotetext}{\footnotesize}{\small}{}{\PatchFailed}
\normalsize
\setlength{\parindent}{1.2em}
% operator spacing shrinks by 1mu at most, so tight lines keep their operators
\medmuskip=4mu plus 2mu minus 1mu
\thickmuskip=5mu plus 3mu minus 1mu
% headings on the 1.2 scale, semibold
\usepackage{titlesec}
\titleformat*{\section}{\fontsize{15.84}{19}\fontseries{sb}\selectfont}
\titleformat*{\subsection}{\fontsize{13.2}{16}\fontseries{sb}\selectfont}
% no heading below subsection: \subsubsection runs in, like \paragraph
\titleformat{\subsubsection}[runin]{\normalfont\bfseries}{\thesubsubsection}{0.5em}{}[.]
% a part (the complete edition) opens as a title page does: numeral, name,
% subtitle; its title reads "Name: subtitle"
\titleformat{\part}[display]{\raggedright}%
  {\fontsize{47.3}{50}\selectfont\thepart}{6mm}{\arb@parttitle}
\titlespacing*{\part}{0pt}{0pt}{12mm}
\def\arb@parttitle#1{\arb@partsplit#1\@nil}
\def\arb@partsplit#1: #2\@nil{{\fontsize{22.8}{27}\selectfont#1\par}%
  \vspace{3mm}{\fontsize{13.2}{17}\selectfont#2\par}}
% captions small, flush left, bold label; table captions above, wherever the
% source puts them
\usepackage{float}
\floatstyle{plaintop}
\restylefloat{table}
\usepackage[font=small,labelfont=bf,labelsep=period,
  justification=raggedright,singlelinecheck=false]{caption}
\captionsetup[table]{position=top}
% tables are small, set in \@floatboxreset, after the float's \normalsize
\def\arb@table{table}
\g@addto@macro\@floatboxreset{\ifx\@captype\arb@table\small\fi}
% a table's paragraph columns are set ragged, as narrow measures want; \\ still ends
% the row
\g@addto@macro\@arrayparboxrestore{\raggedright\let\\\tabularnewline}
% links: a conifer underline on screen, nothing in print
\hypersetup{pdfborder={0 0 0.6},pdfborderstyle={/S/U/W 0.6},
  allbordercolors={0.110 0.353 0.227},pdflinkmargin=0pt}
% contents: no link underlines; part and section entries in the heading
% weight, ragged right
\patchcmd{\l@part}{\bfseries}{\fontseries{sb}\selectfont}{}{\PatchFailed}
\patchcmd{\l@section}{\bfseries}{\fontseries{sb}\selectfont}{}{\PatchFailed}
\patchcmd{\l@part}{\rightskip\@pnumwidth}{\rightskip\@pnumwidth\@plus1fil}{}{\PatchFailed}
\patchcmd{\l@section}{\rightskip\@pnumwidth}{\rightskip\@pnumwidth\@plus1fil}{}{\PatchFailed}
\patchcmd{\l@part}{\hfil}{\hfill}{}{\PatchFailed}
\patchcmd{\l@section}{\hfil}{\hfill}{}{\PatchFailed}
\let\arb@tableofcontents\tableofcontents
\renewcommand{\tableofcontents}{{\hypersetup{pdfborderstyle={/S/U/W 0}}%
  \arb@tableofcontents}}
% abstract: a flush-left head like the contents', body-size text
\AtBeginDocument{\renewenvironment{abstract}{\section*{\abstractname}}{\par}}
% identity from the title line: name, subtitle, accession numeral
\providecommand{\accession}{}
\def\arb@name#1#2\@nil{\def\arbvolname{#2}}
\def\arb@parsetitle\textbf#1\\[#2]#3#4\@nil{\arb@name#1\@nil\def\arb@sub{#4}}
\providecommand{\arbvolname}{}
\renewcommand{\maketitle}{\expandafter\arb@parsetitle\@title\@nil
  \hypersetup{pdftitle={\arbvolname: \arb@sub},pdfauthor={Marek}}%
  \thispagestyle{empty}\begingroup\raggedright
  {\fontsize{13.2}{16}\selectfont\scshape The Zcash Arboretum\par}
  \vspace*{50mm}
  {\fontsize{47.3}{50}\selectfont\accession\par}
  \vspace{6mm}
  {\fontsize{22.8}{27}\selectfont\arbvolname\par}
  \vspace{3mm}
  {\fontsize{13.2}{17}\selectfont\arb@sub\par}
  \vspace{18mm}
  {\fontsize{13.2}{17}\selectfont Marek\par}
  {\small\@author\par}
  \endgroup\clearpage}
\usepackage{fancyhdr}
\pagestyle{fancy}
\fancyhf{}
\fancyhead[L]{\small\scshape\ifdefempty{\accession}{}{\accession\quad}\arbvolname}
\fancyhead[R]{\small\itshape\nouppercase{\leftmark}}
\fancyfoot[C]{\small\thepage}
\renewcommand{\headrulewidth}{0pt}
\makeatother
\renewcommand{\sectionmark}[1]{\markboth{\thesection\ \ #1}{}}

\newcommand{\F}{\mathbb{F}}
\newcommand{\Z}{\mathbb{Z}}
\newcommand{\N}{\mathbb{N}}
\newcommand{\G}{\mathbb{G}}
\newcommand{\E}{\mathbb{E}}
\newcommand{\PP}{\mathbb{P}}
\renewcommand{\Pr}{\mathrm{Pr}}
\newcommand{\Fp}{\F_p}
\newcommand{\Fq}{\F_q}
\newcommand{\ip}[2]{\langle #1,\, #2 \rangle}
\newcommand{\Comm}{\mathsf{Commit}}
\newcommand{\Open}{\mathsf{Open}}
\newcommand{\Setup}{\mathsf{Setup}}
\newcommand{\Verify}{\mathsf{Verify}}
\newcommand{\negl}{\mathsf{negl}}
\newcommand{\poly}{\mathsf{poly}}
\newcommand{\PPT}{\mathsf{PPT}}
\newcommand{\Adv}{\mathcal{A}}
\newcommand{\Prov}{\mathcal{P}}
\newcommand{\Ver}{\mathcal{V}}
\newcommand{\Sim}{\mathcal{S}}
\newcommand{\Ext}{\mathcal{E}}
\newcommand{\Rel}{\mathcal{R}}
\newcommand{\Lang}{\mathcal{L}}
\newcommand{\nfsym}{\mathsf{nf}}
\newcommand{\cmsym}{\mathsf{cm}}
\newcommand{\cmx}{\mathsf{cmx}}
\newcommand{\cvsym}{\mathsf{cv}}
\newcommand{\rk}{\mathsf{rk}}
\newcommand{\Pallas}{\ensuremath{\mathsf{Pallas}}}
\newcommand{\Vesta}{\ensuremath{\mathsf{Vesta}}}
\newcommand{\Ep}{\ensuremath{\mathcal{E}}}
\newcommand{\NoteCommit}{\ensuremath{\mathsf{NoteCommit}}}
\newcommand{\Sinsemilla}{\ensuremath{\mathsf{Sinsemilla}}}
\newcommand{\ShortCommit}{\ensuremath{\mathsf{ShortCommit}}}
\newcommand{\CommitIvk}{\ensuremath{\mathsf{Commit}^{\mathsf{ivk}}}}
\newcommand{\Extract}{\ensuremath{\mathsf{Extract}}}
\newcommand{\Acc}{\ensuremath{\mathsf{Acc}}}
\newcommand{\GroupHash}{\ensuremath{\mathsf{GroupHash}}}
\newcommand{\ask}{\ensuremath{\mathsf{ask}}}
\newcommand{\aksym}{\ensuremath{\mathsf{ak}}}
\newcommand{\nksym}{\ensuremath{\mathsf{nk}}}
\newcommand{\ivk}{\ensuremath{\mathsf{ivk}}}
\newcommand{\fvk}{\ensuremath{\mathsf{fvk}}}
\newcommand{\ovk}{\ensuremath{\mathsf{ovk}}}
\newcommand{\dk}{\ensuremath{\mathsf{dk}}}
\newcommand{\pkd}{\ensuremath{\mathsf{pk_d}}}
\newcommand{\gd}{\ensuremath{\mathsf{g_d}}}
\newcommand{\rcv}{\ensuremath{\mathsf{rcv}}}
\newcommand{\rcm}{\ensuremath{\mathsf{rcm}}}
\newcommand{\rivk}{\ensuremath{\mathsf{r_{ivk}}}}
\newcommand{\rsk}{\ensuremath{\mathsf{rsk}}}

\renewenvironment{abstract}{%
  \small\begin{center}{\bfseries\abstractname}\end{center}\medskip\noindent\ignorespaces
}{\par\medskip}

% Crosslink notation (one meaning per symbol; see architecture.md, Notation)
% chain kinds, genesis blocks and prefix orders; italic x ranges over the two
% kinds (no star placeholder)
\newcommand{\bc}{\mathrm{bc}}
\newcommand{\bft}{\mathrm{bft}}
\newcommand{\Origin}{\mathcal{O}}
\newcommand{\Obc}{\Origin_{\bc}}
\newcommand{\Obft}{\Origin_{\bft}}
\newcommand{\preceqbc}{\preceq_{\bc}}
\newcommand{\preceqbft}{\preceq_{\bft}}
\newcommand{\precbc}{\prec_{\bc}}
\newcommand{\precbft}{\prec_{\bft}}
\newcommand{\agrees}{\mathrel{\preceq\mkern-9mu\succeq}}
% C truncated by k: C\trunc{k}, or C\trunc[\bc]{k}; an ordinary atom, so a
% following relation keeps its space
\NewDocumentCommand{\trunc}{o m}{\mathord{\lceil}\IfValueT{#1}{_{#1}}^{#2}}
\newcommand{\lca}{\mathsf{lca}}
\newcommand{\score}{\mathsf{score}}
\newcommand{\work}{\mathsf{work}}
\newcommand{\ch}{\mathsf{ch}}
% the sub-protocols and the paper's protocols
\newcommand{\Piorigbc}{\Pi_{\mathrm{origbc}}}
\newcommand{\Piorigbft}{\Pi_{\mathrm{origbft}}}
\newcommand{\Pibc}{\Pi_{\bc}}
\newcommand{\Pibft}{\Pi_{\bft}}
\newcommand{\Pilc}{\Pi_{\mathrm{lc}}}
\newcommand{\Pisac}{\Pi_{\mathrm{sac}}}
% the paper's ledgers; the option is the node: \LOGfin[i]^t
\newcommand{\LOG}{\mathsf{LOG}}
\NewDocumentCommand{\LOGfin}{o}{\LOG_{\mathrm{fin}\IfValueT{#1}{,#1}}}
\NewDocumentCommand{\LOGda}{o}{\LOG_{\mathrm{da}\IfValueT{#1}{,#1}}}
\NewDocumentCommand{\LOGlc}{o}{\LOG_{\mathrm{lc}\IfValueT{#1}{,#1}}}
\NewDocumentCommand{\LOGba}{o}{\LOG_{\mathrm{ba}\IfValueT{#1}{,#1}}}
\newcommand{\sanitize}{\mathsf{sanitize}}
\newcommand{\flatten}{\mathsf{flatten}}
\newcommand{\Ch}{\mathsf{Ch}}
\newcommand{\GST}{\mathsf{GST}}
\newcommand{\GAT}{\mathsf{GAT}}
\newcommand{\Tconfirm}{T_{\mathrm{confirm}}}
\newcommand{\Env}{\mathcal{Z}}
% the BFT interface
\newcommand{\proofsym}{\mathsf{proof}}
\newcommand{\bftlastfinal}{\mathsf{bft\text{-}last\text{-}final}}
\newcommand{\origbftlastfinal}{\mathsf{origbft\text{-}last\text{-}final}}
% the construction
\newcommand{\isstalledblock}{\mathsf{is\text{-}stalled\text{-}block}}
\newcommand{\iscoinbaseonlyblock}{\mathsf{is\text{-}coinbase\text{-}only\text{-}block}}
\newcommand{\contextbft}{\mathsf{context\_bft}}
\newcommand{\headersbc}{\mathsf{headers\_bc}}
\newcommand{\snapshot}{\mathsf{snapshot}}
\newcommand{\LF}{\mathsf{LF}}
\newcommand{\candidate}{\mathsf{candidate}}
\newcommand{\finalitydepth}{\mathsf{finality\text{-}depth}}
% the ledgers: \fin_i^t, (\ba_\mu)_i^t
\newcommand{\fin}{\mathsf{fin}}
\newcommand{\ba}{\mathsf{ba}}
\newcommand{\updatefin}{\mathsf{update\text{-}fin}}
\newcommand{\Bckpt}{B_{\mathrm{ckpt}}}
% lower volumes and ZIPs: the Wallet Guide's confirmation policy, ZIP 237
\newcommand{\ctr}{c_{\mathrm{tr}}}
\newcommand{\cuntr}{c_{\mathrm{untr}}}
\newcommand{\removed}{\mathsf{removed}}

\title{\textbf{\Huge Crosslink Guide}\\[6pt]\large The Crosslink~2 construction: trailing finality over proof of work}
\author{m@rek.onl}
\date{}
\begin{document}
\maketitle
\thispagestyle{empty}

\begin{abstract}
This volume of \emph{The Zcash Arboretum} describes Crosslink~2, the
construction of the design book \emph{Zcash Trailing Finality Layer} of
N.~Wilcox and D.-E.~Hopwood, which adds a Byzantine-fault-tolerant finality
protocol to Zcash's proof-of-work best chain. It states the
availability-finality dilemma and exactly what the ebb-and-flow theorems of
Neu, Tas and Tse prove of snap-and-chat; constructs Crosslink~2 from a
best-chain protocol and a BFT protocol taken as black boxes with stated
interfaces, with its cross-link fields, every validity and honest-behaviour
rule, and Stalled Mode; and defines the finalised and bounded-available
chains. It proves that the finalised chain never rolls back and agrees across
all honest nodes and times whenever the best chain has Prefix Agreement at the
confirmation depth or the BFT protocol has Final Agreement, the second with a
corrected proof and a counterexample showing that the Linearity rule is
necessary for it; proves, for a node whose finalised chain is a prefix of its
best chain, the bounded-availability statement that the book leaves unproved,
which, for some best-chain protocol of the interface, requires the confirmation
depth not to exceed the finalisation gap bound, and leaves the other case open;
and
gives each liveness argument its exact status. It records which theorems survive
Shielded Labs' design, describes the finaliser set as far as it is designed,
states what ``final'' means to a client and how it differs from confirmation
depth, and collects the status of every goal and the open problems. No ZIP
specifies Crosslink, the design runs only on Shielded Labs' test network, and
every component of Crosslink is designed but unspecified or an open problem. The
volume assumes the \emph{Math Guide}, the \emph{Crypto Guide} and the
\emph{Consensus Guide}, and cites the \emph{Ironwood Guide}, the \emph{FlyClient
Guide} and the \emph{Wallet Guide}. It is non-normative: the Zcash Protocol
Specification and the Zcash Improvement Proposals are authoritative, and neither
specifies Crosslink.
\end{abstract}
\clearpage \tableofcontents \clearpage

\section{Introduction}\label{sec:intro}

This section fixes the subject of the volume and its status, the classes and
status words that every later claim carries, the sources and the rule by which
they are cited, and the notation, before any construction.

\subsection{Scope and status}\label{sec:intro-scope}

This volume of \emph{The Zcash Arboretum} describes Crosslink~2, the
construction of the design book \emph{Zcash Trailing Finality Layer} of
N.~Wilcox and D.-E.~Hopwood. Crosslink~2 is a hybrid of Zcash's proof-of-work
best chain and a \emph{Byzantine-fault-tolerant} (BFT) protocol: a protocol
whose participants decide by votes and which remains safe while fewer than a
fixed fraction of them deviate arbitrarily from it, a property that ``The
ebb-and-flow model'' (\S\ref{sec:model}) states formally. The book describes
itself as ``an early and incomplete protocol design proposal'' (chapter
``Status and Next Steps''), and its Crosslink~2 chapters are unchanged since
13 December 2024. The volume is non-normative: the Zcash Protocol
Specification and the Zcash Improvement Proposals (ZIPs) are authoritative
for every deployed rule, and no Crosslink rule is normative anywhere.

Every status statement of the volume is made as of 16 October 2026, in the
present tense. As of that date, no ZIP specifies Crosslink: the main branch
of the ZIP repository contains no Crosslink, trailing-finality or NU8
deployment ZIP. Draft Crosslink texts exist only in an open, unmerged draft
pull request to that repository, opened on 17 March 2026, in which every text
is marked ``ZIP: Unassigned'', and on branches of two forks of the
repository; they are not ZIPs. These drafts weaken or drop several rules of the
book, and the author of the pull request describes its texts as requiring
refinement or rewriting (3 June 2026); the volume does not use any of them as
design. The design runs only on
Shielded Labs' test network, relaunched on 1 October 2026: there is no
deployment on Mainnet or Testnet, no activation decision and no production
value of any parameter, and the small parameter values of the prototype that
runs the test network, which its authors describe as unverified, are not
given here.

The only analyses of the design are the following. The ebb-and-flow paper of
J.~Neu, E.~N.~Tas and D.~Tse proves its theorems for snap-and-chat, the
construction from which Crosslink descends, in a permissioned model
(\S\ref{sec:model}). The book gives its own proofs and arguments. One
unrefereed mechanism-design review examines a Shielded Labs design text of
11 December 2025, which predates the overview (\S\ref{sec:fs-review}). Four
partial formal models exist, one of which, an unmerged Lean formalisation,
leaves every Crosslink theorem unproved (\S\ref{sec:formal}). No security audit
firm has reviewed the design, and no paper on arXiv or in the IACR ePrint
archive analyses Crosslink.

The network upgrade NU7, whose deployment ZIP, ZIP 259, has status Draft,
does not include Crosslink; ZIP 218 (Draft), which ZIP 259 lists, sets a
25-second block target spacing, so every Crosslink parameter counted in
blocks covers a third of its earlier duration, and the volume gives no value
of any such parameter.
ZIP 218, ``Motivation'', states that it is complementary to, and does not
compete with, finality mechanisms such as Crosslink, and since 16 October 2026
its footnote on Crosslink cites the book.

The volume assumes the \emph{Math Guide}, the \emph{Crypto Guide} and the
\emph{Consensus Guide}, and cites the \emph{Ironwood Guide}, the
\emph{FlyClient Guide} and the \emph{Wallet Guide} for the objects they
construct, each by italic volume name and section title. From the \emph{Math
Guide}: probability and asymptotic notation (\S``Probability, asymptotics,
and computation'', with \S``Asymptotic notation'' and \S``The security
parameter''). From the \emph{Crypto Guide}: signature schemes and strong
unforgeability (\S``Syntax and security goal'', Definitions ``Digital
signature scheme'' and ``Strong unforgeability''). From the \emph{Consensus
Guide}: work, the best chain and confirmations (\S``Work, and the best
chain''), Proposition ``No finite count is final'' (\S``Confirmation
requirements''), rollback limits (\S``Rollback limits, coinbase
maturity, and finality''), and the boundary that leaves the policy of
treating a transaction as final to the layers above it (\S``Layer
boundary''). From the \emph{Ironwood Guide}: block chains and
height (\S``Scope and status'', Definition ``Block chain, transaction,
consensus rules'') and anchors (\S``Anchors''). From the \emph{FlyClient
Guide}: the epistemic classes (\S``Scope and status'', Definition
``Epistemic classes''), the block header (\S``Header-only verification'',
Definition ``Block header''), the persistence assumption (\S``The
variable-difficulty backbone model'') and the commitment hook of the block
header (\S``The block commitment hash'' and \S``The specified commitment and
its consumers''). From the \emph{Wallet Guide}: the confirmation policy
(\S``The confirmation policy''), anchor height (\S``Target and anchor
heights''), expiry (\S``Expiry (ZIP 203)'') and terminal states
(\S``Status tracking and terminal states'').

Shielded Labs' design departs from the book in several rules; ``The Shielded
Labs design'' (\S\ref{sec:sl}) states the departures and which theorems
survive them.

\subsection{Epistemic classes}\label{sec:intro-classes}

\begin{definition}[Epistemic classes]\label{def:classes}
A \emph{component} is a data structure, rule, parameter or check of a
protocol, or a hypothesis of a result about it. A component is
\begin{enumerate}[label=(\roman*)]
\item \emph{specified} when normative text of the protocol specification or
  of a ZIP fixes it, whatever the status of that ZIP;
\item \emph{designed but unspecified} when a published design with an
  analysis exists but no ZIP or section of the protocol specification fixes
  it for Zcash;
\item an \emph{open problem} when no design exists.
\end{enumerate}
Each result, assumption and classified remark carries a status line whose
status word is one of the following, used verbatim:
\begin{enumerate}[label=(\alph*)]
\item ``proved (ebb-and-flow paper)'', ``proved (TFL book)'' or ``proved (this
  volume)'', for an implication with a complete proof by the source named,
  followed by ``, of a designed-but-unspecified construction'' when its
  object is a component of Crosslink;
\item ``argued (TFL book)'', ``argued (Shielded Labs)'' or ``argued
  (ebb-and-flow paper)'', for a statement that the source named supports by a
  prose argument only;
\item ``machine-checked fragment'', for a statement that a machine checks on a
  formal model or a bounded instance, which covers part of a result;
\item ``open problem'', for a statement that is neither proved nor argued, or
  a component with no design;
\item ``assumption'', on every Assumption environment: a hypothesis that is
  neither argued nor open, on which the results that name it are
  conditional.
\end{enumerate}
\end{definition}

The three classes are those of the \emph{FlyClient Guide}, \S``Scope and
status'', Definition ``Epistemic classes''; the class designed but
unspecified has the sense of that volume, not the sense in which the
\emph{Ironwood Guide} uses the same words. A statement that a source makes
without proof or argument is recorded as that source's statement in a remark
and carries no status word. No Crosslink rule is specified: the class
specified applies in this volume only to the lower layers on which Crosslink
rests. A property proved of a designed-but-unspecified component establishes
an implication about the design and does not make the component specified.

Unqualified, ``Crosslink~2'' denotes the construction of the book. Shielded
Labs' design overview receives one section for its departures from the
construction, ``The Shielded Labs design'' (\S\ref{sec:sl}), which states
what it keeps, drops and adds, and which theorems survive; ``The finaliser
set'' (\S\ref{sec:fs}) describes its finaliser set. Shielded Labs' document on
finality semantics, the chapters of its design book that describe earlier
designs, and the draft texts on fork branches are not cited as design.

\subsection{Sources and conventions}\label{sec:intro-sources}

Every external source is cited by author, title and date. A source whose
title names an implementation, a function of code or the test network, or
whose only title is the name of a repository, is cited instead by author,
kind and date. The design book is attributed to its authors, not to an
organisation.

The design book is N.~Wilcox and D.-E.~Hopwood, \emph{Zcash Trailing
Finality Layer}, whose Crosslink~2 chapters are unchanged since 13 December
2024; it is called \emph{the book} below and is cited by chapter and section
title, verbatim. \emph{Crosslink~1} is the book's earlier construction. Its
ledgers survive in the chapter ``Security Analysis of Crosslink 2'' below
that chapter's marker ``Not updated for Crosslink 2'', and in the chapter
``Potential Changes to Crosslink'', which declares itself out of date. Only
the ledgers of the chapter ``The Crosslink 2 Construction'', the locally
finalised chain and the locally bounded-available chain (\S\ref{sec:ledgers}),
are those of Crosslink~2.

The model source is J.~Neu, E.~N.~Tas and D.~Tse, ``Ebb-and-Flow Protocols:
A Resolution of the Availability-Finality Dilemma'', arXiv 2009.04987, in its
third revision of 4 February 2021, which appeared at the IEEE Symposium on
Security and Privacy 2021. It is called \emph{the paper} below and is cited by
section, definition and theorem number. Its theorems are for snap-and-chat in
a permissioned model, and its statement for proof of work is informal
(\S\ref{sec:model-pow}).

Shielded Labs' design is cited from two documents: Shielded Labs,
\emph{Zcash Crosslink Design Overview}, draft 10 of 14 September 2026,
edited to 5 October 2026, called \emph{the overview} below and cited by
section title; and Shielded Labs, \emph{Rewards Distribution}, 22 September
2026. Shielded Labs' public \emph{Crosslink FAQ}, modified 30 July 2026, is
cited only for statements on rewards and the staking calendar
(\S\ref{sec:fs}).

The mechanism-design review is a review by N.~Della Penna of a Shielded Labs
Crosslink design text of 11 December 2025, which predates the overview, with
a companion document of formal proofs, both dated January 2026, published on
5 February 2026 and unrefereed. Its title names an implementation, so it is
cited by author, kind and date; it is called \emph{the mechanism-design review}
below and serves as framing only (\S\ref{sec:fs-review}). Shielded Labs states
that it will seek assessments from security audit firms before activation; none
has been announced. By the same rule, the blog post of J.~McGee for Shielded
Labs of 29 May 2026 and the Quint models of Informal Systems and of Z.~Manian
are cited by author, kind and date, and the Lean pull request to the book's
repository and A.~Garcia's TLA+ specification by author, title and date
(\S\ref{sec:formal}).

ZIP text is that of the main branch of the ZIP repository on 16 October
2026. A ZIP is cited by its number and the title of a section; the protocol
specification is cited as protocol specification, \S``Title''. The lower
volumes predate that text: the \emph{Consensus Guide} fixes its boundary at
the draft NU7 rules of 6 October 2026, and later NU7 changes to ZIPs 218 and
237 are absent from the lower volumes. Where a statement of a lower volume
differs from the current text of a ZIP, this volume cites the ZIP.

Prose uses the spellings ``finaliser'', ``finalised'' and ``finalisation''.
Verbatim titles, rule names and quotations keep the spelling of their source,
as in the book's section ``Locally finalized chain''. The PCZT roles ``IO
Finalizer'' and ``Spend Finalizer'' of the \emph{Wallet Guide} are proper
names unrelated to the finalisers of this volume.

Each symbol has one meaning, and Table~\ref{tab:notation} indexes the symbols
with the subsection that defines each; the table defines nothing. The letter
$\sigma$ denotes only the confirmation depth of Crosslink~2
(\S\ref{sec:cons-params}), and no signature is written $\sigma$. The paper's
security parameter, which is also its confirmation time and which the paper
writes $\sigma$, is written $\lambda$, as in the \emph{Math Guide},
\S``The security parameter'', Definition ``Security parameter''; $\lambda$
never denotes a rate. The finalisation gap bound $L$ is an italic capital,
never calligraphic. The bounded-available depth $\mu$ keeps the book's
meaning and is defined locally (\S\ref{sec:led-ba}). A variable $x$ over the
two chain kinds, $\bc$ and $\bft$, replaces the book's star placeholder,
since the series reserves the star for the encoding of curve points. The
$(c, L)$-adversary of the \emph{FlyClient Guide} does not occur. The word
``anchor'' means only a Zcash anchor (\emph{Ironwood Guide},
\S``Anchors''). The word ``epoch'' is always qualified, as in BFT epoch. The
word ``final'', applied to a best-chain block or a transaction, is defined
once, in ``Finality as a status'' (\S\ref{sec:k-final}); it is distinct from
the final treestate of a block (\emph{Ironwood Guide}, \S``Anchors'',
Definition ``Treestate and anchor'') and from the finality of a bft-block in
a context (\S\ref{sec:iface-bft}).

% [sec:intro-sources] no worked numbers: the section cites no compute/ output.
\begin{table}[tbp]
\centering
\footnotesize \renewcommand{\arraystretch}{1.1}
\begin{tabular}{@{}p{0.29\linewidth}p{0.52\linewidth}p{0.11\linewidth}@{}}
\toprule
Symbol & Meaning & Defined in \\
\midrule
$\lambda$
  & security parameter and confirmation time of the paper, which the paper
  writes $\sigma$; never a rate
  & \S\ref{sec:model-execution} \\
$\beta$, $\Delta$, $\GST$, $\GAT$
  & adversarial fraction, delay bound and transitions of the paper; the
  paper's node and corruption counts have no symbol
  & \S\ref{sec:model-execution} \\
$\beta_1$, $\beta_2$
  & fractions tolerated under P1 and P2 by an ebb-and-flow protocol
  & \S\ref{sec:model-security} \\
$(\Adv_k(\beta), \Env_k)$, $\Adv_k^*$
  & environment of property P$k$ of the paper, $k \in \{1, 2\}$, and its
  worst case
  & \S\ref{sec:model-execution} \\
$\LOG$, $\LOG_i^t$, $\Tconfirm$
  & a ledger, the ledger of node $i$ at time $t$, and the transaction
  confirmation time
  & \S\ref{sec:model-security} \\
$\LOGfin$
  & finalised ledger: of the paper's protocols, and, for Crosslink~2, the
  transaction sequence of $\fin$
  & \S\ref{sec:model-security}, \S\ref{sec:led-ba} \\
$\LOGda$
  & dynamically available ledger of the paper
  & \S\ref{sec:model-security} \\
$\Pilc$, $\Pisac$, $\LOGlc$
  & longest-chain protocol of snap-and-chat, snap-and-chat itself, and the
  ledger of the longest-chain protocol
  & \S\ref{sec:model-sac} \\
$\Piorigbc$, $\Piorigbft$, $\Pibc$, $\Pibft$
  & the two black-box protocols and their Crosslink~2 transforms
  & \S\ref{sec:iface-hybrid} \\
$V_i^t$
  & a view: the value of node $i$ at time $t$
  & \S\ref{sec:iface-chains} \\
$x$
  & variable over the chain kinds $\bc$ and $\bft$, in place of the book's
  star placeholder
  & \S\ref{sec:iface-chains} \\
$\Origin_x$
  & genesis block of the $x$-chains
  & \S\ref{sec:iface-chains} \\
$C\trunc{k}$
  & $C$ without its last $k$ blocks, floored at genesis (the book's
  notation)
  & \S\ref{sec:iface-chains} \\
$B \preceq_x C$
  & $B$ is a prefix of $C$, equality included
  & \S\ref{sec:iface-chains} \\
$B \agrees C$
  & agreement: one of $B$ and $C$ is a prefix of the other; not transitive
  & \S\ref{sec:iface-chains} \\
$\lca$
  & last common ancestor
  & \S\ref{sec:iface-chains} \\
$\mathsf{height}$
  & number of predecessors (\emph{Ironwood Guide})
  & \S\ref{sec:iface-chains} \\
\bottomrule
\end{tabular}
\caption{Notation: each symbol with its one meaning and the subsection that
defines it. The table defines nothing; a letter bound inside a single
definition, construction or proof keeps its meaning only there.}
\label{tab:notation}
\end{table}

\begin{table}[tbp]
\ContinuedFloat
\centering
\footnotesize \renewcommand{\arraystretch}{1.1}
\begin{tabular}{@{}p{0.29\linewidth}p{0.52\linewidth}p{0.11\linewidth}@{}}
\toprule
Symbol & Meaning & Defined in \\
\midrule
$\ch_i^t$
  & best chain of node $i$ at time $t$
  & \S\ref{sec:iface-bc} \\
$\score$, $\work$
  & best-chain order, accumulated work for Zcash; work of a block
  & \S\ref{sec:iface-bc} \\
$\iscoinbaseonlyblock$
  & coinbase-only predicate
  & \S\ref{sec:iface-bc} \\
$k$
  & depth in Prefix Consistency and Prefix Agreement
  & \S\ref{sec:iface-bc} \\
$n$, $q$, $m$
  & number of voting units of a BFT epoch, the quorum $\lceil 2n/3 \rceil$,
  and the number of units cast non-honestly; $n$ has no other meaning
  & \S\ref{sec:iface-bft} \\
$\proofsym_P$
  & notarisation proof for a proposal $P$
  & \S\ref{sec:iface-bft} \\
$\bftlastfinal$, $\origbftlastfinal$
  & last final ancestor of a valid bft-block, in $\Pibft$ and in
  $\Piorigbft$
  & \S\ref{sec:iface-bft} \\
$\sigma$
  & bc-confirmation depth; never a signature
  & \S\ref{sec:cons-params} \\
$L$
  & finalisation gap bound
  & \S\ref{sec:cons-params} \\
$\isstalledblock$
  & stalled-block predicate
  & \S\ref{sec:cons-params} \\
$\contextbft$, $\headersbc$
  & the two cross-link fields
  & \S\ref{sec:cons-struct} \\
$\snapshot$, $\LF$, $\candidate$, $\finalitydepth$
  & functions derived from the cross-link fields
  & \S\ref{sec:cons-struct} \\
$\fin_i^t$, $\updatefin$
  & locally finalised chain of node $i$ at time $t$; its update rule
  & \S\ref{sec:led-fin} \\
$\mu$, $(\ba_\mu)_i^t$
  & bounded-available depth; locally bounded-available chain
  & \S\ref{sec:led-ba} \\
$\LOGba$
  & for Crosslink~2, where the book's names are used, the transaction
  sequence of $\ba_\mu$
  & \S\ref{sec:led-ba} \\
$\Bckpt$
  & checkpoint bft-block of a release
  & \S\ref{sec:led-sync} \\
$h_1$, $h_2$
  & roster height and start height of the Shielded Labs bootstrap; no value
  is given
  & \S\ref{sec:fs-roster} \\
$N$
  & active roster size of the Shielded Labs design; no value is given; the
  only meaning of italic $N$, the natural numbers being $\N$
  & \S\ref{sec:fs-roster} \\
$h_{\mathrm{exp}}$
  & expiry height of a non-coinbase transaction, $0$ disabling expiry
  & \S\ref{sec:fs-windows}, \S\ref{sec:k-wallet} \\
$\removed$
  & value removed from circulation at a height (ZIP 237)
  & \S\ref{sec:fs-rewards} \\
$\tau$
  & threshold fraction, $\tau \ge 2/3$, of certificates in the
  mechanism-design review
  & \S\ref{sec:fs-review} \\
$\ctr$, $\cuntr$, $z$
  & the confirmation policy of the \emph{Wallet Guide}
  & \S\ref{sec:k-wallet} \\
\bottomrule
\end{tabular}
\caption[]{Notation, continued.}
\end{table}

\section{The ebb-and-flow model}\label{sec:model}

This section states the model of the ebb-and-flow paper of Neu, Tas and Tse
(\S\ref{sec:intro-sources}), its security definitions, the snap-and-chat
construction and exactly what its theorems prove and for which protocols,
the paper's informal statement for proof of work, and four deficiencies of
snap-and-chat, among those the TFL book records, from which it departs. No
theorem of this section is about Crosslink~2 or about proof of work.

\subsection{The availability-finality dilemma}\label{sec:model-dilemma}

A \emph{ledger} is a finite sequence of transactions output by a node; a
\emph{ledger protocol} is a protocol in which every node outputs a ledger at
every time. Informally, a ledger protocol is safe when the ledgers of honest
nodes never conflict, one always being a prefix of the other, and live when
every transaction given to honest nodes eventually enters the ledger of every
honest node; Definition~\ref{def:ntt-security} makes both precise. A protocol
has \emph{dynamic participation} when honest nodes may stop and resume taking
part, asleep and awake, without notice to the other nodes. A \emph{network
partition} is an interval of time during which messages between two sets of
honest nodes are not delivered. A BFT protocol (\S\ref{sec:intro-scope}) decides
by the votes of a fixed set of participants and stays safe while fewer than a
fixed fraction of them deviate arbitrarily from it.

The \emph{availability-finality dilemma} is the statement that no single
ledger protocol is both live under dynamic participation and safe under
network partition: a node cannot distinguish asleep nodes from nodes cut off
by a partition, so it either continues to extend its ledger, losing safety
under partition, or waits for a quorum, losing liveness when many nodes
sleep. The ebb-and-flow paper does not prove this statement; it attributes
the formalisation to Pass and Shi, ``The sleepy model of consensus''
(ASIACRYPT 2017), and its connection with the CAP theorem of Gilbert and
Lynch, ``Brewer's conjecture and the feasibility of consistent, available,
partition-tolerant web services'' (SIGACT News, 2002), that no replicated
service is both consistent and available under network partition, to
Lewis-Pye and Roughgarden, ``Resource pools and the CAP theorem'' (2020)
(Neu, Tas and Tse, Section~I-A). The volume uses it with that attribution, not
as a theorem.

The paper's resolution (its Section~I-B) gives every node two ledgers: a
\emph{finalised} ledger, safe under network partition, and a
\emph{dynamically available} ledger, live under dynamic participation, of
which the finalised ledger is a prefix. Each client chooses which of the two
ledgers to trust.

\subsection{Execution model and environments}\label{sec:model-execution}

The paper's security parameter is also its confirmation time, a number of
slots, and the paper writes it $\sigma$. It is written $\lambda$ here, as in
the \emph{Math Guide}, Definition ``Security parameter'' (\S``The security
parameter''); $\sigma$ is reserved for the confirmation depth of
Crosslink~2, and $\lambda$ is never a rate. The paper's counts of nodes and
of corrupted nodes are not given symbols: its model is stated here with
fractions.

\begin{definition}[Ebb-and-flow execution model]\label{def:ntt-exec}
The model has the following components (Neu, Tas and Tse, Section~III-A).
\begin{enumerate}[label=(\roman*)]
\item A fixed, finite set of nodes. Time proceeds in \emph{slots}, and nodes
  have synchronised clocks; bounded clock offsets are counted as network
  delay.
\item A public-key infrastructure that gives each node a unique
  cryptographic identity, and a random oracle that is the source of the
  protocol's randomness.
\item An \emph{adversary}, a probabilistic polynomial-time algorithm. Before
  the execution starts it corrupts a set of nodes, which are then
  \emph{adversarial}: they surrender their state to the adversary and
  deviate arbitrarily under its control. The other nodes are \emph{honest}
  and follow the protocol. Corruption is static: no node is corrupted later.
\item \emph{Sleep.} For every slot and every honest node the adversary
  chooses whether the node is \emph{awake} or \emph{asleep}, subject to the
  environment's constraints. An awake honest node executes the protocol in
  that slot; an asleep one does not, and the messages that would have
  reached it are queued and delivered in the first slot in which it is awake.
  Adversarial nodes are always awake.
\item \emph{Two transitions.} One asynchrony-to-synchrony transition, at the
  time $\GST$, after which honest-to-honest messages arrive within a delay
  bound of $\Delta$ slots; and one sleep-to-awake transition, at the time
  $\GAT$, after which every honest node is awake.
\item A finite time horizon polynomial in the security parameter $\lambda$.
\item The setting is permissioned: the identities of all nodes are common
  knowledge, while no honest node knows which nodes are awake.
\end{enumerate}
\end{definition}

The paper studies a single transition to synchrony, ``we follow the
long-standing practice in the BFT literature and study only a single such
transition'' (Section~III-A), and leaves proof-of-stake and proof-of-work
settings to later work (Section~III).

\begin{definition}[Environment for P1]\label{def:ntt-env1}
The pair $(\Adv_1(\beta), \Env_1)$ is the environment of the finality
property P1: the adversary $\Adv_1$ corrupts a fraction $\beta$ of the nodes;
before $\GST$ it delays messages arbitrarily, and after $\GST$ it delivers
every message between honest nodes within $\Delta$ slots; before $\GAT$ it
decides which honest nodes are awake and when, and after $\GAT$ every honest
node is awake. The adversary chooses $\GST$ and $\GAT$; honest nodes do not
know them, and they may be causal functions of the protocol's randomness.
\end{definition}

Without the transition $\GAT$ this adversary falls under the CAP theorem, so
that no protocol secure against it exists (Neu, Tas and Tse, footnote~5).

\begin{definition}[Environment for P2]\label{def:ntt-env2}
The pair $(\Adv_2(\beta), \Env_2)$ is the environment of the dynamic
availability property P2: at all times the adversary $\Adv_2$ delivers every
message between honest nodes within $\Delta$ slots, and it decides which
honest nodes are awake and when, subject to the constraint that at all times
at most a fraction $\beta$ of the awake nodes are adversarial and at least
one honest node is awake.
\end{definition}

\begin{remark}[Strict thresholds]\label{rem:ntt-strict}
The paper writes $\Adv_1^* := \Adv_1(1/3)$ and $\Adv_2^* := \Adv_2(1/2)$
(Section~III-C), and its appendix uses the first strictly: under
$(\Adv_1^*, \Env_1)$ the corrupted fraction satisfies $\beta < 1/3$. Read in
the same way, $(\Adv_2^*, \Env_2)$ requires that fewer than half of the awake
nodes be adversarial at every time. The tolerable fractions are therefore
strictly below one third and one half. A finite set of voters needs the
one-third bound in integer form, which the volume gives for Crosslink~2 with
the interface of its BFT protocol.
\end{remark}

\subsection{Ledger security and ebb-and-flow protocols}
\label{sec:model-security}

For ledgers $\LOG$ and $\LOG'$, the relation $\LOG \preceq \LOG'$ means that
$\LOG$ is a prefix of $\LOG'$, equality included. For a ledger protocol,
$\LOG_i^t$ is the ledger of node $i$ at time $t$.

\begin{definition}[Ledger security]\label{def:ntt-security}
Let $\Tconfirm$ be a polynomial function of $\lambda$. A ledger protocol
outputting ledgers $\LOG$ is \emph{secure after time $T$}, with transaction
confirmation time $\Tconfirm$, if $\LOG$ satisfies the following
(Neu, Tas and Tse, Definition~1).
\begin{enumerate}[label=(\roman*)]
\item \emph{Safety.} For all times $t \ge t' \ge T$ and all honest nodes $i$
  and $j$ awake at $t$ and $t'$ respectively, $\LOG_i^t \preceq \LOG_j^{t'}$ or
  $\LOG_j^{t'} \preceq \LOG_i^t$.
\item \emph{Liveness.} If a transaction is received by an awake honest node
  at a time $t \ge T$, then for every time $t' \ge t + \Tconfirm$ and every
  honest node $j$ awake at $t'$, the transaction is in $\LOG_j^{t'}$.
\end{enumerate}
The ledger is \emph{safe after $T$} or \emph{live after $T$} if the
respective property holds; it is \emph{safe} or \emph{live} if that property
holds after $T = 0$.
\end{definition}

\begin{definition}[Flexible protocol]\label{def:ntt-flexible}
A \emph{flexible protocol} is a pair of ledger protocols that take the same
input transactions and each output a ledger (Neu, Tas and Tse,
Definition~2).
\end{definition}

\begin{definition}[Ebb-and-flow protocol]\label{def:ntt-ebbflow}
A \emph{$(\beta_1, \beta_2)$-secure ebb-and-flow protocol} is a flexible
protocol whose two ledgers are an \emph{available ledger} $\LOGda$ and a
\emph{finalised ledger} $\LOGfin$, written $\LOGda[i]^t$ and $\LOGfin[i]^t$
at node $i$ and time $t$, such that, for the security parameter
$\lambda := \Tconfirm$, the following hold (Neu, Tas and Tse,
Definition~3).
\begin{enumerate}[label=(\roman*)]
\item \emph{P1, finality.} Under $(\Adv_1(\beta_1), \Env_1)$, the ledger
  $\LOGfin$ is safe at all times, and there is a constant multiple of
  $\max\{\GST, \GAT\} + \lambda$ after which $\LOGfin$ is live, except with
  probability $\negl(\lambda)$.
\item \emph{P2, dynamic availability.} Under $(\Adv_2(\beta_2), \Env_2)$,
  the ledger $\LOGda$ is secure, except with probability $\negl(\lambda)$.
\item \emph{Prefix.} For every honest node $i$ and every time $t$,
  $\LOGfin[i]^t \preceq \LOGda[i]^t$.
\end{enumerate}
A function is $\negl(\lambda)$ in the sense of the \emph{Math Guide},
Definition ``Negligible function'' (\S``Polynomial, exponential, and
negligible functions''): it is eventually below $\lambda^{-c}$ for every
constant $c > 0$, which is the paper's own definition.
\end{definition}

\begin{remark}[Classical regimes and optimal resilience]
\label{rem:ntt-extremes}
Setting $\GAT = 0$ in Definition~\ref{def:ntt-env1} gives the classical
partially synchronous network, in which every honest node is awake; setting
$\GST = 0$ and $\GAT = \infty$ gives a synchronous network with dynamic
participation. Both follow from the definitions. The paper adds that, for
its snap-and-chat construction (\S\ref{sec:model-sac}), whose instance with
permissioned Sleepy and Streamlet is $(1/3, 1/2)$-secure
(Theorem~\ref{thm:sac-secure}), $\LOGfin$, respectively $\LOGda$, then has
the optimal resilience achievable in that environment (Section~I-D), and
states without proof, as a standard bound, that no ebb-and-flow protocol
tolerates $\Adv_1(\beta_1)$ with $\beta_1 \ge 1/3$ in a partially synchronous
network, or $\Adv_2(\beta_2)$ with $\beta_2 \ge 1/2$ in a synchronous network
(Section~III-C). These
optimality statements are recorded as the paper makes them; no result of
this volume uses them.
\end{remark}

\subsection{The snap-and-chat construction}\label{sec:model-sac}

A subscripted $\Pi$ names a protocol. A \emph{longest-chain protocol} is a
ledger protocol whose nodes extend, and confirm blocks at a fixed depth in,
the longest chain they know.

\begin{construction}[Snap-and-chat]\label{def:sac}
The protocol $\Pisac$ composes a longest-chain protocol $\Pilc$ (in the
paper, the permissioned form of Sleepy, after Pass and Shi) with a BFT
protocol (in the paper, Streamlet, after Chan and Shi, ``Streamlet: Textbook
streamlined blockchains'', AFT 2020), its \emph{BFT sub-protocol}
(Neu, Tas and Tse, Section~III-B and Figure~5).
\begin{enumerate}[label=(\roman*)]
\item At node $i$ and time $t$ the protocol $\Pilc$ outputs a chain of
  confirmed blocks, whose sequence of transactions is the ledger
  $\LOGlc[i]^t$.
\item A proposal of the BFT sub-protocol carries a reference to a block of
  $\Pilc$; the chain that this block tips is the proposal's
  \emph{snapshot}. The BFT sub-protocol outputs at node $i$ and time $t$ the
  chain $\Ch_i^t$ of BFT blocks finalised in that node's view.
\item The sequence $\flatten(\Ch_i^t)$ is the sequence of transactions of
  the snapshot chains of the blocks of $\Ch_i^t$, concatenated in chain order
  (the first reading of Remark~\ref{rem:sac-deficiencies}, item~(3), which in
  the paper's setting coincides with the second). The function $\sanitize$ maps
  a sequence of transactions to its subsequence that keeps the first valid
  occurrence of each transaction and removes duplicates and transactions
  invalid at their position.
\item The ledgers are
  \begin{gather*}
    \LOGfin[i]^t := \sanitize\bigl(\flatten(\Ch_i^t)\bigr), \qquad
    \LOGda[i]^t := \sanitize\bigl(\LOGfin[i]^t \,\|\, \LOGlc[i]^t\bigr),
  \end{gather*}
  where $\|$ is concatenation.
\item \emph{Boycott rule}, the only change to the BFT sub-protocol: an
  honest node votes for a proposed BFT block only if it views the block's
  snapshot as confirmed in its own longest chain (Algorithm~1).
\end{enumerate}
\end{construction}

The paper's ledgers are formally sequences of blocks of $\Pilc$, from which
the transaction ledgers are obtained by sanitisation (its footnote~6). The
Prefix property of Definition~\ref{def:ntt-ebbflow} holds by construction,
since $\LOGda[i]^t$ begins with $\LOGfin[i]^t$.

\begin{theorem}[Snap-and-chat security]\label{thm:sac-secure}
The protocol $\Pisac$ with permissioned Sleepy as $\Pilc$ and Streamlet as
its BFT sub-protocol is a $(1/3, 1/2)$-secure ebb-and-flow protocol, in the
model of Definitions~\ref{def:ntt-exec}, \ref{def:ntt-env1}
and~\ref{def:ntt-env2} with the thresholds read strictly as in
Remark~\ref{rem:ntt-strict}.

Status: proved (ebb-and-flow paper).
\end{theorem}

The theorem is Neu, Tas and Tse, Section~III-C, Theorem~1, with its proof in
Appendix~B; the paper is peer-reviewed (IEEE Symposium on Security and
Privacy 2021). Its assumptions are those of the paper's model: a
permissioned node set, static corruption, a single $\GST$ and a single
$\GAT$, a public-key infrastructure and a random oracle; and a lottery
probability of $\Pilc$ below the bound of Theorem~\ref{thm:sleepy-secure},
which the proof uses for P1 and which implies the bound of Sleepy that it
uses for P2. The paper's statement leaves this hypothesis implicit. It is not
a theorem about Crosslink~2 or about proof of work.

\begin{theorem}[Security of permissioned Sleepy]\label{thm:sleepy-secure}
For the static version of permissioned Sleepy, with a lottery probability per
node and slot below the explicit bound of the paper's Theorem~2, a function
of the delay bound $\Delta$, the number of nodes and the corrupted fraction,
and for every $\GST$ and $\GAT$ of the environment $(\Adv_1^*, \Env_1)$, the
protocol $\Pilc$ is secure after a constant multiple of
$\max\{\GST, \GAT\} + \lambda$, with confirmation time $\lambda$, except with
probability $e^{-\Omega(\sqrt{\lambda})}$ (Definition~\ref{def:ntt-security};
\emph{Math Guide}, \S``Asymptotic notation'').

Status: proved (ebb-and-flow paper).
\end{theorem}

The theorem is Neu, Tas and Tse, Appendix~B, Theorem~2, with its proof in
Appendix~C, where the security parameter is written $\sigma$; the constant
depends on the lottery probability, the number of nodes, the corrupted fraction
and $\Delta$. It concerns Sleepy's permissioned lottery in the model of
Definition~\ref{def:ntt-exec}, not proof of work. The TFL book uses it only
in the sense that it ``essentially says that under certain conditions Prefix
Agreement holds except with probability $e^{-\Omega(\sqrt{\sigma})}$''
(chapter ``The Crosslink 2 Construction'', section ``Model for best-chain
protocols''), and the volume uses it in no other sense.

\begin{lemma}[Safety of the finalised ledger]\label{lem:sac-fin-safe}
In $\Pisac$ with permissioned Sleepy as $\Pilc$, its lottery probability
below the bound of Theorem~\ref{thm:sleepy-secure}, and Streamlet as its BFT
sub-protocol, under $(\Adv_2^*, \Env_2)$ the ledger $\LOGfin$ is a safe
prefix of the output of $\Pilc$ in the view of every honest node at all
times, except with probability $e^{-\Omega(\sqrt{\lambda})}$.

Status: proved (ebb-and-flow paper).
\end{lemma}

The lemma is Neu, Tas and Tse, Appendix~B, Lemma~5: under
$(\Adv_2^*, \Env_2)$ fewer than half of the awake nodes are adversarial, so
finalising a BFT block needs at least one honest vote, and by the boycott
rule an honest node votes only for a snapshot confirmed in its own view. The
proof derives from the security of $\Pilc$ under $(\Adv_2^*, \Env_2)$ that
any two snapshots confirmed in honest views, at any times, are
prefix-comparable, except with that probability; with the honest vote and
the boycott rule, the same then holds for the snapshots of finalised BFT
blocks even if those blocks conflict. Section~III-C states the step as
``since $\Pilc$ is safe, the fact that one honest node sees that snapshot as
confirmed implies that every honest node sees the same snapshot as
confirmed''. As for Theorem~\ref{thm:sac-secure}, the paper leaves the
lottery bound implicit.

\begin{proposition}[Case analysis of snap-and-chat]\label{prop:sac-cases}
In $\Pisac$ with permissioned Sleepy and Streamlet:
\begin{enumerate}[label=(\roman*)]
\item if no two honest views contain conflicting final BFT blocks, then
  $\LOGfin$ is safe, whether or not $\Pilc$ is safe;
\item under $(\Adv_2^*, \Env_2)$, with the lottery probability of $\Pilc$
  below the bound of Theorem~\ref{thm:sleepy-secure}, whether or not the BFT
  sub-protocol is safe, $\LOGfin$ is a prefix of the output of $\Pilc$ in every
  honest view, except with probability $e^{-\Omega(\sqrt{\lambda})}$.
\end{enumerate}

Status: proved (ebb-and-flow paper).
\end{proposition}

The first case follows from the safety of the BFT sub-protocol in the proof
of Theorem~1 (Neu, Tas and Tse, Appendix~B, Lemma~1 and the ledger
extraction); the second is Lemma~\ref{lem:sac-fin-safe}; Section~III-C and
Figure~9 illustrate both. In the first case the safe BFT sub-protocol orders
conflicting snapshots linearly and sanitisation removes the transactions
that the conflicts invalidate. In the second, fewer than half of the awake
nodes are adversarial, so $\Pilc$, with that lottery bound, is secure except
with that probability and
finalising a snapshot needs an honest vote, as in Lemma~\ref{lem:sac-fin-safe}.
The TFL book states the aim of Crosslink~2 beyond the second case: safety of the
finalised ledger even when the BFT protocol is entirely subverted, provided the
adversary has under half of the best-chain protocol's hash rate (chapter ``Notes
On Snap-and-Chat'', section ``Comment on security assumptions'').

\begin{remark}[Other BFT protocols and the scope of the theorems]
\label{rem:ntt-scope}
The paper replaces Streamlet by HotStuff and by PBFT (Section~III-D). The
HotStuff variant is proved only in an appendix of the arXiv version
(Appendix~D), which the proceedings version omits; status: proved
(ebb-and-flow paper), in the arXiv version only. The PBFT variant is argued,
``intuitively'', in prose; status: argued (ebb-and-flow paper). Within the
paper's model its theorems are therefore settled for snap-and-chat with
Sleepy and Streamlet, and with HotStuff in the arXiv appendix. They do not
extend to the paper's informal statement for proof of work, nor to PBFT.
\end{remark}

\subsection{The proof-of-work setting}\label{sec:model-pow}

In the paper's proof-of-work setting nodes are of two kinds:
\emph{miners}, measured by hash rate, extend the proof-of-work longest chain
$\Pilc$, and \emph{validators}, with unique cryptographic identities, run the
BFT sub-protocol (Neu, Tas and Tse, Section~V-B).

\begin{proposition}[Permissioned finality for proof of work]\label{prop:ntt-pow}
Let snap-and-chat use the proof-of-work longest chain as $\Pilc$ and
Streamlet, HotStuff or PBFT, adapted as in the paper's Sections~III-B
and~III-D, as its BFT sub-protocol, in a network in which communication is
asynchronous until $\GST$ and synchronous after it, validators are always awake,
and miners may sleep and wake at any time.
\begin{enumerate}[label=(\roman*)]
\item \emph{P1, finality.} The ledger $\LOGfin$ is safe at all times and
  live after $\GST$ if fewer than one third of the validators are
  adversarial, under half of the awake hash rate is adversarial, and the
  awake hash rate is bounded away from zero.
\item \emph{P2, dynamic availability.} If $\GST = 0$, the ledger $\LOGda$ is
  safe and live at all times if fewer than two thirds of the validators are
  adversarial, under half of the awake hash rate is adversarial, and the
  awake hash rate is bounded away from zero.
\end{enumerate}

Status: argued (ebb-and-flow paper).
\end{proposition}

The arXiv version of the paper states the proposition as ``Theorem
(Informal, Permissioned Finality for Permissionless PoW)'' in Section~V-B,
with the validator bounds written as 33 and 66 per cent; its only support is
the sentence ``Security is proved analogously to Section III-C'', and no
proof is given. The proceedings version contains only a prose summary.

\begin{remark}[The open question of the paper]\label{rem:ntt-open}
The hypotheses of Proposition~\ref{prop:ntt-pow} beyond the permissioned
case are, under P1, that under half of the awake hash rate is adversarial
and that the awake hash rate is bounded away from zero, without which
$\Pilc$ may not be live and $\LOGfin$ has nothing to finalise; and, under
P2, that fewer than two thirds of the validators are adversarial, without
which adversarial validators finalise blocks that are not on the longest
chain. The paper leaves open ``whether these additional requirements are
fundamental or a limitation of the snap-and-chat construction''
(Section~V-B).

Status: open problem.
\end{remark}

\subsection{From snap-and-chat to Crosslink 2}\label{sec:model-amend}

\begin{remark}[Deficiencies of snap-and-chat]\label{rem:sac-deficiencies}
Among the problems of snap-and-chat for Zcash that the TFL book records
(chapter ``Notes On Snap-and-Chat''), four determine the construction of
Crosslink~2.
\begin{enumerate}[label=(\arabic*)]
\item \emph{Scope of the one-honest-vote argument.} The argument of
  Lemma~\ref{lem:sac-fin-safe} needs an honest vote in every finalising
  quorum, hence fewer than two thirds of the votes adversarial, and the
  security of $\Pilc$, which in the paper's model needs fewer than half of
  the awake nodes adversarial; the book notes that it ``only applies when the
  number of malicious nodes $f$ is such that $n/3 < f < n/2$'', the range in
  which the BFT sub-protocol may be unsafe. It does not reach the target of
  Crosslink~2: safety of the finalised ledger with every validator
  subverted, provided the adversary has under half of the hash rate (section
  ``Comment on security assumptions'').
\item \emph{An assumption on the best chain.} The safety of $\Pilc$, from
  which the paper derives that a snapshot confirmed in one honest view is
  confirmed in all (Section~III-C), is to be supplemented by
  evidence that any party can check: a proposal carries the headers of a
  fixed number of best-chain blocks above its snapshot, whose proof of work
  confirms it (section ``Finalization Availability''; chapter ``The
  Crosslink 2 Construction'', section ``Structural additions'').
\item \emph{Sanitisation.} Sanitisation admits two readings: concatenate the
  transactions of all snapshots and keep each transaction that is valid at
  its position; or concatenate the blocks of all snapshots, remove duplicate
  blocks, and only then apply the same filter. The two coincide if double
  spends and missing inputs are the only reasons for a transaction to be
  invalid at a position, since a transaction omitted for either reason
  cannot later be reinstated. Zcash anchors violate the condition: the
  anchor rule admits the final treestate of any earlier block
  (\emph{Ironwood Guide}, \S``Anchors''), and the same treestate can be
  produced by unrelated transactions, so a transaction invalid for a missing
  anchor can become valid at a later position; expiry, by contrast, is never
  undone. The ledgers are therefore to be plain block chains, and
  sanitisation disappears (section ``Subtlety in the definition of
  sanitization''; chapter ``The Crosslink 2 Construction'', section
  ``Locally finalized chain'').
\item \emph{Spending finalised outputs.} Snap-and-chat cannot let a
  transaction spend a finalised output that comes from another fork's
  snapshot, because the best chain does not determine the finalised ledger,
  and validity of a block may depend only on the block and its ancestors.
  Each best-chain block is therefore to commit to a BFT block from which the
  finalised prefix is computable, and that commitment is not to move
  backwards along a chain (sections ``Spending finalized outputs'' and
  ``Finalization Availability'', the latter stating the rule
  ``Extension'').
\end{enumerate}

Status: argued (TFL book).
\end{remark}

\begin{remark}[The model of Crosslink 2]\label{rem:cl-model}
The TFL book takes neither synchrony nor partial synchrony as an implicit
assumption of the communication model of Crosslink~2: unless a sub-protocol
states otherwise, messages may be delayed arbitrarily or dropped, and each
sub-protocol imports its timing assumptions explicitly. This choice is
designed but unspecified. Its reason is argued (TFL book): the single-$\GST$
argument of partial synchrony is sound for the safety of a non-terminating
protocol, but cannot correctly be applied to its liveness, and properties
that hold after a first transition to synchrony need not hold after later
ones (chapter ``The Crosslink 2 Construction'', admonition ``Discussion of
incorrect applications of the GST formalization of partial synchrony to
continuously operating protocols.''). No $(\beta_1, \beta_2)$-security
theorem in the sense of Definition~\ref{def:ntt-ebbflow} exists for
Crosslink~2. Its safety statements are implications between properties of
single executions, without probability, and are not comparable with that
definition.
\end{remark}

\section{The model and the two sub-protocols}\label{sec:iface}

This section defines Crosslink~2 as a transformation of two black-box
protocols, fixes its communication and execution model and the notation for
chains, and states the interface that each black box must satisfy: the
best-chain protocol with Prefix Consistency and Prefix Agreement at a depth
$k$, and the BFT protocol with voting units, notarisation, the one-third
bound, its last final block and Final Agreement. The Zcash instances of the
score and of header-checkable work are specified facts; the prefix properties
of Zcash proof of work are an open problem, stated once, at the end of the
best-chain interface.

\subsection{The hybrid construction}\label{sec:iface-hybrid}

\begin{definition}[Crosslink 2]\label{def:crosslink}
\emph{Crosslink~2} is the hybrid construction that keeps a best-chain
protocol $\Piorigbc$ and a BFT protocol $\Piorigbft$ as black boxes and
transforms them into protocols $\Pibc$ and $\Pibft$ by structural additions,
changed validity rules and changed rules of honest behaviour, so that
best-chain blocks and BFT blocks commit to each other; these mutual
commitments are the \emph{cross-links}. Every node of Crosslink~2 takes part
in both $\Pibc$ and $\Pibft$, keeping a view of each; the roles of
bft-proposer, bft-validator and bc-block producer are optional, and every
party acting in one of them is a node of Crosslink~2.
\end{definition}

The definition follows the TFL book, chapter ``The Crosslink 2
Construction'', section ``Subprotocols''. Here $\Pi$ names a protocol; the
\emph{Crypto Guide} uses the letter for a scheme. The lineage of the
construction is the ebb-and-flow model and snap-and-chat
(\S\ref{sec:model}); its additions answer the deficiencies of
Remark~\ref{rem:sac-deficiencies}. In the classes of
Definition~\ref{def:classes}, Crosslink~2 is designed but unspecified.

\subsection{Executions, chains and prefix order}\label{sec:iface-chains}

\begin{definition}[Communication model]\label{def:comm}
Unless a sub-protocol states a stronger model, messages between participants
may be arbitrarily delayed or dropped, each message is received at most
once, and messages are non-malleably authenticated as coming from their
sender wherever a sub-protocol needs it. A sub-protocol that needs a
stronger model imports it explicitly.
\end{definition}

\begin{definition}[Execution and view]\label{def:exec}
An \emph{execution} of a protocol is the set of all events, sends and
receipts of messages and decisions of participants, of one run from its
initiation up to a given time; a prefix of an execution is an execution.
Times are totally ordered with least element $0$, and every execution starts
at time $0$. Up to each time an execution has finitely many events, and the
view and the best chain (Definition~\ref{def:bc-protocol}) of a node change
only at its events, so that its best chain changes at finitely many times in
every interval $[0, t]$. A node $i$ is \emph{honest at time $t$} in an
execution iff it has followed the protocol up to and including $t$ and the
adversary has not, up to $t$, cut it off from the rest of the network, for
example by an eclipse; a node so cut off counts as not honest even when it
follows the protocol (TFL book, admonition ``Discussion of [GKL2020]'s
communication model and network partition.''). A \emph{time series} $S$
assigns a value $S^t$ to each time $t$. A \emph{view} $V$ assigns to each node
$i$ and time $t$ a value $V_i^t$, here a chain.
\end{definition}

The total order of times is a modelling convenience; a happens-before order
on events would serve as well (TFL book, chapter ``The Crosslink 2
Construction'', section ``Conventions'').

\begin{definition}[Chain]\label{def:chain}
Let $x$ range over the two chain kinds $\bc$ and $\bft$. An
\emph{$x$-chain} is a non-empty sequence of $x$-blocks that starts at the
\emph{genesis block} $\Origin_x$, written $\Obc$ and $\Obft$, each later
block committing by a collision-resistant hash to its predecessor, its
\emph{parent}. The last block is the chain's \emph{tip}, and a block is
identified with the chain it tips. The \emph{height} $\mathsf{height}(B)$
of a block $B$ is its number of predecessors, so that the genesis block has
height $0$ (\emph{Ironwood Guide}, Definition ``Block chain, transaction,
consensus rules'', in \S``Scope and status'').
\end{definition}

The kind $\bc$ comprises the blocks of $\Piorigbc$ and of $\Pibc$, the kind
$\bft$ those of $\Piorigbft$ and of $\Pibft$; a property relative to one
protocol names it, as in ``valid in $\Piorigbft$''. The TFL book writes a
star for a metavariable over protocol names and as a wildcard within them;
the variable $x$ over the two kinds replaces it. The identification of a
block with its chain is sound under collision resistance of the parent hash:
exactly one chain has a given tip. The book uses heights without defining
them; Definition~\ref{def:chain} supplies the definition.

\begin{definition}[Truncation and depth]\label{def:trunc}
For an $x$-chain $C$ and $k \in \N$, the chain $C\trunc[x]{k}$, written
$C\trunc{k}$ where the kind is clear, is the prefix of $C$ whose tip has
height $\max\{0, \mathsf{height}(C) - k\}$: the chain $C$ without its last
$k$ blocks, or the genesis chain if $C$ has at most $k$ blocks. The block at
\emph{depth} $d$ of $C$ is the tip of $C\trunc{d}$; the tip of $C$ is at
depth $0$. For an object $P$ that names a parent $x$-block $B$ without being
an $x$-block itself, and for $k \ge 1$, the truncation is extended by
$P\trunc[x]{k} := B\trunc[x]{k - 1}$.
\end{definition}

The block at depth $d$ of $C$ is the last one that no rollback of $d$ blocks
can remove; the genesis chain cannot roll back. For $d \le \mathsf{height}(C)$
it has height $\mathsf{height}(C) - d$, the depth convention of the
\emph{Ironwood Guide}, \S``Anchors''.

\begin{remark}[Depth and confirmations]\label{rem:depth-conf}
The \emph{Consensus Guide}, Definition ``Confirmations'' (\S``Work, and the
best chain''), counts confirmations inclusively: a transaction in a block of
the best chain has as many confirmations as there are blocks from that block
to the tip, both included. A transaction in the block at depth
$d \le \mathsf{height}(C)$ of the best chain $C$ therefore has $d + 1$
confirmations, and for $k \le \mathsf{height}(C)$ the chain $C\trunc{k}$ holds
exactly the blocks of $C$ whose transactions have at least $k + 1$
confirmations.
For $k > \mathsf{height}(C)$ the chain $C\trunc{k}$ is the genesis chain,
whose block has fewer than $k + 1$ confirmations. The ebb-and-flow paper
uses ``depth'' for what this volume calls height.
\end{remark}

\begin{definition}[Prefix order and agreement]\label{def:prefix}
For $x$-chains $B$ and $C$, the relation $B \preceq_x C$ holds iff $B$ is an
initial segment of $C$, equality included, and $B \prec_x C$ iff
$B \preceq_x C$ and $B \ne C$; the subscript is omitted where the kind is
clear. The chains $B$ and $C$ \emph{agree}, written $B \agrees C$, iff
$B \preceq_x C$ or $C \preceq_x B$; they \emph{conflict} iff they do not
agree. A time series $S$ of $x$-chains, or a sequence of them indexed by an
ordered set, is \emph{$x$-linear} iff $t \le u$ implies
$S^t \preceq_x S^u$.
\end{definition}

Agreement is reflexive and symmetric but not transitive: the genesis block
agrees with every block, including two blocks that conflict. It is therefore
never written as an equivalence.

\begin{definition}[Last common ancestor]\label{def:lca}
For $x$-chains $A$ and $B$, the \emph{last common ancestor} $\lca(A, B)$ is
the longest $x$-chain $C$ with $C \preceq_x A$ and $C \preceq_x B$.
\end{definition}

The chain $\lca(A, B)$ exists, because $\Origin_x$ is a prefix of both $A$
and $B$ and the common prefixes are finitely many; and it is unique, because
the prefixes of $A$ are its initial segments, one of each length, so two
common prefixes of the same length coincide. The book uses the function
without defining it. If $A \preceq_x B$, then $\lca(A, B) = A$.

\begin{definition}[Agreement on a view]\label{def:agree-view}
An execution has \emph{Agreement} on a view $V$ iff for all times $t$ and
$u$ and all nodes $i$ and $j$, $i = j$ permitted, such that $i$ is honest at
$t$ and $j$ is honest at $u$, the chains $V_i^t$ and $V_j^u$ agree.
\end{definition}

\begin{lemma}[Linear prefix]\label{lem:linear-prefix}
If $A \preceq_x C$ and $B \preceq_x C$, then $A \agrees B$.

Status: proved (TFL book).
\end{lemma}

The TFL book states and proves the lemma in chapter ``The Crosslink 2
Construction'', section ``Conventions''; the proof below expands its
one-line argument.

\begin{proof}
The chains $A$ and $B$ are initial segments of $C$, and the initial segments
of a sequence are totally ordered by length: the shorter of $A$ and $B$ is an
initial segment of the longer.
\end{proof}

\subsection{The best-chain protocol}\label{sec:iface-bc}

The interface of this subsection is the TFL book's model for best-chain
protocols (chapter ``The Crosslink 2 Construction'', section ``Model for
best-chain protocols''). It is stated for $\Piorigbc$ and holds for $\Pibc$
with the validity rules that the construction adds.

\begin{definition}[Best-chain protocol]\label{def:bc-protocol}
A \emph{best-chain protocol} $\Piorigbc$ has the following interface.
\begin{enumerate}[label=(\roman*)]
\item \emph{Views and validity.} The view of a node is a set of
  $\bc$-chains rooted at $\Obc$. A validity rule on bc-blocks, the consensus
  rules, depends only on the content of the block and its ancestors, and a
  non-genesis block is valid only if its parent is valid; a \emph{valid
  chain} is a chain of valid blocks.
\item \emph{Extension.} Informally, it is hard to extend a valid chain unless
  one is selected by a random process that chooses block producers in
  proportion to a resource, with an approximately known and consistent time
  distribution, under an assumption on the share of the resource that honest
  nodes hold.
\item \emph{Score and best chain.} A function $\score$ maps valid chains into
  a set with a strict total order $<$, and $\score(H\trunc{1}) < \score(H)$ for
  every valid non-genesis chain $H$.
  The \emph{best chain} $\ch_i^t$ of node $i$ at time $t$ is a valid chain of
  highest score in its view, any rule breaking ties, and
  $\ch_i^0 := \Obc$; the best chain of an honest node is valid.
\item \emph{Transactions.} A valid block commits to a collection of
  transactions. The consensus rules on transactions are unchanged by
  Crosslink~2, including the rules that depend on the height of the block or
  on earlier blocks: transactions are neither reordered nor sanitised, and a
  block enters a chain of Crosslink~2 with its entire parent chain, so
  heights are preserved.
\item \emph{Headers.} Each bc-block has a \emph{bc-header} that commits to
  it. Header validity is necessary, but not sufficient, for block validity.
  For proof of work, the proof of work and the difficulty of a block are
  checkable from its header, given a valid parent chain.
\end{enumerate}
\end{definition}

Clause~(ii) is informal and enters no proof. The book asks of a
proof-of-work protocol that the score be its accumulated work.

% cited formula, not in compute/OUTPUT.md: work of a block, 2^256 over the
%   256-bit target threshold plus one (Consensus Guide, 'Work of a block')
\begin{remark}[The Zcash instance]\label{rem:bc-zcash}
For Zcash the score of clause~(iii) and the header check of clause~(v) are
specified facts. The \emph{work} $\work(B)$ of a block $B$ is
$\lfloor 2^{256} / (\mathsf{ToTarget}(\mathsf{nBits}) + 1) \rfloor$
(\emph{Consensus Guide}, Definition ``Work of a block'', in \S``Work, and the
best chain''; protocol specification, \S``Definition of Work''), where
$\mathsf{ToTarget}$ converts the header field $\mathsf{nBits}$ to its target
threshold, and the best valid block chain is the valid chain of greatest total
work, ties between leaf blocks going to the block received first
(\emph{Consensus Guide}, Definition ``Best valid block chain''; protocol
specification, \S``The Block Chain''). The score of a chain $C$ is therefore
\begin{gather*}
  \score(C) = \sum_{B \preceq_{\bc} C} \work(B).
\end{gather*}
For a valid block at height $h$, the field $\mathsf{nBits}$ equals
$\mathsf{ToCompact}(\mathsf{Threshold}(h))$ (protocol specification,
\S``Block Header Encoding and Consensus''), where $\mathsf{Threshold}(h) \le
\mathsf{PoWLimit} < 2^{256}$ (\S``Difficulty adjustment''), except that on
Testnet from height $299{,}188$ a minimum-difficulty block, whose header
timestamp $\mathsf{nTime}$ exceeds that of its predecessor by more than a
threshold,
has $\mathsf{nBits} = \mathsf{ToCompact}(\mathsf{PoWLimit})$ (ZIP 205,
``Change to difficulty adjustment on Testnet''; ZIP 208, ``Minimum difficulty
blocks on Testnet''); since $\mathsf{ToCompact}$ rounds down (\S``nBits
conversion''), $\mathsf{ToTarget}(\mathsf{nBits}) \le \mathsf{PoWLimit}$ in
both cases, so $\work(B) \ge 1$ for every valid block, and $\score(H) =
\score(H\trunc{1}) + \work(H) > \score(H\trunc{1})$ for every valid non-genesis
chain $H$; the tie rule is one that the interface allows. Header-checkable proof
of work is the condition of the \emph{FlyClient Guide}, Definition ``Block
header'' (\S``Header-only verification''): a header satisfies proof of work, and
its field $\mathsf{nBits}$ equals the value that the difficulty adjustment
assigns from the preceding headers, or, for a minimum-difficulty block on
Testnet, the encoding of $\mathsf{PoWLimit}$, so the check needs the headers of
the ancestors (\emph{Consensus Guide}, \S``Difficulty adjustment''; protocol
specification, \S``Difficulty adjustment''). ZIP 218 (Draft) changes the
averaging window of that adjustment from NU7 (ZIP 218, ``Averaging window'');
the volume states no window length.
\end{remark}

\begin{definition}[Coinbase-only block]\label{def:coinbase-only}
A \emph{coinbase transaction}, in the TFL book's abstraction, is a
bc-transaction that only distributes newly issued funds and has no inputs.
The predicate $\iscoinbaseonlyblock(B)$ on bc-blocks holds iff $B$ has
exactly one transaction and that transaction is a coinbase transaction; the
block is then \emph{coinbase-only}.
\end{definition}

The abstraction is narrower than the deployed coinbase transaction, which
also collects the fees of the block's other transactions and has consensus
rules of its own (\emph{Ironwood Guide}, Definition ``Coinbase
transaction'', in \S``The transaction format''; protocol specification,
\S``Coinbase Transactions'').

\begin{remark}[Switching depth]\label{rem:bc-switch}
The book asks informally of $\Piorigbc$ that, unless an adversary can censor
the node's knowledge of other chains, it be hard to make a node switch to a
best chain whose last common ancestor with its previous best chain lies more
than $\sigma$ blocks below that previous chain's tip; the volume writes a
generic depth $k$. The requirement enters no proof. Its precise form is
Prefix Consistency (Definition~\ref{def:pc}, below), whose instance $i = j$
bounds the depth of a node's own switches; Prefix Agreement
(Definition~\ref{def:pa}) gives no such bound.
\end{remark}

\begin{definition}[Prefix Consistency]\label{def:pc}
For $k \in \N$, an execution of a best-chain protocol has \emph{Prefix
Consistency} at depth $k$ iff for all times $t \le u$, every node $i$ honest
at $t$ and every node $j$ honest at $u$, $i = j$ permitted,
\begin{gather*}
  \ch_i^t\trunc[\bc]{k} \preceq_{\bc} \ch_j^u .
\end{gather*}
\end{definition}

The property is the consistency of Pass, Seeman and shelat, ``Analysis of the
Blockchain Protocol in Asynchronous Networks'' (EUROCRYPT 2017), Section~3.3,
which, unlike the common-prefix property of Garay, Kiayias and Leonardos,
``The Bitcoin backbone protocol: analysis and applications'' (EUROCRYPT
2015), also compares a node with its own later chains. The book
assumes it of executions of $\Piorigbc$; later sections state it of $\Pibc$.

\begin{definition}[Prefix Agreement]\label{def:pa}
For $k \in \N$, an execution of a best-chain protocol has \emph{Prefix
Agreement} at depth $k$ iff it has Agreement on the view
$(i, t) \mapsto \ch_i^t\trunc[\bc]{k}$ (Definition~\ref{def:agree-view}):
for all times $t$ and $u$ and all nodes $i$ honest at $t$ and $j$ honest at
$u$, $i = j$ permitted, $\ch_i^t\trunc[\bc]{k} \agrees \ch_j^u\trunc[\bc]{k}$.
\end{definition}

Prefix Consistency compares a truncated chain with an untruncated one,
Prefix Agreement two truncated chains. Both are stated at a generic depth
$k$, so that one definition serves every depth at which later sections
instantiate them. Prefix Consistency at depth $k$ implies Prefix Agreement at
depth $k$: for $t \le u$, both $\ch_i^t\trunc{k}$ (by Prefix Consistency)
and $\ch_j^u\trunc{k}$ (Definition~\ref{def:trunc}) are prefixes of
$\ch_j^u$, so they agree by Lemma~\ref{lem:linear-prefix}; the case $u < t$
is symmetric. The converse fails: a node whose best chain moves from
$\Obc a_1 a_2$ to $\Obc a_1'$ keeps Prefix Agreement at depth $1$ and breaks
Prefix Consistency.

\begin{remark}[Properties of executions]\label{rem:exec-props}
Prefix Consistency and Prefix Agreement, and the properties of the BFT
protocol and of the ledgers stated later in terms of executions, are
unconditional properties of single executions. The probability that an
execution has them is a separate question and is not part of the
definitions: a statement ``every execution with property $\mathcal{P}$ has
property $\mathcal{Q}$'' bounds the probability of failure of $\mathcal{Q}$
by that of $\mathcal{P}$ under any distribution of executions (TFL book,
section ``Model for best-chain protocols''). Assuming Prefix Consistency at
depth $k$ for proof of work implicitly excludes network partitions longer
than about $k$ block times: after the best chain on one side of a full
partition has grown by more than $k$ blocks, its truncation contains blocks
that the other side cannot know. A BFT protocol, by contrast, can remain safe
under asynchrony. A node that the adversary cuts off from the rest of the
network, for example by an eclipse, is counted as not honest
(Definition~\ref{def:exec}), even though it follows the protocol, so that the
properties make no statement about it.
This modelling matches the protocol specification, \S``The Block Chain'':
``the vast majority of well-connected nodes should eventually agree on their
best valid block chain up to that height''.
\end{remark}

% cited constants, not in compute/OUTPUT.md: target spacing 75 s before NU7
%   (protocol specification, 'Constants') and 25 s from NU7 (ZIP 218,
%   'Block target spacing')
\begin{remark}[Status of the prefix properties for Zcash]\label{rem:pow-prefix}
For Zcash proof of work, Prefix Consistency and Prefix Agreement at any
depth, at the 75-second block spacing before NU7 or the 25-second spacing of
ZIP 218 (Draft), ``Block target spacing'', are assumptions, not results. No
lower volume proves a prefix property: the analysis of the \emph{Consensus
Guide}, \S``Confirmation requirements'', rests on its Assumption ``Constant
difficulty'' (\S``The catch-up probability'') and its first-seen tie model,
in which a private branch displaces the public one only once it carries
strictly more work, and the \emph{FlyClient Guide}'s Assumption ``Persistence
and liveness'' (\S``The variable-difficulty backbone model'') is unproved in its
model. No source proves either property for Zcash, and no source computes an
adequate depth under the 25-second spacing. Every result of this volume whose
premise is a prefix property is therefore conditional for Zcash.

Status: open problem.
\end{remark}

\subsection{The BFT protocol}\label{sec:iface-bft}

The interface of this subsection is the TFL book's model for BFT protocols
(chapter ``The Crosslink 2 Construction'', section ``Model for BFT
protocols''). It is stated for $\Piorigbft$ and holds for $\Pibft$ with the
validity rules that the construction adds.

\begin{definition}[BFT protocol]\label{def:bft-protocol}
A \emph{BFT protocol} $\Piorigbft$ has the following interface.
\begin{enumerate}[label=(\roman*)]
\item \emph{Views and validity.} The view of a node is a set of
  $\bft$-chains rooted at $\Obft$. A validity rule on bft-blocks depends only
  on the content of the block and its ancestors, and a non-genesis block is
  valid only if its parent is valid. A \emph{bft-valid block} is a bft-block
  that satisfies the rule, and a \emph{bft-valid chain} a chain of such
  blocks.
\item \emph{BFT epochs and proposals.} An execution proceeds in a sequence of
  \emph{BFT epochs}. In each BFT epoch an honest proposer for that BFT epoch
  may make a \emph{bft-proposal}, which names its parent bft-block and its BFT
  epoch and is signed by its proposer; the signature is part of the proposal.
\item \emph{Proposal validity.} A validity rule on proposals depends only on
  the content of the proposal, its parent block and the validity of the
  proposer's signature; a \emph{valid proposal} satisfies it.
\end{enumerate}
\end{definition}

A proposal names a parent bft-block without being one, so
Definition~\ref{def:trunc} gives $P\trunc[\bft]{k} = B\trunc[\bft]{k-1}$ for
a proposal $P$ with parent $B$ and $k \ge 1$.

\begin{assumption}[Strong unforgeability of proposer signatures]
\label{ass:suf-cma}
The signature scheme of proposers, used for the signature on a proposal (the
\emph{inner} signature) and for the signature on the block formed from it
(the \emph{outer} signature), is strongly unforgeable under chosen-message
attack in the sense of the \emph{Crypto Guide}, Definition ``Strong
unforgeability'' (\S``Syntax and security goal'').

Status: assumption.
\end{assumption}

No lower volume proves strong unforgeability for a concrete scheme, and no
source fixes a signature scheme for proposers, so the property is a
hypothesis.

\begin{definition}[Voting units and notarisation]\label{def:units}
Each BFT epoch has a fixed number $n \in \N^{+}$ of \emph{voting units},
distributed among the validators of the BFT protocol, which use them to vote
for proposals of that BFT epoch. The holders of a BFT epoch's units form its
\emph{finaliser set}, and each holder is a \emph{finaliser}. A unit is
\emph{cast} for a proposal $P$ at a time iff $P$ is a valid proposal and a
ballot for $P$ authenticated by the unit's holder exists at that time. A
\emph{notarisation proof} $\proofsym_P$ for $P$ is obtainable from, and only
from, knowledge of ballots cast for $P$ that together satisfy a
\emph{notarisation rule}; the rule requires at least two thirds of the units
of $P$'s BFT epoch to be cast for $P$, and it may require more.
\emph{Soundness} of notarisation proofs is the clause ``only from'': a valid
proof for $P$ exists only if ballots meeting the rule have been cast.
\end{definition}

The term ``finaliser'' is Shielded Labs'; the volume uses it for the holders
of voting units throughout. In this volume $n$ denotes only the number of
voting units of a BFT epoch.

\begin{definition}[Non-honest casting]\label{def:nonhonest}
A voting unit is \emph{cast non-honestly} for a proposal of its BFT epoch
iff it is cast by a party other than its holder, including after compromise
of the holder's key; or it is \emph{double-cast}, that is, ballots cast it
for two distinct proposals of the BFT epoch; or its holder, following the
conditions for honest voting that the protocol prescribes, would not have
cast it in its view. For $\Pibft$ those conditions are the criteria of
Construction~\ref{def:honest-voting}. Several ballots casting one unit for
one proposal are not non-honest, provided the unit is not counted twice
towards the notarisation rule.
\end{definition}

\begin{assumption}[One-third bound]\label{ass:one-third}
In the execution considered, for every BFT epoch, strictly fewer than one
third of its voting units are ever cast non-honestly.

Status: assumption.
\end{assumption}

The assumption is the book's ``one-third bound on non-honest voting'', a
property of executions of $\Pibft$. The word ``ever'' ranges over the whole
execution: a key compromised long after its BFT epoch can still cast that
BFT epoch's units. The keys of honest validators must therefore remain secret
indefinitely, and a rotated key must be securely deleted (TFL book, chapter
``The Crosslink 2 Construction'', section ``Model for BFT protocols'',
admonition ``Security caveat'').

\begin{assumption}[Unforgeable ballots]\label{ass:ballots}
In every execution considered, no voting unit is cast by a party other than
its holder: ballot authentication is unforgeable, and no holder's key is
compromised.

Status: assumption.
\end{assumption}

The volume adds this hypothesis. The book states only that a ballot is
authenticated by the holder of the unit and fixes no scheme. Shielded Labs'
overview signs ballots with Ed25519 (\S\ref{sec:fs-keys}), an instance of
the \emph{Crypto Guide}'s Construction ``EdDSA, schematically''
(\S``RedDSA: re-randomisable Schnorr for Orchard''); Brendel, Cremers,
Jackson and Zhao, ``The Provable Security of Ed25519: Theory and Practice''
(IEEE Symposium on Security and Privacy 2021), analyse the unforgeability of
each variant of Ed25519 separately, and which variant the design uses is
unverified. No specification fixes a ballot scheme.
Assumption~\ref{ass:one-third} does not imply it: a forged cast
counts as non-honest under Definition~\ref{def:nonhonest}, and the bound
limits how many units are cast non-honestly without excluding any one such
cast. Of the results of this volume, only the withholding power of a
coalition of finalisers rests on it.

% quorum.py: q = ceil(2n/3): overlap 2q - n >= ceil(n/3) > ceil(n/3) - 1 for every n <= 1000: ok
% quorum.py: 2f+1 with f = floor((n-1)/3) is below 2n/3 exactly unless n = 3f + 1 (n <= 1000): ok
\begin{definition}[Integer threshold]\label{def:quorum}
For a BFT epoch with $n$ voting units, the notarisation rule requires at
least
\begin{gather*}
  q := \lceil 2n/3 \rceil
\end{gather*}
units cast for the proposal, and Assumption~\ref{ass:one-third} reads
$m \le \lceil n/3 \rceil - 1$, where $m$ is the number of units of the BFT
epoch ever cast non-honestly.
\end{definition}

The two forms are exact, because unit counts are integers: an integer is at
least $2n/3$ iff it is at least $\lceil 2n/3 \rceil$, and below $n/3$ iff it
is at most $\lceil n/3 \rceil - 1$. The book leaves the integer form
implicit. A count of $2f + 1$ units with $f = \lfloor (n - 1)/3 \rfloor$,
the usual form for protocols that tolerate $f$ faults, does not meet the
rule unless $n = 3f + 1$: writing $n = 3f + 1 + s$ with $s \in \{0, 1, 2\}$,
the inequality $2f + 1 \ge 2n/3 = 2f + (2 + 2s)/3$ holds iff $s = 0$.

\begin{definition}[bft-block]\label{def:bft-block}
A \emph{bft-block} is a pair $(P, \proofsym_P)$ of a proposal $P$ and a
notarisation proof for it, signed again by the proposer of $P$ (the outer
signature). It is valid in $\Piorigbft$ iff $P$ is a valid proposal,
$\proofsym_P$ is a valid proof that ballots cast for $P$ satisfy the
notarisation rule, and the outer signature on $(P, \proofsym_P)$ is valid.
The parent reference of a proposal hashes the whole parent block: proposal,
proof and outer signature. Neither $\proofsym_P$ nor the outer signature is
unique for a given $P$. A bft-block is \emph{in honest view} iff some party
observes it at a time at which that party is honest.
\end{definition}

Since each block commits to its proposal and each proposal hashes its whole
parent block with a collision-resistant hash, every bft-block determines its
chain of ancestors, although one proposal may underlie several valid blocks.
By Assumption~\ref{ass:suf-cma}, the outer signature is not malleable by a
third party.

\begin{lemma}[Per-BFT-epoch uniqueness]\label{lem:uniqueness}
In an execution that satisfies Assumption~\ref{ass:one-third} and in which
notarisation proofs are sound, all bft-valid blocks of one BFT epoch in
honest view carry the same proposal.

Status: proved (this volume), of a designed-but-unspecified construction.
\end{lemma}

\begin{proof}
Let $B$ and $B'$ be bft-valid blocks of a BFT epoch with $n$ voting units,
carrying proposals $P$ and $P'$, and suppose $P \ne P'$. By
Definition~\ref{def:bft-block} both blocks carry valid notarisation proofs,
so by soundness and Definition~\ref{def:quorum} there are sets $S$ and $S'$
of at least $q = \lceil 2n/3 \rceil$ units cast for $P$ and for $P'$
respectively. Then
\begin{gather*}
  |S \cap S'| \ge |S| + |S'| - n \ge 2q - n \ge \lceil n/3 \rceil .
\end{gather*}
The last inequality holds for every $n$: writing $n = 3a + r$ with
$r \in \{0, 1, 2\}$, one has $q = 2a + r$, so $2q - n = a + r$, while
$\lceil n/3 \rceil$ is $a$ for $r = 0$ and $a + 1$ otherwise. Every unit in
$S \cap S'$ is cast for the two distinct proposals $P$ and $P'$ of one BFT
epoch, so it is double-cast, hence cast non-honestly
(Definition~\ref{def:nonhonest}). The number $m$ of units of the BFT epoch
ever cast non-honestly is thus at least $\lceil n/3 \rceil$, contradicting
$m \le \lceil n/3 \rceil - 1$ (Definition~\ref{def:quorum}). Hence $P = P'$.
\end{proof}

% quorum.py: q = ceil(2n/3): overlap 2q - n >= ceil(n/3) > ceil(n/3) - 1 for every n <= 1000: ok
The proof does not use that the blocks are in honest view. The book's
version, the theorem ``On bft-valid blocks for a given epoch in honest view''
(chapter ``Security Analysis of Crosslink 2''), is adapted from Lemma~1 of
Chan and Shi's Streamlet paper; it lies below the chapter's marker that the
text after it is not updated for Crosslink~2, but it uses no ledger notation,
and it argues with the real bound $2n/3$. The proof above gives the integer
form, which the quorum script also checks for every $n \le 1000$.

\begin{remark}[Honestly submitted blocks]\label{rem:honestly-submitted}
An \emph{honest proposal} is a proposal made for a BFT epoch by its proposer
at a time at which that proposer is honest, and a bft-block is \emph{honestly
submitted} iff its proposer submits it at a time at which that proposer is
honest. Honesty at a time covers all earlier conduct
(Definition~\ref{def:exec}), and the protocol lets a proposer make at most one
proposal and at most one outer signature per BFT epoch, so, when each BFT
epoch has a single proposer, at most one honestly submitted bft-block exists
per BFT epoch; a dishonest proposer can produce several valid blocks for one
proposal, with different proofs or outer signatures. The book states that the
property matters for liveness and marks its use there as unchecked (``TODO check
that we are correctly using this in the liveness analysis'').

Status: open problem.
\end{remark}

\begin{definition}[Last final block]\label{def:blf}
The BFT protocol has an efficiently computable function
\begin{gather*}
  \bftlastfinal \colon \{\text{bft-blocks}\} \to
    \{\text{bft-blocks}\} \cup \{\bot\}
\end{gather*}
that satisfies the following requirements.
\begin{enumerate}[label=(\roman*)]
\item The value $\bftlastfinal(C)$ is $\bot$ iff $C$ is not bft-valid.
\item If $C$ is bft-valid, then $\bftlastfinal(C) \preceq_{\bft} C$, so that
  $\bftlastfinal(C)$ is bft-valid.
\item If $C$ is bft-valid, then for every bft-valid $D$ with
  $C \preceq_{\bft} D$, $\bftlastfinal(C) \preceq_{\bft} \bftlastfinal(D)$.
\item The genesis block satisfies $\bftlastfinal(\Obft) = \Obft$.
\end{enumerate}
For $\Piorigbft$ the function is written $\origbftlastfinal$.
\end{definition}

\begin{definition}[Finality in context]\label{def:final-context}
A bft-block $B$ is \emph{final in the context of} a bft-block $C$ iff
$B \preceq_{\bft} \bftlastfinal(C)$. No bft-block is final other than in a
context.
\end{definition}

Each bft-valid block $C$ determines its last final block $\bftlastfinal(C)$,
but no bft-block is objectively final. The notion concerns bft-blocks only;
the volume defines finality of a best-chain block separately, in the client
contract.

\begin{definition}[Final Agreement]\label{def:fa}
An execution of a BFT protocol has \emph{Final Agreement} iff for all
bft-valid blocks $C$ in honest view at a time $t$ and $C'$ in honest view at
a time $t'$, $\bftlastfinal(C) \agrees \bftlastfinal(C')$.
\end{definition}

Final Agreement can hold in an asynchronous communication model. If
Assumption~\ref{ass:one-third} is broken at any time of an execution, Final
Agreement may fail from that time on (TFL book, section ``Model for BFT
protocols'').

\begin{remark}[Adapting a concrete protocol]\label{rem:bft-adapt}
A concrete BFT protocol may need protocol-specific adaptation to fit the
interface before Crosslink~2 transforms it. Proposal validity is to be the
strongest property of a proposal that is objectively and feasibly checkable
from the proposal and its parent block when the proposal is made, and it
includes the check of the proposer's signature. Block validity is to be the
strongest property of a block that is objectively and feasibly checkable from
the block and its ancestors when the block is added to a chain, and it
includes the checks of the notarisation proof and of the outer signature.
For a node that has observed a valid block $C$, a rollback of its view past
$\origbftlastfinal(C)$ is to be infeasible, and its view up to that block is
to agree with that of every other honest node; Final Agreement formalises
the latter.
\end{remark}

\begin{example}[Adapted Streamlet]\label{ex:streamlet}
In Streamlet (Chan and Shi, ``Streamlet: Textbook streamlined blockchains'',
AFT 2020) a block is notarised once at least two thirds of the players have
voted for it, and in a notarised chain the prefix up to the second of three
adjacent blocks with consecutive BFT epoch numbers is final. The TFL book adapts
it to the interface as follows (section ``Model for BFT protocols'',
admonition ``Adapting the Streamlet BFT protocol.''). Once a proposal has
votes from two thirds of the units, its proposer aggregates them into a
notarisation proof, signs the resulting pair and submits the bft-block;
validity in $\Piorigbft$ then coincides with notarisation and is objectively
checkable. For the liveness analysis to carry over, the combined latency of
voting, aggregation and submission must be such that every node receives the
notarised block within two BFT epochs. For a valid $C$, the block
$\origbftlastfinal(C)$ is the last valid $B \preceq_{\bft} C$ that is
$\Obft$ or the second of three adjacent blocks with consecutive BFT epoch
numbers. Streamlet's liveness theorem guarantees progress after $\GST$
whenever five consecutive BFT epochs have honest leaders; the book refers to
this condition, and to Streamlet's liveness proof, in the admonition ``Why
does the Linearity rule allow keeping the same snapshot as the parent?'' of
its construction chapter.
\end{example}

\section{The Crosslink 2 construction}\label{sec:cons}

This section constructs $\Pibft$ and $\Pibc$ from the black boxes
$\Piorigbft$ and $\Piorigbc$ of Definition~\ref{def:crosslink}: the
parameters, the two cross-link fields and the functions derived from them,
the validity rules and the honest behaviour on each side, and Stalled Mode.
Rule names are those of the TFL book, chapter ``The Crosslink 2
Construction'', and each rule that answers an item of
Remark~\ref{rem:sac-deficiencies} is tied to that item; the Finality depth
rule serves bounded availability (\S\ref{sec:bounded}). In the classes of
Definition~\ref{def:classes}, every field, function and rule of this section
is designed but unspecified: the book defines it and no ZIP specifies it
(\S\ref{sec:intro-scope}). Components that no source fixes are marked as open
problems where they arise, and each result carries its status. Where no
ambiguity arises, a bc-header is identified with the bc-block it commits to
(Definition~\ref{def:bc-protocol}).

\subsection{Parameters}\label{sec:cons-params}

\begin{definition}[Parameters of Crosslink 2]\label{def:params}
Crosslink~2 has three parameters, fixed by the protocol:
\begin{enumerate}[label=(\roman*)]
\item the \emph{bc-confirmation depth} $\sigma \in \N^{+}$, the same for every
  node;
\item the \emph{finalisation gap bound} $L \in \N^{+}$;
\item a predicate $\isstalledblock$ on bc-blocks; a bc-block $H$ is
  \emph{stalled} iff $\isstalledblock(H)$ holds.
\end{enumerate}
\end{definition}

Every node uses the one value of $\sigma$, both in the rules below and in
computing its finalised chain. The book requires $L$ to be ``significantly
greater than $\sigma$'' (section ``Parameters'') and ``in practice'' at least
$2\sigma$ (section ``Stalled Mode''). This volume does not assume
$\sigma \le L$: Proposition~\ref{prop:short-L} shows that, for some
best-chain protocol of the interface and every $L < \sigma$, a best chain can
carry a non-stalled tip more than $L$ blocks above the finalised chain, so the
constraint is derived where it is needed.

The values of $\sigma$ and $L$ are an open problem. No value is chosen, no
source computes an adequate $\sigma$ at 75-second or at 25-second spacing
(Remark~\ref{rem:pow-prefix}), and no source gives a rule for choosing $L$;
the book's liveness argument leaves its ``dependence on $L$'' to a more
detailed argument that it does not give (chapter ``Security Analysis of
Crosslink 2''). No value of $\sigma$ or $L$ appears in this volume
(\S\ref{sec:intro-scope}). The predicate $\isstalledblock$ is designed but
unspecified: the book proposes coinbase-only stalled blocks for Zcash and
never fixes the predicate.

\subsection{Structural additions}\label{sec:cons-struct}

\begin{definition}[Cross-link fields]\label{def:fields}
\begin{enumerate}[label=(\roman*)]
\item Each bc-header carries, besides the fields of a $\Piorigbc$ header, a
  field $\contextbft$ that commits to a bft-block, the \emph{context} of the
  block; for a bc-block $H$ it is written $H.\contextbft$.
\item Each proposal and each non-genesis bft-block $B$ carries, besides the
  fields of $\Piorigbft$, a field $B.\headersbc$: a sequence of exactly
  $\sigma$ bc-headers, indexed from $0$, deepest first. The genesis
  bft-block has $\Obft.\headersbc = \mathsf{null}$.
\end{enumerate}
For a bc-chain $T$ of at least $\sigma + 1$ blocks, the \emph{$\sigma$-block
tail} of $T$ is the sequence of the headers of $T\trunc{\sigma - 1}, \dots,
T\trunc{1}, T$; its entry $j$ is the header of $T\trunc{\sigma - 1 - j}$. Every
entry is the header of a non-genesis block, so the deepest has a parent (the
book's admonition ``Why $\sigma + 1$?'').
\end{definition}

The field $\contextbft$ makes every best-chain block commit to a BFT block
from which its finalised prefix is computed, which answers item~(4) of
Remark~\ref{rem:sac-deficiencies}. The field $\headersbc$ is the evidence that
item~(2) asks for: the headers of $\sigma$ best-chain blocks above the
snapshot, which any party can check.

\begin{definition}[Snapshot]\label{def:snapshot}
For a proposal or bft-block $B$,
\[
  \snapshot(B) :=
  \begin{cases}
    \Obc, & \text{if } B.\headersbc = \mathsf{null},\\
    B.\headersbc[0]\trunc[\bc]{1}, & \text{otherwise},
  \end{cases}
\]
the parent of the block whose header is $B.\headersbc[0]$.
\end{definition}

If $B.\headersbc$ is the $\sigma$-block tail of a bc-chain $T$, then
$\snapshot(B) = T\trunc{\sigma}$: entry $0$ is the header of
$T\trunc{\sigma - 1}$, and truncations compose, since by
Definition~\ref{def:trunc} both $(T\trunc{\sigma-1})\trunc{1}$ and
$T\trunc{\sigma}$ are the prefix of $T$ whose tip has height
$\max\{0, \mathsf{height}(T) - \sigma\}$.
The snapshot then lies $\sigma$ blocks below the block of the newest carried
header, and $\mathsf{height}(\snapshot(B)) = \mathsf{height}(T) - \sigma$.

\begin{definition}[Last final context]\label{def:lf}
For a bc-block $H$,
\[
  \LF(H) := \bftlastfinal(H.\contextbft),
\]
the last block final in the context of $H.\contextbft$
(Definitions~\ref{def:blf} and~\ref{def:final-context}), and
$\snapshot(\LF(H))$ is the \emph{final snapshot} of $H$. For a bft-block $C$,
$\snapshot(\bftlastfinal(C))$ is the final snapshot of $C$.
\end{definition}

In the book's prose the final snapshot of a node's best tip is also called
its \emph{finalisation point}; Definition~\ref{def:candidate} lowers it to a
$\sigma$-confirmed prefix of the tip.

\begin{definition}[Candidate finalisation point]\label{def:candidate}
For a bc-block $H$,
\[
  \candidate(H) := \lca\bigl(\snapshot(\LF(H)),\, H\trunc{\sigma}\bigr).
\]
For the tip $H$ of a node's best chain, $\candidate(H)$ is the node's
\emph{candidate finalisation point}.
\end{definition}

When the final snapshot of $H$ is a prefix of $H$, as a rule of $\Pibc$ below
requires, the last common ancestor reduces to the lower of its two arguments;
the reduction is proved with that rule.

\begin{definition}[Finality depth]\label{def:fdepth}
For a bc-block $H$,
\[
  \finalitydepth(H) := \mathsf{height}(H) - \mathsf{height}(\snapshot(\LF(H))).
\]
\end{definition}

The finality depth is objective, a function of $H$ through its context and
not of any node's view, as a consensus rule must be. It counts the blocks of
$H$'s chain above the final snapshot when that snapshot is a prefix of $H$,
which the Last Final Snapshot rule below requires; without that rule a
difference of heights would measure no distance along $H$'s chain. This
form supersedes that of the book's earlier Crosslink~1, which measures from
$\lca(\snapshot(\LF(H)), H)$; when the final snapshot is a prefix of $H$ the
last common ancestor is the final snapshot itself, and the two forms
coincide.

\begin{remark}[Purpose of the additions]\label{rem:struct-purpose}
The book gives four reasons for the additions (section ``Structural
additions'' and its admonitions).
\begin{enumerate}[label=(\alph*)]
\item The field $\headersbc$ of a proposal or bft-block shows any party that
  sees it that $\sigma$ bc-blocks of valid proof of work existed above its
  snapshot when the proposal was made; snap-and-chat's snapshot reference
  carries no such evidence (item~(2) of Remark~\ref{rem:sac-deficiencies}).
\item Full headers, not hashes, let a validator check proof of work and
  difficulty from headers alone before it downloads any block, which limits
  denial of service: the proof of work of a Zcash block is checkable from its
  header (\emph{FlyClient Guide}, Definition ``Block header'',
  \S``Header-only verification''), and its target is a function of the
  headers of its ancestors and, on Testnet, of its own $\mathsf{nTime}$
  (\emph{Consensus Guide}, \S``Difficulty adjustment'';
  Remark~\ref{rem:bc-zcash}).
\item The genesis bft-block carries no headers, since there are not $\sigma$
  bc-headers to carry at genesis; $\snapshot(\Obft)$ is defined to be $\Obc$
  instead of making $\Obft$ refer to $\Obc$, which would close a hash cycle
  with $\Obc.\contextbft = \Obft$.
\item One field $\contextbft$ suffices, with no separate field for the last
  final block, because a context determines its last final ancestor through
  the efficiently computable function $\bftlastfinal$
  (Definition~\ref{def:blf}).
\end{enumerate}

Status: argued (TFL book).
\end{remark}

\begin{remark}[Encoding]\label{rem:encoding}
No source fixes the encoding of $\contextbft$ and $\headersbc$ or the place of
the cross-link data in a Zcash block. The specified
place for a commitment to new block data is the terminator of the block
commitment hash (\emph{FlyClient Guide}, Remark ``Extension of the block
commitment hash'', \S``The block commitment hash''), which no upgrade through
NU7 uses. Shielded Labs' \emph{Zcash Crosslink Design Overview} (draft 10,
14 September 2026), section ``3. The proof-of-work side: the fat pointer and
where it lives'', weighs five placements of its own cross-link data, from a
commitment through that terminator alone to a change of the header format,
judges the commit-only option ``probably insufficient, since some data must
actually be stored'', and lists the choice among the five as open (Part VI,
``Open questions'', ``Header and encoding'').

Status: open problem.
\end{remark}

\begin{figure}[ht]
\centering
\begin{tikzpicture}[x=12mm, y=10mm,
  blk/.style={draw, rectangle, minimum width=5.5mm, minimum height=5.5mm,
    inner sep=0pt, fill=white},
  fin/.style={blk, fill=tierin!25}, hdr/.style={blk, fill=tierfull!30},
  dts/.style={inner sep=1pt}, lnk/.style={-{Stealth[length=2mm]}, thick},
  snp/.style={-{Stealth[length=2mm]}, dashed},
  ctx/.style={-{Stealth[length=2mm]}, thick, tierrecv},
  lbl/.style={font=\footnotesize},
]
% bc-chain
\node[blk] (g) at (0,0) {};
\node[dts] (d1) at (1,0) {$\cdots$};
\node[fin] (s1) at (2,0) {};
\node[dts] (d2) at (3,0) {$\cdots$};
\node[blk] (s2) at (4,0) {};
\node[hdr] (h0) at (5,0) {};
\node[dts] (d3) at (6,0) {$\cdots$};
\node[hdr] (hs) at (7,0) {};
\node[dts] (d4) at (8,0) {$\cdots$};
\node[blk] (H) at (9,0) {};
\draw[lnk] (d1) -- (g);
\draw[lnk] (s1) -- (d1);
\draw[lnk] (d2) -- (s1);
\draw[lnk] (s2) -- (d2);
\draw[lnk] (h0) -- (s2);
\draw[lnk] (d3) -- (h0);
\draw[lnk] (hs) -- (d3);
\draw[lnk] (d4) -- (hs);
\draw[lnk] (H) -- (d4);
\node[lbl, below=1mm of g] {$\Obc$};
\node[lbl, below=1mm of s1, align=center]
  {$\snapshot(B_1)$\\$=\snapshot(\LF(H))$};
\node[lbl, below=1mm of s2] {$\snapshot(B_2)$};
\node[lbl, below=1mm of h0] {$[0]$};
\node[lbl, below=1mm of hs] {$[\sigma - 1]$};
\node[lbl, below=1mm of H] {$H$};
% the sigma carried headers of B2
\draw (4.7,0.45) -- (4.7,0.6) -- (7.3,0.6) -- (7.3,0.45);
\node[lbl, above] (hl) at (6,0.6) {$B_2.\headersbc$};
% bft-chain
\node[blk] (o) at (1,2.4) {};
\node[fin] (b1) at (3.5,2.4) {};
\node[blk] (b2) at (6,2.4) {};
\draw[lnk] (b1) -- (o);
\draw[lnk] (b2) -- (b1);
\node[lbl, above=1mm of o] {$\Obft$};
\node[lbl, above=1mm of b1] {$B_1 = \LF(H)$};
\node[lbl, above=1mm of b2] {$B_2$};
\draw[snp] (o) -- (g);
\draw[snp] (b1) -- (s1);
\draw[snp] (b2) -- (s2);
\draw[dotted] (b2) -- (hl);
\draw[ctx] (H) to[out=90, in=0]
  node[lbl, pos=0.4, right=1mm] {$H.\contextbft$} (b2);
\end{tikzpicture}
\caption{The cross-link structure. Solid black arrows point from child to
parent on the bc-chain (below) and the bft-chain $\Obft$, $B_1$, $B_2$
(above); dashed arrows point from a bft-block to its snapshot. The bracket
marks the $\sigma$ headers $B_2.\headersbc[0]$ to $B_2.\headersbc[\sigma - 1]$
(orange) that $B_2$ carries, and $\snapshot(B_2)$ is the parent of
$B_2.\headersbc[0]$. The green arrow is $H.\contextbft = B_2$; in the case drawn
the last block final in the context of $B_2$ is $B_1$ (blue), so $\LF(H) = B_1$
and the final snapshot of $H$ is $\snapshot(B_1)$ (blue). The Linearity rule
appears as $\snapshot(B_1)$ lying at or below $\snapshot(B_2)$ on one bc-chain.}
\label{fig:crosslinks}
\end{figure}

\subsection{Validity rules of the BFT protocol}\label{sec:cons-bft-rules}

\begin{definition}[Validity in $\Pibft$]\label{def:bft-validity}
\emph{Genesis bft-block rule}: the genesis bft-block $\Obft$ is
bft-block-valid, and $\Obft.\headersbc = \mathsf{null}$. A proposal $B$,
respectively a non-genesis bft-block $B$, is \emph{bft-proposal-valid},
respectively \emph{bft-block-valid}, iff the following hold.
\begin{enumerate}[label=(\roman*)]
\item \emph{Inherited origbft rules}: $B$ satisfies the proposal-validity,
  respectively block-validity, rules of $\Piorigbft$
  (Definitions~\ref{def:bft-protocol} and~\ref{def:bft-block}). These rules
  are assumed to include the \emph{Parent rule}, that the parent of $B$ is
  bft-block-valid; for a block they include a valid notarisation proof
  (Definition~\ref{def:units}) and a valid outer signature.
\item \emph{Linearity rule}:
  $\snapshot(B\trunc[\bft]{1}) \preceqbc \snapshot(B)$.
\item \emph{Tail Confirmation rule}: $B.\headersbc$ is the $\sigma$-block
  tail of a bc-valid chain.
\end{enumerate}
Proposal validity and block validity may differ, as they may in
$\Piorigbft$. The function $\bftlastfinal$ of $\Pibft$ is the function
$\origbftlastfinal$ of $\Piorigbft$, read with bft-block-validity and
$\Obft$: finality in context (Definition~\ref{def:final-context}) is
unchanged.
\end{definition}

For $\Pibft$, the bft-valid blocks of Definition~\ref{def:bft-protocol}, on
which Definitions~\ref{def:blf} and~\ref{def:fa} and Lemma~\ref{lem:uniqueness}
are stated, are the bft-block-valid blocks.

For a proposal $P$ with parent block $B'$, $P\trunc[\bft]{1} = B'$
(Definition~\ref{def:trunc}), so the Linearity rule compares a proposal with
its parent block, and a block with its parent. The rule is relative to the
parent's snapshot even when the parent is not final in the context of $B$,
because it must hold if the parent is later final in the context of a
descendant; equality is permitted, so a block may repeat its parent's
snapshot. The Tail Confirmation rule makes the blocks behind the carried
headers, not only the headers, part of bft-validity: bc-validity of the
chain needs the blocks themselves, and honest validators download and check
them before voting.

Along a valid bft-chain the Linearity rule makes the snapshots bc-linear:
for bft-blocks $C$ and $C'$,
\begin{equation}\label{eq:cons-snap-linear}
  C \preceqbft C' \text{ and } C' \text{ bft-block-valid}
  \implies \snapshot(C) \preceqbc \snapshot(C').
\end{equation}
The proof is by induction on the number of blocks of $C'$ not in $C$. If
there are none, $C = C'$. Otherwise $C'$ is not $\Obft$, its parent
$C'\trunc[\bft]{1}$ is bft-block-valid by the Parent rule and has $C$ as a
prefix, so the induction hypothesis gives $\snapshot(C) \preceqbc
\snapshot(C'\trunc[\bft]{1})$, and the Linearity rule for $C'$ gives
$\snapshot(C'\trunc[\bft]{1}) \preceqbc \snapshot(C')$.

\begin{remark}[Simultaneous definition of validity]
\label{rem:validity-recursion}
Definition~\ref{def:bft-validity} refers to bc-validity through the Tail
Confirmation rule, and Definition~\ref{def:bc-validity} refers to
bft-validity through the Valid context rule; the two are one simultaneous
definition of bft-validity and bc-validity. Each rule of either definition
refers only to blocks that the block under test commits to by hash, directly
or through a chain of commitments: its ancestors, the blocks behind its
carried headers, its context, and their ancestors. Validity is defined by
recursion along this commitment relation, restricted to its well-founded
part: a block from which some chain of commitments is infinite, in particular
a block that commits through a cycle to its own hash, is neither
bft-block-valid nor bc-block-valid. The genesis case serves
the same end: $\Obc$ commits to $\Obft$ through $\Obc.\contextbft$, so
$\Obft$ carries no bc-header.
\end{remark}

\begin{remark}[Objectivity of validity]\label{rem:objective}
Validity in $\Pibft$ depends only on a proposal or block and on the blocks
it commits to by hash, never on an observer's best chain, and it is checked
independently of the ballots for a proposal. A proposal that fails a rule
of Definition~\ref{def:bft-validity} remains invalid however many ballots
name it: no voting unit counts as cast for it (Definition~\ref{def:units}), and
no honest node accepts a block formed from it as bft-block-valid, whatever
notarisation proof the block carries.
\end{remark}

\begin{remark}[Purpose of the Linearity rule]\label{rem:linearity-purpose}
Besides the route to Assured Finality under Final Agreement
(Theorem~\ref{thm:fa-af}; Remark~\ref{rem:linearity-book}), the book gives
the following reasons for the rule (admonitions ``What does the Linearity
rule do?'' and ``Why does the Linearity rule allow keeping the same snapshot
as the parent?'').
\begin{enumerate}[label=(\alph*)]
\item It replaces the Increasing Score rule of Crosslink~1, which required
  each snapshot of a valid bft-chain to equal its parent's or to have higher
  score, and it implies that rule, since score strictly increases along a
  bc-chain (Definition~\ref{def:bc-protocol}).
\item It excludes proposals of a bc-valid chain that forks from a much
  earlier block, where the difficulty could have been much lower.
\item It limits the disruption that an adversary that has subverted $\Pibft$
  can cause to the best chain: the finalised chain is a $\Pibc$ chain, so its
  safety is no worse than that of $\Pibc$ alone for a rollback of any depth.
\item Either the snapshot advances on that of the parent bft-block, or it
  needs no further validation, having been validated already.
\item It removes sanitisation: by~\eqref{eq:cons-snap-linear} the snapshots
  of a valid bft-chain extend one another, so the finalised history is one
  bc-chain whose transactions keep their consensus rules, which answers
  item~(3) of Remark~\ref{rem:sac-deficiencies}.
\item It permits equality so that an honest proposer can always make a valid
  proposal, by repeating its parent's snapshot; liveness of $\Piorigbft$ may
  require honest proposals at a minimum rate. The book prefers this to null
  proposals, which would need rules of their own in $\Piorigbft$.
\end{enumerate}

Status: argued (TFL book).
\end{remark}

\subsection{Honest proposal and honest voting}\label{sec:cons-bft-honest}

\begin{construction}[Honest proposal]\label{def:honest-proposal}
An honest proposer $i$ that makes a proposal $P$ at time $t$ proceeds as
follows.
\begin{enumerate}[label=(\roman*)]
\item It makes no proposal before its best chain $\ch_i^t$ has at least
  $\sigma + 1$ blocks.
\item It sets $P.\headersbc$ to the $\sigma$-block tail of $\ch_i^t$ if $P$
  then satisfies the Linearity rule, and otherwise to the field $\headersbc$
  of the parent bft-block of $P$.
\end{enumerate}
\end{construction}

When $\ch_i^t$ has exactly $\sigma + 1$ blocks the first branch gives the
snapshot $\Obc$, which satisfies the Linearity rule against $\Obft$ because
$\snapshot(\Obft) = \Obc$. When the parent is $\Obft$ the first branch always
applies, since $\Obc$ is a prefix of every bc-chain, so the second branch
never copies the null field of $\Obft$. The first branch satisfies the Tail
Confirmation rule because an honest node's best chain is bc-valid
(Definition~\ref{def:bc-protocol}); the second satisfies the Linearity rule
with equality and the Tail Confirmation rule because the parent does. The
first branch can fail the Linearity rule after a reorganisation of the
proposer's best chain, since the block $\sigma$ deep on the new chain need
not descend from the parent's snapshot.

\begin{construction}[Honest voting]\label{def:honest-voting}
An honest validator $v$ that considers a proposal $P$ at time $t$ proceeds as
follows.
\begin{enumerate}[label=(\roman*)]
\item It updates its views of both sub-protocols with $P.\headersbc$: it
  downloads the bc-blocks of those headers that it lacks and checks their
  bc-validity, and for each such block it validates the bft-chain that the
  block's $\contextbft$ references, if it has not seen it, and recursively
  the bc-chains carried by that bft-chain's blocks. Non-recursive checks come
  first: the notarisation proof of a bft-block before its ancestry, and the
  proof of work of carried headers, with their connection to an already
  validated bc-chain, before any block is downloaded.
\item It then votes for $P$ only if both criteria hold:
  \begin{itemize}
  \item \emph{Valid proposal criterion}: $P$ is bft-proposal-valid;
  \item \emph{Confirmed best-chain criterion}:
    $\snapshot(P) \preceqbc \ch_v^t\trunc{\sigma}$.
  \end{itemize}
\end{enumerate}
\end{construction}

The Confirmed best-chain criterion places $\snapshot(P)$ at depth at least
$\sigma$ in the validator's best chain after the update; the block of the
newest carried header, $P.\headersbc[\sigma - 1]$, need not lie on that
chain, which may fork above the snapshot. Validity of the carried blocks is
checked in step~(i) and not inferred from the criterion, because the
criterion concerns the validator's best chain while validity is objective
(Remark~\ref{rem:objective}). Snap-and-chat's voting condition asked only
that the validator view the snapshot as confirmed, and its proposals carried
no headers (item~(2) of Remark~\ref{rem:sac-deficiencies}). Steps~(i)
and~(ii) are the honest-voting conditions of $\Pibft$ to which
Definition~\ref{def:nonhonest} refers.

\begin{remark}[Refusal to vote]\label{rem:only-if}
Honest voting is an ``only if'' condition, not ``if and only if'': a
validator may refuse to vote for a proposal that meets both criteria, on any
good ground, including information obtained outside the protocol, at a
possible cost to the liveness of $\Pibft$. The book declines to make refusal
on long rollbacks a voting rule only because such a rule would be stateful
(chapter ``Questions About Crosslink''). The book's argument in section
``Complementarity'' of the chapter ``The Arguments for Bounded Availability
and Finality Overrides'', that validators asked to stop validating cause a
finalisation stall which Stalled Mode then contains, depends on this
allowance.
\end{remark}

\begin{remark}[Validation cost]\label{rem:validation-cost}
Honest voting validates recursively: a referenced bft-chain in full, the
bc-chain carried by each of its blocks, the bft-chains those bc-blocks
reference, and so on. The book bounds the cost only by the validation order
and by the Linearity rule (admonition ``Wait what, how much validation is
that?''). Non-recursive checks come first; the proof of work of carried
headers, and their connection to a validated bc-chain, is checked before
blocks are downloaded, so bc-block validation proceeds no faster than the
network finds proof of work; and by~\eqref{eq:cons-snap-linear} the
snapshots of one valid bft-chain lie on one bc-chain, so one bc-chain needs
validation per bft-chain, which is, if $\Pibft$ is safe and the adversary
has no overwhelming advantage in proof of work, effectively one bc-chain
overall with short side branches. No bound on the validation cost is proved.

Status: argued (TFL book).
\end{remark}

\subsection{Validity rules of the best-chain protocol}\label{sec:cons-bc-rules}

\begin{definition}[Validity in $\Pibc$]\label{def:bc-validity}
This definition completes the simultaneous definition of
Remark~\ref{rem:validity-recursion}. \emph{Genesis bc-block rule}:
$\Obc.\contextbft = \Obft$; hence $\LF(\Obc) = \Obft$
(Definition~\ref{def:blf}), and the final snapshot of $\Obc$ is
$\snapshot(\Obft) = \Obc$. A non-genesis bc-block $H$ is
\emph{bc-block-valid} iff the following hold.
\begin{enumerate}[label=(\roman*)]
\item \emph{Inherited origbc rules}: $H$ satisfies the block-validity rules
  of $\Piorigbc$ (Definition~\ref{def:bc-protocol}), which include validity
  of its parent.
\item \emph{Valid context rule}: $H.\contextbft$ is bft-block-valid.
\item \emph{Extension rule}: $\LF(H\trunc[\bc]{1}) \preceqbft \LF(H)$.
\item \emph{Last Final Snapshot rule}: $\snapshot(\LF(H)) \preceqbc H$.
\item \emph{Finality depth rule}: $\finalitydepth(H) \le L$ or
  $\isstalledblock(H)$.
\end{enumerate}
A bc-block is \emph{bc-valid} iff it is $\Obc$ or bc-block-valid, and a
\emph{bc-valid chain} is a chain of bc-valid blocks; by the Inherited origbc
rules a bc-valid block tips a bc-valid chain. The book classes every change
as one to non-contextual validity; contextual validity, the validity of a
block's transactions against the state of its chain, is that of $\Piorigbc$.
\end{definition}

The Valid context rule makes $\LF(H)$ a bft-block-valid block, not $\bot$
(Definition~\ref{def:blf}). The Extension rule makes the commitment of
item~(4) of Remark~\ref{rem:sac-deficiencies} monotone along a chain: the
last final context of a block extends that of its parent. The Last Final
Snapshot rule places the final snapshot on the block's own chain, so that its
transactions can spend outputs of the finalised prefix, and makes the finality
depth a distance along that chain. The Finality depth rule forces the
predicate $\isstalledblock$, not coinbase-only blocks; coinbase-only blocks
are the policy that the book proposes for Zcash.

% ledger_checks.py: 1. candidate is the shorter of the final snapshot and the truncated tip (3255 cases, sigma <= 7): ok
\begin{lemma}[Candidate under the Last Final Snapshot rule]
\label{lem:candidate-lfs}
Let $H$ be a bc-block with $\snapshot(\LF(H)) \preceqbc H$, as the Last Final
Snapshot rule requires. Then $\candidate(H)$ is the shorter of
$\snapshot(\LF(H))$ and $H\trunc{\sigma}$.

Status: proved (this volume), of a designed-but-unspecified construction.
\end{lemma}

\begin{proof}
Both arguments of the last common ancestor are prefixes of $H$: the final
snapshot by hypothesis, and $H\trunc{\sigma}$ by Definition~\ref{def:trunc}.
By Lemma~\ref{lem:linear-prefix} they agree; let $A$ be the shorter and $B$
the other, so $A \preceqbc B$. Then $A$ is a common prefix of both, and every
common prefix of both is a prefix of $A$ and so no longer than $A$; hence
$\lca(A, B) = A$ (Definition~\ref{def:lca}).
\end{proof}

A script checks the lemma on instances.

\begin{lemma}[Prefix Lemma]\label{lem:lf-monotone}
If $H_1$ and $H_2$ are bc-valid blocks with $H_1 \preceqbc H_2$, then
$\LF(H_1) \preceqbft \LF(H_2)$.

Status: proved (TFL book), of a designed-but-unspecified construction.
\end{lemma}

\begin{proof}
The book proves the lemma by induction on the distance between $H_1$ and
$H_2$ with the Extension rule (chapter ``Security Analysis of Crosslink 2'').
If $H_1 = H_2$ there is nothing to prove. Otherwise $H_2$ is not genesis and
$H_1 \preceqbc H_2\trunc[\bc]{1}$, which is bc-valid by the Inherited origbc
rules and lies one block closer to $H_1$; the induction hypothesis gives
$\LF(H_1) \preceqbft \LF(H_2\trunc[\bc]{1})$, the Extension rule for $H_2$
gives $\LF(H_2\trunc[\bc]{1}) \preceqbft \LF(H_2)$, and the prefix order is
transitive.
\end{proof}

With~\eqref{eq:cons-snap-linear}, the lemma makes final snapshots monotone
along a valid chain: if $H_1 \preceqbc H_2$ are bc-valid, then
$\snapshot(\LF(H_1)) \preceqbc \snapshot(\LF(H_2))$, since $\LF(H_2)$ is
bft-block-valid by the Valid context rule and Definition~\ref{def:blf}.

\subsection{Honest block production and fork choice}\label{sec:cons-bc-honest}

\begin{construction}[Honest block production]\label{def:honest-production}
An honest producer of a bc-block $H$ follows the rules of
Definition~\ref{def:bc-validity}; in particular it produces a stalled block
when the Finality depth rule requires one. To choose $H.\contextbft$ it
considers the bft-block-valid blocks $T$ of its view whose last final block
$\bftlastfinal(T)$ extends the last final context of the parent,
\[
  \bigl\{\, T : T \text{ bft-block-valid and }
    \LF(H\trunc[\bc]{1}) \preceqbft \bftlastfinal(T) \,\bigr\},
\]
the \emph{candidate contexts} of $H$. Among the candidate contexts $T$ for
which $H$ with $H.\contextbft = T$ satisfies the rules of
Definition~\ref{def:bc-validity}, it chooses one of the longest, $C$,
breaking ties by the greatest $\score(\snapshot(\bftlastfinal(C)))$ and then
by the smallest hash, and sets $H.\contextbft := C$.
\end{construction}

The restriction to choices that satisfy the rules is this volume's reading;
the book's rule chooses among all candidate contexts, a set defined without
the Last Final Snapshot and Finality depth rules, of which a longest member
can make $H$ invalid. The book's reasons for the rule are argued (TFL
book): choosing the longest chain follows Streamlet, which builds on the
longest notarised chain rather than the longest finalised one; and
the tie-break by the score of the final snapshot inhibits an adversary from
finalising a chain of its own with lower score, while an adversary chain of
greater score has already won the race of proof of work.

\begin{proposition}[Existence of a valid context]\label{prop:context-exists}
Let $G$ be a bc-valid block, and let $H$ be a bc-block with parent $G$ that
satisfies the Inherited origbc rules.
\begin{enumerate}[label=(\roman*)]
\item The context $G.\contextbft$ is a candidate context of $H$. With
  $H.\contextbft = G.\contextbft$, block $H$ satisfies the Valid context,
  Extension and Last Final Snapshot rules, $\finalitydepth(H) =
  \finalitydepth(G) + 1$, and $H$ is bc-valid if it is stalled whenever
  this depth exceeds $L$.
\item Every candidate context satisfies the Valid context and Extension
  rules as context of $H$. A candidate context $T$ whose final snapshot
  conflicts with $G$ violates the Last Final Snapshot rule, and the rules of
  $\Pibft$ do not exclude that a longest candidate context is such a $T$; so
  the literal longest choice need not give a valid block.
\end{enumerate}

Status: proved (this volume), of a designed-but-unspecified construction.
\end{proposition}

\begin{proof}
(i) If $G = \Obc$, then $G.\contextbft = \Obft$, which is valid by the
Genesis bft-block rule; otherwise $G.\contextbft$ is valid by the Valid
context rule for $G$. In both cases $\bftlastfinal(G.\contextbft) = \LF(G)$,
so the context is a candidate. With it, $\LF(H) = \LF(G)$, so the Extension
rule holds with equality; the final snapshot of $H$ is that of $G$, which is
a prefix of $G$ (the Last Final Snapshot rule for $G$, or
$\snapshot(\Obft) = \Obc$ for $G = \Obc$) and hence of $H$; and
$\finalitydepth(H) = \mathsf{height}(G) + 1 - \mathsf{height}(\snapshot(\LF(G)))
= \finalitydepth(G) + 1$. The Finality depth rule then holds if $H$ is
stalled whenever that depth exceeds $L$.

(ii) A candidate context $T$ is bft-block-valid, and as context of $H$ it
gives $\LF(H) = \bftlastfinal(T) \succeq_{\bft} \LF(G)$. Let $S$ be its final
snapshot and suppose $S$ conflicts with $G$. If $S \preceqbc H$, then either
$S = H$, so $G \preceqbc S$, or $S \preceqbc G$; both contradict the conflict.
Hence $S \not\preceqbc H$ and the Last Final Snapshot rule fails.

For a witness, take an instance of $\Piorigbft$ in which every
bft-block-valid block is final in its own context, $\bftlastfinal(C) = C$,
which meets Definition~\ref{def:blf}, and a block $G$ whose context
$B_1 := G.\contextbft = \LF(G)$ has $\snapshot(B_1) \precbc G$. Let a
bc-valid branch fork from $G$'s chain at a block $D$ with
$\snapshot(B_1) \preceqbc D \precbc G$ and carry $\sigma + 1$ blocks
$Y_1, \dots, Y_{\sigma+1}$ above $D$, each citing $B_1$ and stalled where the
Finality depth rule requires: the Valid context rule holds since $B_1$ is
valid, the Extension rule holds for $Y_1$ since $\LF(D) \preceqbft \LF(G) = B_1$
(Lemma~\ref{lem:lf-monotone}) and for the others with equality, and the Last
Final Snapshot rule holds since $\snapshot(B_1) \preceqbc D$. Let $B_2$, a child
of $B_1$ valid in $\Piorigbft$, carry the $\sigma$-block tail of $Y_{\sigma+1}$.
Then $\snapshot(B_2) = Y_1$, which conflicts with $G$; the block $B_2$ satisfies
the Linearity rule, since $\snapshot(B_1) \preceqbc D \precbc Y_1$, and the
Tail Confirmation rule, so it is bft-block-valid, and $\bftlastfinal(B_2) = B_2
\succeq_{\bft} B_1 = \LF(G)$ makes it a candidate context of $H$ with final
snapshot $Y_1$. In a view whose bft-blocks are those of the chain of $B_2$ it is
the only longest candidate context, so the literal choice selects it.
\end{proof}

The construction is accordingly read as a choice among the candidate contexts
that satisfy the consensus rules, which by~(i) is non-empty for a block $H$
that is stalled whenever $\finalitydepth(G) + 1 > L$; such a block exists
under the producibility constraint of Remark~\ref{rem:producible}. The volume
does not claim that the literal choice is valid whenever $\Pibft$ is not
subverted: Assumption~\ref{ass:one-third} limits non-honest casting, not the
honesty of the nodes whose best chains confirmed a snapshot, and no claim is
made about the executions in which a longest candidate context has a final
snapshot off the producer's chain.

\begin{definition}[Unchanged fork choice]\label{def:fork-choice}
Crosslink~2 has \emph{unchanged fork choice}: no honest node, and in
particular no honest bc-block producer, uses information from $\Pibft$ other
than the rules of Definition~\ref{def:bc-validity} to decide which bc-valid
chain to follow. An honest node's best chain $\ch_i^t$ is a bc-valid chain of
highest score in its view, as in $\Piorigbc$ (Definition~\ref{def:bc-protocol}),
with validity that of Definition~\ref{def:bc-validity}.
\end{definition}

The definition states, under a neutral name, the book's admonition titled
``Attention'' (section ``$\Pibc$ honest block production''). The book's
reason is that the rules depending on $\Pibft$ were constructed to preserve
the safety of $\Pibc$ relative to $\Piorigbc$, and that any further
constraint could let an adversary that subverts $\Pibft$ steer the best
chain in ways the safety argument does not consider. A chain whose context
lags far behind the newest BFT blocks can therefore still become the best
chain (chapter ``Questions About Crosslink''). The book's liveness argument
rests on this definition.

\begin{remark}[Bidirectional dependence]\label{rem:bidirectional}
Crosslink~2 makes each sub-protocol depend on the other: the Valid context,
Extension, Last Final Snapshot and Finality depth rules make bc-validity
depend on $\Pibft$, and the Tail Confirmation rule and honest voting make
$\Pibft$ depend on $\Pibc$. The ebb-and-flow paper of Neu, Tas and Tse, in
its Appendix~E, ``Bouncing Attack on Casper FFG'', warns that such a
bidirectional interaction is ``intricate to reason about and a gateway for
liveness attacks''; in Casper FFG it arises because finalisation constrains
the fork choice. The book's mitigation is Definition~\ref{def:fork-choice}:
the fork choice of $\Pibc$ is not modified (chapter ``Security Analysis of
Crosslink 2''). Asked whether the fork choice should be changed to follow the
last final snapshot, the book answers ``Probably not. I don't know how to
repair the safety and liveness arguments'' (chapter ``Questions About
Crosslink'').

Status: argued (TFL book).
\end{remark}

\subsection{Stalled Mode}\label{sec:cons-stalled}

\begin{definition}[Stalled Mode]\label{def:stalled}
A bc-block $H$ is \emph{in Stalled Mode} iff $\finalitydepth(H) > L$;
stalled blocks are those of Definition~\ref{def:params}. By the Finality
depth rule, every bc-valid block in Stalled Mode is stalled. Entry into
Stalled Mode is per block and objective: whether $H$ is in Stalled Mode is a
function of $H$, and no network-wide flag exists. Exit is automatic: along a
chain, Stalled Mode ends at the first block that cites a context whose final
snapshot lies within $L$ blocks below it.
\end{definition}

The book states exit in the same terms: Stalled Mode ends ``as soon as the
finalization point catches up to within $L$ blocks (if it does without an
intentional rollback)'' (chapter ``Notes On Snap-and-Chat'').

\begin{remark}[Proposed Zcash policy]\label{rem:stalled-policy}
For Zcash the book proposes that stalled blocks be coinbase-only,
$\isstalledblock = \iscoinbaseonlyblock$ (Definition~\ref{def:coinbase-only}),
and never fixes the predicate (chapter ``The Arguments for Bounded
Availability and Finality Overrides'', section ``How to block hazards''). A
coinbase-only block spends no funds: its one transaction collects the block
subsidy and, there being no other transaction, no fees (\emph{Ironwood
Guide}, Definition ``Coinbase transaction'', \S``The transaction format'').
Mining pools therefore cannot distribute rewards in stalled blocks. Whether
shielded coinbase outputs are also disabled in Stalled Mode is an open
problem: the book says only that disabling them ``might'' be wanted.
\end{remark}

\begin{remark}[Producibility of stalled blocks]\label{rem:producible}
The only formal constraint on $\isstalledblock$, imposed by the book's
liveness argument (chapter ``Security Analysis of Crosslink 2'', section
``Liveness argument''), is that the set of bc-blocks $H$ with
$\iscoinbaseonlyblock(H)$ and $\isstalledblock(H)$ be producible at all
times: a bc-block producer can always produce a block in it.
Proposition~\ref{prop:context-exists}(i) supplies the context such a block
needs. Under the proposed policy the set is the set of coinbase-only blocks.
\end{remark}

\begin{remark}[Independence of Stalled Mode]\label{rem:stalled-independent}
Stalled Mode is technically independent of the rest of Crosslink~2: omitting
the Finality depth rule removes it and leaves every other rule in place. The
book argues that the protocol then keeps its security advantages over
snap-and-chat, stays simpler and still lets transactions spend finalised
outputs (item~(4) of Remark~\ref{rem:sac-deficiencies}); the incentive for
producers to advance the finalisation point is weaker, and honest producers
still advance it (chapter ``Security Analysis of Crosslink 2'', section
``Liveness argument'').

Status: argued (TFL book).
\end{remark}

% ledger_checks.py: 2. finality depth grows by at most one per block and candidate is monotone (2000 trials, 116000 blocks); the first block deeper than L has depth L + 1 (31622 cases): ok
\begin{lemma}[Growth of finality depth]\label{lem:fd-step}
Let $H$ be a non-genesis bc-valid block. Then
\[
  \finalitydepth(H) \le \finalitydepth(H\trunc[\bc]{1}) + 1
  \quad\text{and}\quad \candidate(H\trunc[\bc]{1}) \preceqbc \candidate(H).
\]

Status: proved (this volume), of a designed-but-unspecified construction.
\end{lemma}

\begin{proof}
Write $G = H\trunc[\bc]{1}$, which is bc-valid, and let $S_G$ and $S_H$ be the
final snapshots of $G$ and $H$. By Lemma~\ref{lem:lf-monotone},
$\LF(G) \preceqbft \LF(H)$, and $\LF(H)$ is bft-block-valid by the Valid
context rule and Definition~\ref{def:blf}; so~\eqref{eq:cons-snap-linear}
gives $S_G \preceqbc S_H$ and hence $\mathsf{height}(S_G) \le
\mathsf{height}(S_H)$. Since $\mathsf{height}(H) = \mathsf{height}(G) + 1$,
\begin{align*}
  \finalitydepth(H) &= \mathsf{height}(G) + 1 - \mathsf{height}(S_H)\\
  &\le \mathsf{height}(G) + 1 - \mathsf{height}(S_G) = \finalitydepth(G) + 1.
\end{align*}
For the second claim, $\candidate(G) \preceqbc S_G \preceqbc S_H$, and
$\candidate(G) \preceqbc G\trunc{\sigma} = H\trunc{\sigma + 1} \preceqbc
H\trunc{\sigma}$. So $\candidate(G)$ is a common prefix of $S_H$ and
$H\trunc{\sigma}$, and so is $\candidate(H)$, the longest common prefix.
Both are prefixes of $H\trunc{\sigma}$, so they agree
(Lemma~\ref{lem:linear-prefix}), and the longer is $\candidate(H)$; hence
$\candidate(G) \preceqbc \candidate(H)$.
\end{proof}

Entry into Stalled Mode is therefore gradual. Along a bc-valid chain the
finality depth starts at $\finalitydepth(\Obc) = 0$ and grows by at most one
per block, so the first block in Stalled Mode, the first block that the
Finality depth rule forces to be stalled, has finality depth exactly $L + 1$.
A script checks the lemma and this consequence on random chains.

\begin{proposition}[Branch dependence of exit]\label{prop:exit-branch}
Let $H$ be a bc-valid block.
\begin{enumerate}[label=(\roman*)]
\item If $H$ is not in Stalled Mode, its final snapshot lies on its own
  chain at most $L$ blocks below it: $\snapshot(\LF(H)) \preceqbc H$ and
  $\mathsf{height}(H) - \mathsf{height}(\snapshot(\LF(H))) \le L$. Stalled Mode
  can therefore end on a branch only at a block whose final snapshot lies on
  that branch within $L$ of the block.
\item Let $X$ and $Y$ be conflicting bc-chains with fork point
  $D = \lca(X, Y)$, and suppose that $\Pibft$ has Final Agreement
  (Definition~\ref{def:fa}) and that some bft-block-valid block $C^*$ in
  honest view has final snapshot $S^*$ with $D \precbc S^* \preceqbc X$.
  If $H \preceqbc Y$, $\mathsf{height}(H) > \mathsf{height}(D) + L$ and an
  honest node has validated $H$, then $H$ is stalled.
\end{enumerate}

Status: proved (this volume), of a designed-but-unspecified construction.
\end{proposition}

\begin{proof}
(i) The Last Final Snapshot rule gives the prefix relation, and a block not
in Stalled Mode has finality depth at most $L$.

(ii) To validate $H$ the honest node obtained $H.\contextbft$ and checked it
(the Valid context rule), so $H.\contextbft$ is a bft-block-valid block in
honest view (Definition~\ref{def:bft-block}). Final Agreement makes
$A := \bftlastfinal(C^*)$ and $B := \LF(H)$ agree. First, $S^*$ is not a
prefix of $Y$: otherwise it would be a common prefix of $X$ and $Y$, hence a
prefix of $D$, against $D \precbc S^*$. If $A \preceqbft B$,
then~\eqref{eq:cons-snap-linear}, with $B$ valid, gives
$S^* \preceqbc \snapshot(B)$; the Last Final Snapshot rule gives
$\snapshot(B) \preceqbc H \preceqbc Y$, so $S^* \preceqbc Y$, which is
excluded; this case does not occur. Hence $B \preceqbft A$,
and~\eqref{eq:cons-snap-linear}, with $A$ valid as a prefix of the valid
$C^*$, gives $\snapshot(B) \preceqbc S^* \preceqbc X$. With
$\snapshot(B) \preceqbc H \preceqbc Y$, the final snapshot of $H$ is a common
prefix of $X$ and $Y$, so $\mathsf{height}(\snapshot(B)) \le \mathsf{height}(D)$
and $\finalitydepth(H) \ge \mathsf{height}(H) - \mathsf{height}(D) > L$. The
Finality depth rule makes $H$ stalled.
\end{proof}

Under the hypotheses of~(ii), Stalled Mode persists on $Y$ above
$\mathsf{height}(D) + L$, whatever the score of $Y$. Unchanged fork choice
(Definition~\ref{def:fork-choice}) keeps an honest node on $Y$ while $Y$ has
the highest score in its view, and every block that node produces on $Y$
above that height is stalled. The proposition makes no claim about which
assumptions exclude its hypothesis, a final snapshot beyond the fork point on
the branch that honest nodes do not follow.

\begin{proposition}[Forced entry into Stalled Mode]\label{prop:forced-entry}
\begin{enumerate}[label=(\roman*)]
\item Let $G \preceqbc H$ be bc-valid blocks with $G \neq \Obc$, such that
  every block $H'$ with $G \precbc H' \preceqbc H$ has the final snapshot of
  $G$, as when each of those blocks repeats its parent's $\contextbft$. Then
  $\finalitydepth(G) \ge 1$ and
  \[
    \finalitydepth(H) = \finalitydepth(G) + \mathsf{height}(H) -
    \mathsf{height}(G),
  \]
  so $H$ is stalled whenever $\mathsf{height}(H) - \mathsf{height}(G) \ge L$.
\item A bc-valid block $H'$ with $\finalitydepth(H') \le L$, as when it cites
  a context whose final snapshot lies on its chain at most $L$ blocks below
  it, is not in Stalled Mode, whatever the finality depth of its parent, and
  no block of a bc-valid chain within $L - \finalitydepth(H')$ blocks above
  $H'$ is in Stalled Mode.
\item If more than one third of the voting units of a BFT epoch are cast for
  no proposal of that BFT epoch, then no proposal of the BFT epoch is
  notarised, and no bft-block of the BFT epoch is bft-block-valid.
\end{enumerate}

Status: proved (this volume), of a designed-but-unspecified construction.
\end{proposition}

\begin{proof}
(i) Let $S$ be the final snapshot of $G$; it is a prefix of $G$ by the Last
Final Snapshot rule. It is not $G$ itself: $\snapshot(\Obft) = \Obc \neq G$,
so $S = G$ would make $\LF(G)$ carry the header of a child of $G$, and $G$
would commit to its own hash through $G.\contextbft$, its ancestor $\LF(G)$
and that child, a cycle excluded as in Remark~\ref{rem:validity-recursion}.
Hence $\mathsf{height}(S) < \mathsf{height}(G)$ and $\finalitydepth(G) \ge 1$.
The block $H$ has final snapshot $S$, so $\finalitydepth(H) = \mathsf{height}(H)
- \mathsf{height}(S) = \finalitydepth(G) + \mathsf{height}(H) -
\mathsf{height}(G)$, which exceeds $L$ when $\mathsf{height}(H) -
\mathsf{height}(G) \ge L$; the Finality depth rule makes
$H$ stalled.

(ii) The first claim is Definition~\ref{def:stalled}. For a block $H''$ that
lies $k$ blocks above $H'$ on a bc-valid chain, $k$ applications of
Lemma~\ref{lem:fd-step} give $\finalitydepth(H'') \le \finalitydepth(H') +
k$, which is at most $L$ when $k \le L - \finalitydepth(H')$.

(iii) Let $n$ be the number of voting units of the BFT epoch and $q =
\lceil 2n/3 \rceil$ the threshold of Definition~\ref{def:quorum}. For any
proposal $P$ of the BFT epoch, the units cast for $P$ are among those cast
for some proposal of the BFT epoch, fewer than $n - n/3 = 2n/3 \le q$; so
fewer than $q$ units are cast for $P$. A notarisation proof for $P$ is
obtainable only from ballots that meet the notarisation rule
(Definition~\ref{def:units}), which requires $q$ units cast for $P$, so none
exists, and a bft-block $(P, \proofsym_P)$ needs a valid notarisation proof
(Definition~\ref{def:bft-block}).
\end{proof}

By~(i) and~(ii), forcing Stalled Mode against producers that cite newer
contexts requires producing every block of the best chain above the last
block $H'$ that cites a newer context, $L - \finalitydepth(H') + 1$
consecutive blocks, at most $L$ since $\finalitydepth(H') \ge 1$ by~(i); a
party with more than half of the hash rate can do so, since it can reject
every other producer's blocks (\emph{Consensus Guide}, Remark ``Majority is a
different problem'', in \S``Reward-normalised exposure''), and a smaller
party only by chance, with a probability that no result of this volume
bounds; by~(iii), holders of more than one third of a BFT epoch's units who
withhold their ballots, and whose units no one else casts, prevent every block
of that BFT epoch, so finalisation cannot advance through a block of it. The
book's sentence that Stalled Mode ``can only be entered by consensus of a
larger validator set, or if there is an availability failure of the
finalization protocol'' contrasts the design with an earlier proposal to
trigger Stalled Mode by a key (chapter ``The Arguments for Bounded
Availability and Finality Overrides'', section ``Complementarity''); by~(i),
block producers can force entry while $\Pibft$ is live, so the sentence is
not a property of the construction.

\begin{remark}[Escape from Stalled Mode]\label{rem:escape}
A bc-block producer is constrained to produce stalled blocks exactly while no
candidate context in its view has a final snapshot on the new block's chain
at most $L$ blocks below the new block; the book states this only ``roughly
speaking'', as holding while its view of the finalisation point is not
advancing. The book cautions that an adversary that has subverted $\Pibft$
without stopping the finalisation point from advancing can always avoid the
constraint (section ``Stalled Mode''): a block is in Stalled Mode only when its
final snapshot lies more than $L$ blocks below it, and such an adversary can
keep citing contexts whose final snapshots lie within $L$ of its blocks.

Status: argued (TFL book).
\end{remark}

\section{The finalised and bounded-available chains}\label{sec:ledgers}

This section defines the state that a Crosslink~2 node keeps besides its
views of $\Pibc$ and $\Pibft$: the locally finalised chain $\fin$ with its
no-rollback update and hazard record, the locally bounded-available chain
$\ba_\mu$, the syncing gate that withholds both from clients, and the
computation that a node performs on each change of its best chain, with the
lag of $\fin$ behind the best tip. The lemmas proved here are those that the
safety proofs use. Every definition and construction of the section is
designed but unspecified in the classes of Definition~\ref{def:classes}; it
comes from the sections ``Parameters'', ``Locally finalized chain'',
``Locally bounded-available chain'' and ``Syncing and checkpoints'' of the
TFL book's chapter ``The Crosslink 2 Construction'', and
Construction~\ref{def:node-update} orders the definitions of
\S\ref{sec:cons} and of this section as a node applies them. Open problems are
marked where they arise. Throughout, $i$ is a node and $\ch_i^t$ its best chain
at time $t$ (Definition~\ref{def:bc-protocol}).

\subsection{The locally finalised chain}\label{sec:led-fin}

\begin{definition}[Locally finalised chain]\label{def:fin}
Each node $i$ keeps a bc-chain $\fin_i^t$, its \emph{locally finalised
chain} at time $t$, as follows.
\begin{enumerate}[label=(\roman*)]
\item The chain starts at $\fin_i^0 = \Obc$.
\item \emph{Update rule} $\updatefin$: at each time $t$ at which the best
  chain of node $i$ changes, with $s$ the time of the previous change, or
  $0$ if there is none,
  \[
    \fin_i^t :=
    \begin{cases}
      \candidate(\ch_i^t), & \text{if } \fin_i^s \preceqbc
        \candidate(\ch_i^t),\\
      \fin_i^s, & \text{otherwise}.
    \end{cases}
  \]
  If moreover $\candidate(\ch_i^t)$ conflicts with $\fin_i^s$, node $i$
  records a \emph{finalisation safety hazard} that holds $\ch_i^t$ and the
  history of its updates of $\fin_i$ including and since the last one whose
  value was a prefix of $\candidate(\ch_i^t)$.
\item Between changes of its best chain, $\fin_i$ is constant.
\end{enumerate}
\end{definition}

An update never replaces $\fin_i$ by a chain that does not extend it. A
candidate that is a strict prefix of $\fin_i^s$, as after a short
reorganisation, leaves $\fin_i$ unchanged and is no hazard; only a
conflicting candidate is. The book states that a hazard can occur only if
global safety assumptions are violated.

\begin{lemma}[Local finalization linearity]\label{lem:fin-linear}
If node $i$ follows the update rule of Definition~\ref{def:fin} up to time
$u$, then $\fin_i^t \preceqbc \fin_i^u$ for every $t \le u$: the series
$\fin_i$ is bc-linear up to $u$ (Definition~\ref{def:prefix}).

Status: proved (this volume), of a designed-but-unspecified construction.
\end{lemma}

\begin{proof}
The book names the property and asserts that the update guarantees it by
construction (section ``Locally finalized chain''); the proof is the
following. Between changes of $\ch_i$ the value of $\fin_i$ is constant, and at
each change it is replaced by an extension of itself: by $\candidate(\ch_i^t)$
under the case condition $\fin_i^s \preceqbc \candidate(\ch_i^t)$, or by itself.
Induction over the changes in $(t, u]$, with transitivity of $\preceqbc$, gives
$\fin_i^t \preceqbc \fin_i^u$.
\end{proof}

\begin{lemma}[Local fin-depth]\label{lem:fin-depth}
In every execution of Crosslink~2, for every node $i$ honest at time $t$
there is a time $r \le t$ with $\fin_i^t \preceqbc \ch_i^r\trunc{\sigma}$.

Status: proved (TFL book), of a designed-but-unspecified construction.
\end{lemma}

\begin{proof}
The proof is the book's (section ``Locally finalized chain''). For every
bc-valid block $H$, $\candidate(H) \preceqbc H\trunc{\sigma}$, since a last
common ancestor is a prefix of each argument (Definition~\ref{def:lca}).
Node $i$, honest at $t$, has followed the update rule up to $t$
(Definition~\ref{def:exec}), and its best chains are bc-valid
(Definition~\ref{def:bc-protocol}). If $\fin_i$ has not changed up to $t$,
then $\fin_i^t = \Obc$; take $r = 0$, where $\ch_i^0 = \Obc$ and
$\Obc\trunc{\sigma} = \Obc$ (Definitions~\ref{def:bc-protocol}
and~\ref{def:trunc}). Otherwise let $r \le t$ be the time of its last change
up to $t$. A change occurs only in the first case of the update rule, so
$\fin_i^t = \fin_i^r = \candidate(\ch_i^r) \preceqbc \ch_i^r\trunc{\sigma}$.
\end{proof}

\begin{remark}[Purpose of the no-rollback rule]\label{rem:fin-local}
When the best chain of a node reorganises to another fork, even by fewer
than $\sigma$ blocks, $\candidate(\ch_i^t)$ can move back. If $\Pibft$ has
Final Agreement (Definition~\ref{def:fa}), the new candidate should then be
an ancestor of the old one and, as the best chain grows, move forward again
along the same path. Clients of the node must not see that temporary
regression, so $\fin_i$ is node state that never rolls back (the book's
admonition ``Why does $\fin_i$ need to be maintained using local state?'').

Status: argued (TFL book).
\end{remark}

\begin{remark}[Candidate or final snapshot]\label{rem:candidate-choice}
The update uses $\candidate(H)$ rather than the final snapshot
$\snapshot(\LF(H))$ so that the finalised prefix is $\sigma$-confirmed: the
Last Final Snapshot rule places the final snapshot on $H$'s chain without
guaranteeing it a depth of $\sigma$ or more, while the book's safety argument
from Prefix Agreement needs $\candidate(H) \preceqbc H\trunc{\sigma}$, the
property behind Lemma~\ref{lem:fin-depth}. The book records the alternative of
strengthening the Last Final Snapshot rule to $\snapshot(\LF(H)) \preceqbc
H\trunc{\sigma}$ and leaves the choice open: ``Choose between these options
based on what works well for the security proofs and finalization latency''
(admonition ``Why do we need to use $\candidate(H)$ rather than
$\snapshot(\LF(H))$?''). It expects the two to differ rarely in practice.

Status: open problem.
\end{remark}

\subsection{The locally bounded-available chain}\label{sec:led-ba}

\begin{definition}[Bounded-available depth]\label{def:mu}
Each node $i$ chooses a \emph{bounded-available depth} $\mu \in \N$ with
$0 < \mu \le \sigma$, possibly different from the choices of other nodes; the
default is $\mu = \sigma$.
\end{definition}

The book restricts $\mu$ to at most $\sigma$ because a larger value has no
use. A node that chooses $\mu < \sigma$ does so at its own risk: the smaller
depth may increase the risk of rollback against its bounded-available chain,
and it does not affect $\fin_i$ (section ``Parameters'').

\begin{definition}[Locally bounded-available chain]\label{def:ba}
The \emph{locally bounded-available chain} of node $i$ at time $t$ is
\[
  (\ba_\mu)_i^t :=
  \begin{cases}
    \ch_i^t\trunc{\mu}, & \text{if } \fin_i^t \preceqbc \ch_i^t\trunc{\mu},\\
    \fin_i^t, & \text{otherwise}.
  \end{cases}
\]
\end{definition}

\begin{proposition}[Ledger prefix property]\label{prop:ledger-prefix}
For every node $i$ honest at time $t$ and every bounded-available depth
$\mu$, $\fin_i^t \preceqbc (\ba_\mu)_i^t$.

Status: proved (TFL book), of a designed-but-unspecified construction.
\end{proposition}

\begin{proof}
In the first case of Definition~\ref{def:ba} the relation is the case
condition; in the second the two chains are equal. The book states the
property as a theorem and proves it by construction of $\ba_\mu$ (section
``Locally bounded-available chain'').
\end{proof}

\begin{lemma}[Local ba-depth]\label{lem:ba-depth}
In every execution of Crosslink~2, for every bounded-available depth
$\mu \le \sigma$ and every node $i$ honest at time $t$, there is a time
$r \le t$ with $(\ba_\mu)_i^t \preceqbc \ch_i^r\trunc{\mu}$.

Status: proved (TFL book), of a designed-but-unspecified construction.
\end{lemma}

\begin{proof}
The proof is the book's (section ``Locally bounded-available chain''). In
the first case of Definition~\ref{def:ba}, $(\ba_\mu)_i^t =
\ch_i^t\trunc{\mu}$; take $r = t$. Otherwise $(\ba_\mu)_i^t = \fin_i^t$, and
Lemma~\ref{lem:fin-depth} gives $r \le t$ with $\fin_i^t \preceqbc
\ch_i^r\trunc{\sigma}$. Since $\mu \le \sigma$, $\ch_i^r\trunc{\sigma} =
(\ch_i^r\trunc{\mu})\trunc{\sigma - \mu} \preceqbc \ch_i^r\trunc{\mu}$
(Definition~\ref{def:trunc}).
\end{proof}

% ledger_checks.py: 4. fallback arm needed (mu = sigma) after a switch to a fork of sigma - 1 blocks from fin, Prefix Consistency at depth sigma holding, for every 2 <= sigma <= 10: ok
\begin{remark}[The fallback arm]\label{rem:ba-fallback}
The second case of Definition~\ref{def:ba} is needed even when every safety
assumption holds. The book argues it for proof of work (admonition ``Why is
the `otherwise' case in the definition of $(\ba_\mu)_i^t$ necessary?''). After
the best chain switches to a fork on which the difficulty is higher than on
the chain of the previous snapshot, the new chain can reach a higher score
than the old one with fewer than $\sigma$ blocks above the fork point, so
$\ch_i^t\trunc{\mu}$ need not descend from $\fin_i^t$, more often so when
$\mu = \sigma$. Using $\fin_i^t$ in that case is safe, the book argues,
because although $\fin_i^t$ is not a prefix of $\ch_i^t\trunc{\mu}$, so that
its tip has at most $\mu$ confirmations on $\ch_i^t$
(Remark~\ref{rem:depth-conf}; the book's ``$\mu$ confirmations'' means depth
$\mu$), it carries at least the required work above it: score is accumulated
work (Remark~\ref{rem:bc-zcash}; \emph{Consensus Guide}, \S``Work, and the best
chain''), and Zcash's difficulty adjustment damps and clamps the effect of
timestamps, so an adversary gains little by manipulating it (\emph{Consensus
Guide}, \S``Difficulty adjustment''). The argument rests on that damping, not
on the length of the averaging window, which ZIP 218, ``Averaging window'',
changes under NU7; no window length is used here. The \emph{Consensus
Guide}'s Remark ``What per-block retargeting changes'' records that window,
damping and clamp do not prove that two forked branches keep close targets,
so the safety of the second case remains an argument. The book's wording, a
switch of $\fin_i$ ``to a different fork'', is loose: by
Lemma~\ref{lem:fin-linear}, $\fin_i$ never moves to a chain that conflicts
with it; the best chain switches. A script exhibits, for every
$2 \le \sigma \le 10$ with $\mu = \sigma$, a switch of the best chain to a
fork of $\sigma - 1$ blocks above $\fin_i^t$, with Prefix Consistency at
depth $\sigma$ (Definition~\ref{def:pc}) holding, after which
$\ch_i^t\trunc{\mu}$ is a strict prefix of $\fin_i^t$ and
$(\ba_\mu)_i^t = \fin_i^t$ through the second case only.

Status: argued (TFL book).
\end{remark}

\begin{remark}[Ledgers as chains]\label{rem:ledger-names}
The ledgers of Crosslink~2 are the bc-chains $\fin_i^t$ and $(\ba_\mu)_i^t$.
Since no snapshot is sanitised, they need not be expressed as transaction
logs (section ``Parameters''). Where the ledger names of the paper and of the
TFL book are used, $\LOGfin[i]^t$ denotes the transaction sequence of
$\fin_i^t$, the transactions of its blocks from $\Obc$ to its tip in chain order
and, within a block, in block order, and $\LOGba[i]^t$ that of $(\ba_\mu)_i^t$
for node $i$'s choice of $\mu$. The book's security analysis below its marker
``Not updated for Crosslink 2'' and its chapter ``Potential Changes to
Crosslink'' use two other ledger forms, from Crosslink~1, whose finalised ledger
is read from the snapshots of $\LF(\ch_i^t\trunc{\sigma})$; they are not the
ledgers of Crosslink~2 and are not used in this volume.
\end{remark}

\subsection{Syncing and checkpoints}\label{sec:led-sync}

\begin{remark}[Syncing gate]\label{rem:sync}
The book recommends, as a rule for nodes and not as a consensus rule, that
a node embed a checkpoint bft-block $\Bckpt$, fixed in each release, and that
node $i$ expose $\fin_i^t$ and $(\ba_\mu)_i^t$ to its clients only once it is
\emph{synced}, that is, when
\begin{enumerate}[label=(\roman*)]
\item $\Bckpt \preceqbft \LF(\ch_i^t)$,
\item $\snapshot(\Bckpt) \preceqbc \fin_i^t$, and
\item the timestamp of the tip of $\fin_i^t$ is within some threshold of the
  current time
\end{enumerate}
(section ``Syncing and checkpoints''). The gate carries the book's withholding
clauses for the two chains, which Definitions~\ref{def:fin}
and~\ref{def:ba} leave out: neither chain is exposed to clients before the
node is synced (sections ``Locally finalized chain'' and ``Locally
bounded-available chain''). By condition~(i), an observation of this volume,
a node is not synced while its best tip cites a context that does not
descend from $\Bckpt$, such as one on a bft-chain forged with the keys of BFT
epochs before the checkpoint; the book does not state this. The checkpoint and
the threshold are an open problem: neither the book nor Shielded Labs'
\emph{Zcash Crosslink Design Overview} (draft 10, 14 September 2026) fixes them.
\end{remark}

\subsection{Computation of the finalised prefix}\label{sec:led-compute}

\begin{construction}[Node update]\label{def:node-update}
On each change of its best chain, at time $t$ with new tip $H = \ch_i^t$,
node $i$ performs the following steps.
\begin{enumerate}[label=(\arabic*)]
\item It resolves $H.\contextbft$ to a bft-block-valid block, with its
  ancestry, the $\sigma$-block tail that each ancestor carries and,
  recursively, the contexts of those bc-blocks, in the validation order of
  Construction~\ref{def:honest-voting}; validating $H$ under the Valid
  context rule has done so.
\item It computes $\LF(H) = \bftlastfinal(H.\contextbft)$.
\item It computes the final snapshot $\snapshot(\LF(H))$ of $H$.
\item It computes $\candidate(H)$, the lower of the final snapshot and
  $H\trunc{\sigma}$ (Lemma~\ref{lem:candidate-lfs}, since $H$ is bc-valid).
\item It applies the update rule of Definition~\ref{def:fin}.
\item It computes $(\ba_\mu)_i^t$ (Definition~\ref{def:ba}).
\item It exposes $\fin_i^t$ and $(\ba_\mu)_i^t$ only through the syncing
  gate of Remark~\ref{rem:sync}.
\end{enumerate}
\end{construction}

A decision of $\Pibft$ by itself changes nothing at a node: $\fin_i$ changes
only at changes of $\ch_i$, and its new value depends on $\Pibft$ only through
$\LF(\ch_i^t)$. By Lemma~\ref{lem:lf-monotone}, $\LF(\ch_i^t)$ extends the
last final context of every block of $\ch_i^t$, so a decided bft-block moves
$\fin_i$ only when a bc-block whose last final context extends it is, or is
an ancestor of, the best tip.

% ledger_checks.py: 3. candidate below the truncated tip (47850 cases); with the carried headers on the tip's chain, lag = finality depth >= sigma + 1 and >= sigma + 2 confirmations (7330 cases); lag sigma + 1 blocks, sigma + 2 confirmations, attained, for every sigma <= 20; genesis context excluded (lag height(H) <= sigma for a low tip): ok
\begin{proposition}[Finalisation lag]\label{prop:lag}
Let $H$ be a bc-valid block.
\begin{enumerate}[label=(\roman*)]
\item The candidate satisfies $\candidate(H) \preceqbc H\trunc{\sigma}$.
\item If $\LF(H) \neq \Obft$ and the headers that $\LF(H)$ carries are
  headers of ancestors of $H$, then $\candidate(H) = \snapshot(\LF(H))$ and
  \[
    \mathsf{height}(H) - \mathsf{height}(\candidate(H)) = \finalitydepth(H) \ge
    \sigma + 1.
  \]
  Hence at every update of a node with best tip $H$ that sets $\fin$ to
  $\candidate(H)$ under these hypotheses, $\fin$ trails the tip by at least
  $\sigma + 1$ blocks, and the distance from the tip to $\fin$ equals the
  finality depth.
\item The bound is attained: if $T = H\trunc[\bc]{1}$ has at least
  $\sigma + 1$ blocks and $\LF(H)$ carries the $\sigma$-block tail of $T$,
  then $\mathsf{height}(H) - \mathsf{height}(\candidate(H)) = \sigma + 1$. This
  is the case when an honest proposer whose best tip is $T$ carries the
  $\sigma$-block tail of $T$ (Construction~\ref{def:honest-proposal}) and the
  bft-block formed from its proposal is the last block final in the context that
  $H$ cites (Definition~\ref{def:final-context}).
\end{enumerate}

Status: proved (this volume), of a designed-but-unspecified construction.
\end{proposition}

\begin{proof}
(i) A last common ancestor is a prefix of each argument
(Definition~\ref{def:lca}).

(ii) The block $\LF(H)$ is bft-block-valid (the Valid context rule and
Definition~\ref{def:blf}) and is not $\Obft$, so by the Tail Confirmation
rule its field $\headersbc$ is the $\sigma$-block tail of a bc-valid chain
$T'$ of at least $\sigma + 1$ blocks, $T'$ being the block of the newest
carried header. By the paragraph after Definition~\ref{def:snapshot},
$\snapshot(\LF(H)) = T'\trunc{\sigma}$, of height $\mathsf{height}(T') -
\sigma$. By hypothesis $T' \preceqbc H$, and $T' \neq H$: block $H$ commits by
hash to $H.\contextbft$, which has $\LF(H)$ as a prefix and commits to it
through parent links, and $\LF(H)$ commits to the header of $T'$, so $T' = H$
would be a cycle of commitments, excluded as in
Remark~\ref{rem:validity-recursion}. Hence $\mathsf{height}(T') \le
\mathsf{height}(H) - 1$ and
\[
  \mathsf{height}(\snapshot(\LF(H))) \le \mathsf{height}(H) - \sigma - 1.
\]
Both $\snapshot(\LF(H)) \preceqbc T' \preceqbc H$ and $H\trunc{\sigma} \preceqbc
H$ are prefixes of $H$, and $\mathsf{height}(H\trunc{\sigma}) =
\mathsf{height}(H) - \sigma$ because $\mathsf{height}(H) \ge \mathsf{height}(T')
+ 1 \ge \sigma + 1$; so the final snapshot is a strict prefix of
$H\trunc{\sigma}$ (Lemma~\ref{lem:linear-prefix}), and
Lemma~\ref{lem:candidate-lfs} gives $\candidate(H) = \snapshot(\LF(H))$. Then
$\mathsf{height}(H) - \mathsf{height}(\candidate(H)) = \finalitydepth(H) \ge
\sigma + 1$. At an update that sets $\fin_i^t := \candidate(\ch_i^t)$ with
$\ch_i^t = H$, the distance from the tip to $\fin_i^t$ is this number.

(iii) Here the newest carried header is that of $T$, so by the computation
of~(ii), with $T' = T$, $\mathsf{height}(\snapshot(\LF(H))) = \mathsf{height}(T)
- \sigma = \mathsf{height}(H) - \sigma - 1$ and $\candidate(H) =
\snapshot(\LF(H))$; the distance is $\sigma + 1$. In the case described, the
proposal carries the $\sigma$-block tail of $T$ and its block is $\LF(H)$.
\end{proof}

The hypothesis $\LF(H) \neq \Obft$ cannot be dropped: if $\LF(H) = \Obft$,
then $\candidate(H) = \lca(\Obc, H\trunc{\sigma}) = \Obc$, at distance
$\mathsf{height}(H)$ from $H$, which is below $\sigma + 1$ whenever
$\mathsf{height}(H) \le \sigma$. A script checks~(i) for final snapshots on and
off $H$'s chain, (ii) and~(iii) for every $\sigma \le 20$, and the excluded
case.
The latency of $\Pibft$ and contexts older than the newest decision add to
the lag. The proposition counts blocks, not confirmations
(Remark~\ref{rem:depth-conf}). It corrects the book's remark that the
finality depth and the distance from a node's tip to $\fin$ differ only just
after a reorganisation (admonition ``Explain the definition of
finality-depth.''): by~(ii) they can differ only at an update that keeps
$\fin$, or when the headers that $\LF(H)$ carries are not all ancestors of
$H$, and the second case needs no reorganisation, as in the construction of
Proposition~\ref{prop:short-L}.

\begin{remark}[Latency estimate of the book]\label{rem:latency}
The book estimates that the finalised view lies ``about $\mu + 1 + \sigma$
blocks back'', against $\sigma$ blocks for snap-and-chat at the same
confirmation depth, a relative increase of ``at most slightly more than a
factor of 2'' for $\mu = \sigma$ (chapter ``Security Analysis of Crosslink
2'', section ``Finalization latency''). Both figures concern a ledger form of
the book's Crosslink~1 analysis, which reads the BFT context from the block
at depth $\mu$ of the best chain, and are stale for Crosslink~2. In
Crosslink~2, at an update that sets $\fin$ to $\candidate(H)$, $\fin$ lies at
least $\min\{\sigma, \mathsf{height}(H)\}$ blocks below $H$
(Proposition~\ref{prop:lag}(i)). Under the hypotheses of~(ii) it lies at
least $\sigma + 1$ blocks below, attained when a decision is final in context
at once and cited by the next block, so the cost over snap-and-chat's
$\sigma$ is about one block plus the latency of $\Pibft$. No source gives a
validated latency bound for Crosslink~2.
\end{remark}

Figure~\ref{fig:two-ledgers} draws the two chains at an update in Stalled
Mode. Every block $B \preceqbc H$ with $\mathsf{height}(B) >
\mathsf{height}(\snapshot(\LF(H))) + L$ is stalled: $B$ is
bc-valid as a prefix of $H$, and by Lemma~\ref{lem:lf-monotone}
and~\eqref{eq:cons-snap-linear} its final snapshot is a prefix of that of
$H$, so $\finalitydepth(B) > L$ and the Finality depth rule applies.

\begin{figure}[ht]
\centering
\begin{tikzpicture}[x=11mm, y=10mm,
  blk/.style={draw, rectangle, minimum width=5.5mm, minimum height=5.5mm,
    inner sep=0pt, fill=white},
  fin/.style={blk, fill=tierin!30}, stl/.style={blk, fill=tierspend!30},
  dts/.style={inner sep=1pt}, lnk/.style={-{Stealth[length=2mm]}, thick},
  lbl/.style={font=\footnotesize}, bar/.style={line width=2.5pt},
]
\node[blk] (g) at (0,0) {};
\node[dts] (d1) at (1,0) {$\cdots$};
\node[fin] (s) at (2,0) {};
\node[dts] (d2) at (3,0) {$\cdots$};
\node[blk] (l) at (4,0) {};
\node[stl] (f) at (5,0) {};
\node[dts] (d3) at (6,0) {$\cdots$};
\node[stl] (ts) at (7,0) {};
\node[dts] (d4) at (8,0) {$\cdots$};
\node[stl] (tm) at (9,0) {};
\node[dts] (d5) at (10,0) {$\cdots$};
\node[stl] (H) at (11,0) {};
\draw[lnk] (d1) -- (g);
\draw[lnk] (s) -- (d1);
\draw[lnk] (d2) -- (s);
\draw[lnk] (l) -- (d2);
\draw[lnk] (f) -- (l);
\draw[lnk] (d3) -- (f);
\draw[lnk] (ts) -- (d3);
\draw[lnk] (d4) -- (ts);
\draw[lnk] (tm) -- (d4);
\draw[lnk] (d5) -- (tm);
\draw[lnk] (H) -- (d5);
\node[lbl, below=1mm of g] {$\Obc$};
\node[lbl, below=1mm of s, align=center]
  {$\snapshot(\LF(H))$\\$= \candidate(H)$};
\node[lbl, below=1mm of ts] {$H\trunc{\sigma}$};
\node[lbl, below=1mm of tm] {$H\trunc{\mu}$};
\node[lbl, below=1mm of H] {$H$};
% brackets above the chain
\draw (1.7,0.45) -- (1.7,0.6) -- (4.3,0.6) -- (4.3,0.45);
\node[lbl, above] at (3,0.6) {$L$};
\draw (6.7,0.45) -- (6.7,0.6) -- (11.3,0.6) -- (11.3,0.45);
\node[lbl, above] at (9,0.6) {$\sigma$};
\node[lbl, text=tierspend] at (5.5,0.75) {stalled};
\draw (1.7,1.25) -- (1.7,1.4) -- (11.3,1.4) -- (11.3,1.25);
\node[lbl, above] at (6.5,1.4)
  {$\finalitydepth(H) \ge \sigma + 1$};
% the two chains of node i
\draw[bar, tierin] (0,-1.6) -- (2,-1.6);
\node[lbl, anchor=west] at (2.1,-1.6) {$\fin_i^t$};
\draw[bar, tierrecv] (0,-2.0) -- (9,-2.0);
\node[lbl, anchor=west] at (9.1,-2.0) {$(\ba_\mu)_i^t$};
\end{tikzpicture}
\caption{The two chains of node $i$ at an update with best tip
$H = \ch_i^t$ in Stalled Mode, when the headers that $\LF(H)$ carries are
ancestors of $H$ and the update sets $\fin_i^t$ to $\candidate(H)$. Arrows
point from child to parent. The final snapshot $\snapshot(\LF(H))$ is the
lower of itself and $H\trunc{\sigma}$, hence $\candidate(H)$; $\fin_i^t$
(blue) ends there, $\finalitydepth(H) \ge \sigma + 1$ blocks below $H$. The
chain $(\ba_\mu)_i^t$ (green) is $H\trunc{\mu}$, and $\fin_i^t$ is a prefix of
it. The blocks more than $L$ above the final snapshot, up to $H$, are stalled
(red). Positions are symbolic and represent no value of $\sigma$, $\mu$ or
$L$.}
\label{fig:two-ledgers}
\end{figure}

\section{Safety}\label{sec:safety}

This section states the safety goal of Crosslink~2 and proves it from either
of two premises: Prefix Agreement of $\Pibc$ at depth $\sigma$, which needs
no property of $\Pibft$, and Final Agreement of $\Pibft$, which needs the
Linearity rule. It records which rules the two proofs use, the attack that
defeats both premises, and what is known of the safety of the
bounded-available chain. Every object of the section belongs to the
construction of Definition~\ref{def:crosslink}, which is designed but
unspecified in the classes of Definition~\ref{def:classes}; the theorems are
implications proved of that construction, and no rule they use is
specified.

\subsection{Assured Finality}\label{sec:saf-af}

\begin{definition}[Assured Finality]\label{def:af}
Assured Finality holds in an execution of Crosslink~2 iff for all times $t$,
$u$ and all nodes $i$ honest at $t$ and $j$ honest at $u$, $i = j$
permitted, the chains $\fin_i^t$ and $\fin_j^u$ agree; that is, iff
Agreement on the view $(i, t) \mapsto \fin_i^t$ holds
(Definition~\ref{def:agree-view}).
\end{definition}

Assured Finality is the main safety goal of the TFL book for Crosslink~2,
which intends it to hold under assumptions on either sub-protocol (chapter
``The Crosslink 2 Construction'', section ``Locally finalized chain''). For
a single node the goal concerns one time series of bc-chains.

\begin{proposition}[The single-node case]\label{prop:af-single}
Let $S$ be a time series of bc-chains. If $S$ is bc-linear
(Definition~\ref{def:prefix}), then $S^t$ and $S^u$ agree for all times $t$
and $u$. The converse fails: some series whose values agree pairwise is not
bc-linear. Hence Lemma~\ref{lem:fin-linear} implies the case $i = j$ of
Assured Finality, and that case does not imply bc-linearity of a general
series.

Status: proved (this volume).
\end{proposition}

% safety_checks.py: 3. a pairwise-agreeing series need not be linear: ok
\begin{proof}
If $t \le u$ then $S^t \preceqbc S^u$, and if $u \le t$ then
$S^u \preceqbc S^t$; either way the two agree. For the converse let $A$ be a
bc-chain other than $\Obc$, and let $S^0 = A$ and $S^t = A\trunc{1}$ for
$t > 0$. Any two values agree, since $A\trunc{1} \preceqbc A$, and the
series is not bc-linear, since $A \mathrel{\not\preceq}_{\bc} A\trunc{1}$. If
node $i$ is honest at $t$ and at $u$, it has followed the update rule up to
$\max\{t, u\}$ (Definition~\ref{def:exec}), so Lemma~\ref{lem:fin-linear}
makes $\fin_i$ bc-linear up to that time, and the first part makes
$\fin_i^t$ and $\fin_i^u$ agree.
\end{proof}

The TFL book states that the restriction of Assured Finality to the case
$i = j$ ``is equivalent to'' Local finalization linearity (chapter ``The
Crosslink 2 Construction'', section ``Locally finalized chain''). By
Proposition~\ref{prop:af-single} only one direction holds for a general
series, as a script check of the two-element series confirms. For $\fin$
both hold, because $\updatefin$ makes each series $\fin_i$ bc-linear by
construction; the book's inference from Assured Finality to Local
finalization linearity at every honest node uses the false direction, and
its conclusion holds for $\fin$ by Lemma~\ref{lem:fin-linear} alone.

\subsection{The Prefix Agreement route}\label{sec:saf-pa}

\begin{theorem}[Prefix Agreement implies Assured Finality]\label{thm:pa-af}
In every execution of Crosslink~2 in which $\Pibc$ has Prefix Agreement at
depth $\sigma$ (Definition~\ref{def:pa}), Assured Finality holds.

Status: proved (TFL book), of a designed-but-unspecified construction.
\end{theorem}

The TFL book states and proves the theorem in chapter ``Security Analysis of
Crosslink 2'', section ``Safety argument''; its proof follows.

\begin{proof}
Let node $i$ be honest at $t$ and node $j$ honest at $u$. By
Lemma~\ref{lem:fin-depth} there are times $r \le t$ and $s \le u$ with
\[
  \fin_i^t \preceqbc \ch_i^r\trunc{\sigma} \quad\text{and}\quad
  \fin_j^u \preceqbc \ch_j^s\trunc{\sigma},
\]
and node $i$ is honest at $r$ and node $j$ at $s$ (Definition~\ref{def:exec}).
Prefix Agreement at $\sigma$, applied at $(i, r)$ and $(j, s)$, makes
$\ch_i^r\trunc{\sigma}$ and $\ch_j^s\trunc{\sigma}$ agree; by symmetry let the
first be a prefix of the second. Then $\fin_i^t \preceqbc \ch_j^s\trunc{\sigma}$
by transitivity, and $\fin_j^u \preceqbc \ch_j^s\trunc{\sigma}$, so $\fin_i^t$
and $\fin_j^u$ agree by Lemma~\ref{lem:linear-prefix}.
\end{proof}

The proof uses Prefix Agreement of $\Pibc$ at $\sigma$ and $\updatefin$,
through Lemma~\ref{lem:fin-depth}, and nothing about $\Pibft$. For Zcash the
premise is an assumption, not a result (Remark~\ref{rem:pow-prefix}); the
nearest statement in the lower volumes, the persistence clause of Assumption
``Persistence and liveness'' (\emph{FlyClient Guide}, \S``The
variable-difficulty backbone model''), is a prefix property of the backbone
model and is itself adopted there as an assumption.

\begin{corollary}[Safety under subversion of the BFT protocol]
\label{cor:bft-subversion}
In every execution in which $\Pibc$ has Prefix Agreement at $\sigma$,
Assured Finality holds whatever the holders of voting units do, including
executions in which every voting unit is cast non-honestly
(Definition~\ref{def:nonhonest}) and signatures are forged.

Status: proved (this volume), of a designed-but-unspecified construction.
\end{corollary}

\begin{proof}
The proof of Theorem~\ref{thm:pa-af} uses Prefix Agreement of $\Pibc$,
Lemma~\ref{lem:fin-depth} and Lemma~\ref{lem:linear-prefix}. The proof of
Lemma~\ref{lem:fin-depth} uses $\updatefin$, $\ch_i^0 = \Obc$, bc-validity of
honest best chains (Definition~\ref{def:bc-protocol}), whose Valid context
rule makes $\LF(H)$ a bft-block-valid block, and the relation
$\candidate(H) \preceqbc H\trunc{\sigma}$, which holds by
Definitions~\ref{def:candidate} and~\ref{def:lca} whatever block
$\snapshot(\LF(H))$ is. None of these refers to the behaviour of $\Pibft$:
validity in $\Pibft$ is objective (Remark~\ref{rem:objective}), not a matter
of the conduct of the holders of voting units; the proof uses neither Final
Agreement (Definition~\ref{def:fa}), nor Assumption~\ref{ass:one-third}, nor
Assumption~\ref{ass:ballots}, nor Assumption~\ref{ass:suf-cma}. The theorem
therefore holds in every execution with the premise, whatever votes are cast and
whatever signatures exist.
\end{proof}

The corollary makes precise the consequence that the TFL book argues from
the fact recorded in Remark~\ref{rem:objective}: the finalised chain stays
safe under a complete break of $\Pibft$, its signature scheme included,
while $\Pibc$ is safe (chapter ``The Crosslink 2 Construction'', section
``$\Pibft$ changes from $\Piorigbft$''). It meets item~(1) of
Remark~\ref{rem:sac-deficiencies}, with Prefix Agreement of $\Pibc$ at
$\sigma$ in place of the condition that the adversary holds under half of
the hash rate. Whether $\Pibc$ keeps Prefix Agreement while $\Pibft$ is
subverted is a separate question, which the corollary does not address:
validity in $\Pibc$ depends on $\Pibft$ (Remark~\ref{rem:bidirectional}),
and transfer of the prefix properties from $\Piorigbc$ to $\Pibc$ is an open
problem, recorded with the liveness of the sub-protocols.

\subsection{The Final Agreement route}\label{sec:saf-fa}

\begin{theorem}[Final Agreement implies Assured Finality]\label{thm:fa-af}
In every execution of Crosslink~2 in which $\Pibft$ has Final Agreement
(Definition~\ref{def:fa}), Assured Finality holds.

Status: proved (this volume), of a designed-but-unspecified construction.
\end{theorem}

\begin{proof}
Let node $i$ be honest at $t$ and node $j$ honest at $u$. The chain $\Obc$
is a prefix of every bc-chain, so the claim holds if $\fin_i^t = \Obc$ or
$\fin_j^u = \Obc$; let both differ from $\Obc$. Since $\fin_i^0 = \Obc$,
$\updatefin$ (Definition~\ref{def:fin}) has set $\fin_i$ to a candidate at
some time at most $t$; let $r \le t$ be the last such time, so that
$\fin_i^t = \candidate(\ch_i^r)$. By Definitions~\ref{def:candidate}
and~\ref{def:lca},
\[
  \fin_i^t \preceqbc \snapshot(\LF(\ch_i^r)).
\]
Node $i$ is honest at $r$, so $\ch_i^r$ is valid in $\Pibc$
(Definition~\ref{def:bc-protocol}). It is not $\Obc$, whose candidate is
$\Obc$ by the Genesis bc-block rule; so the Valid context rule
(Definition~\ref{def:bc-validity}) makes $C_i = \ch_i^r.\contextbft$ a valid
bft-block. Node $i$ observed it, while honest, in validating $\ch_i^r$, so
the block $C_i$ is in honest view (Definition~\ref{def:bft-block}), and
$\LF(\ch_i^r) = \bftlastfinal(C_i)$ (Definition~\ref{def:lf}). In the same
way there are a time $s \le u$ and a valid bft-block $C_j$ in honest view
with
\[
  \fin_j^u = \candidate(\ch_j^s) \preceqbc \snapshot(\LF(\ch_j^s))
  \quad\text{and}\quad \LF(\ch_j^s) = \bftlastfinal(C_j).
\]
Final Agreement makes $\LF(\ch_i^r)$ and $\LF(\ch_j^s)$ agree; by symmetry
let $\LF(\ch_i^r) \preceqbft \LF(\ch_j^s)$. The block $\LF(\ch_j^s)$ is
valid (Definition~\ref{def:blf}), so by the Parent rule every bft-block $D$
with $\LF(\ch_i^r) \precbft D \preceqbft \LF(\ch_j^s)$ is valid, and the
Linearity rule (Definition~\ref{def:bft-validity}) gives
$\snapshot(D\trunc{1}) \preceqbc \snapshot(D)$ for each. By transitivity
along this path,
\[
  \snapshot(\LF(\ch_i^r)) \preceqbc \snapshot(\LF(\ch_j^s)).
\]
Both $\fin_i^t$ and $\fin_j^u$ are therefore prefixes of
$\snapshot(\LF(\ch_j^s))$, and Lemma~\ref{lem:linear-prefix} makes them
agree.
\end{proof}

The proof uses Final Agreement of $\Pibft$, the Genesis bc-block, Valid
context, Parent and Linearity rules and $\updatefin$. It uses neither $\sigma$
nor Prefix Agreement, nor the Extension, Last Final Snapshot or Finality depth
rules.
The TFL book's printed proof of this theorem (chapter ``Security Analysis of
Crosslink 2'', section ``Safety argument'') is a verbatim copy of its proof
of Theorem~\ref{thm:pa-af}; the admonition ``What does the Linearity rule
do?'' of its construction chapter sketches the route taken above.

% safety_checks.py: 1. Final Agreement and Linearity: all fin values agree (1500 trials, 4500 node runs, 1949 with forks deeper than sigma): ok
A randomised check by script agrees with the theorem: in $1500$ executions
with one bft-chain obeying the Linearity rule and every block final in its
own context, over $4500$ node runs, $1949$ of them with forks deeper than
$\sigma$, so that Prefix Agreement at $\sigma$ fails, all values of $\fin$
agree. The Linearity rule cannot be dropped from the theorem.

\begin{proposition}[Necessity of the Linearity rule]\label{prop:fa-needs-lin}
Let the Linearity rule be removed from Definition~\ref{def:bft-validity}.
For every $\sigma$ there is an execution in which $\Pibft$ has Final
Agreement, the best chains of two honest nodes $i$ and $j$ satisfy the Valid
context, Extension and Last Final Snapshot rules, the voters of every
bft-block they cite meet the Confirmed best-chain criterion
(Construction~\ref{def:honest-voting}), and $\fin_i$ and $\fin_j$ conflict.

Status: proved (this volume), of a designed-but-unspecified construction.
\end{proposition}

\begin{proof}
Let $\Piorigbft$ be a protocol in which every valid bft-block is final in
its own context, $\bftlastfinal(B) = B$, as in a protocol that decides each
block individually; the identity on valid blocks meets Definition~\ref{def:blf}.
Let $X$ and $Y$ be two bc-branches that fork at $\Obc$, with blocks $X_1,
\ldots, X_{\sigma+2}$ and $Y_1, \ldots, Y_{\sigma+2}$, the index being the
height (Figure~\ref{fig:linearity-counterexample}).

\emph{BFT side.} Let $B_0 = \Obft$; let $B_1$ be a bft-block with parent
$B_0$ whose field $\headersbc$ is the $\sigma$-block tail of the chain
$X_{\sigma+1}$, the headers of $X_2, \ldots, X_{\sigma+1}$; and let $B_2$ be
a bft-block with parent $B_1$ whose field $\headersbc$ is the
$\sigma$-block tail of $Y_{\sigma+1}$. Then $\snapshot(B_1) = X_1$ and
$\snapshot(B_2) = Y_1$ (Definition~\ref{def:snapshot}). Both blocks satisfy
the Tail Confirmation rule. The block $B_1$ also satisfies the Linearity
rule, since $\snapshot(B_0) = \Obc \preceqbc X_1$; the block $B_2$ violates
it, since $X_1$ and $Y_1$ conflict, and is valid once that rule is removed.
The only bft-blocks are $B_0$, $B_1$ and $B_2$, on one chain, so any two
values of $\bftlastfinal$ agree and $\Pibft$ has Final Agreement.

\emph{Voting.} At least two thirds of the units of $B_1$'s BFT epoch are
cast by validators whose best chain is $X_{\sigma+1}$, and at least two
thirds of those of $B_2$'s by validators whose best chain is $Y_{\sigma+1}$.
Each such validator meets the Confirmed best-chain criterion, since
$X_1 = X_{\sigma+1}\trunc{\sigma}$ and $Y_1 = Y_{\sigma+1}\trunc{\sigma}$. Let
$B_2$ be of a later BFT epoch than $B_1$, and let $\score(X_{\sigma+1}) <
\score(Y_{\sigma+1})$, which Definition~\ref{def:bc-protocol} permits. A network
partition that hides $Y$ from the validators of $B_1$ and from node $i$ permits
this. The validators of $B_2$ and node $j$ validate $B_1$, and so obtain $X_1,
\ldots, X_{\sigma+1}$, but keep $Y$ as their best chain. It breaks Prefix
Consistency (Definition~\ref{def:pc}), since $X_{\sigma+1}\trunc{\sigma}$ is not
a prefix of the later best chain $Y_{\sigma+1}$ of an honest validator; no
unit is cast non-honestly, so Assumption~\ref{ass:one-third} holds.

\emph{Best chains.} Node $i$ is honest, and its best chain grows from $\Obc$
through $X_1, \ldots, X_{\sigma+2}$; the blocks $X_1, \ldots, X_{\sigma+1}$
cite $\Obft$, and the tip $X_{\sigma+2}$ cites $B_1$. Node $j$ follows $Y$
in the same way, with $Y_{\sigma+2}$ citing $B_2$. On $X$ the Valid context
rule holds, $\Obft$ and $B_1$ being valid; the Extension rule holds, since
$\LF(X_k) = \Obft$ for $k \le \sigma+1$ and $\Obft \preceqbft B_1 =
\LF(X_{\sigma+2})$; the Last Final Snapshot rule holds, since $\snapshot(\Obft)
= \Obc$ and $X_1 \preceqbc X_{\sigma+2}$; and the headers carried by $B_1$ are
those of strict ancestors of $X_{\sigma+2}$, so no block commits to its own
hash. The same holds on $Y$.
The Finality depth rule is met by letting each block whose finality depth
exceeds $L$ be a stalled block (Remark~\ref{rem:producible}), which changes
no value of $\fin$.

\emph{Finalised chains.} For $k \le \sigma+1$, $\candidate(X_k) = \lca(\Obc,
X_k\trunc{\sigma}) = \Obc$, so $\fin_i$ stays $\Obc$; at the tip,
\[
  \candidate(X_{\sigma+2}) = \lca(X_1, X_2) = X_1,
\]
of which $\Obc$ is a prefix, so $\updatefin$ sets $\fin_i = X_1$. In the same
way $\fin_j = Y_1$. The blocks $X_1$ and $Y_1$ conflict.
\end{proof}

% safety_checks.py: 2. without Linearity: fin values conflict at height 1, tips at height sigma + 2, for every sigma <= 20: ok
A script checks the construction for every $\sigma \le 20$: the two values
of $\fin$ conflict at height~$1$, with the tips at height $\sigma + 2$. In
the execution Prefix Agreement at $\sigma$ fails, as Theorem~\ref{thm:pa-af}
requires, since $X_{\sigma+2}\trunc{\sigma} = X_2$ and
$Y_{\sigma+2}\trunc{\sigma} = Y_2$ conflict. The proposition is the only
result in any source that shows a rule of Crosslink~2 necessary for Assured
Finality.

\begin{figure}[htbp]
\centering
\begin{tikzpicture}[
  bcb/.style={draw, minimum width=11mm, minimum height=6mm,
    inner sep=1pt, font=\small},
  bftb/.style={bcb, draw=tierfull, very thick}, finb/.style={bcb, draw=tierin,
  fill=tierin!20, very thick}, par/.style={-{Latex[length=1.6mm]}},
  ctx/.style={-{Latex[length=1.6mm]}, tierrecv, very thick},
  snap/.style={-{Latex[length=1.6mm]}, tierfull, dashed, thick}]
\node[bcb] (O) at (0,0) {$\Obc$};
\node[finb] (X1) at (1.6,1.5) {$X_1$};
\node[bcb] (X2) at (3.0,1.5) {$X_2$};
\node (Xd) at (4.1,1.5) {$\cdots$};
\node[bcb] (Xs) at (5.2,1.5) {$X_{\sigma+1}$};
\node[bcb] (Xt) at (6.7,1.5) {$X_{\sigma+2}$};
\node[finb] (Y1) at (1.6,-1.5) {$Y_1$};
\node[bcb] (Y2) at (3.0,-1.5) {$Y_2$};
\node (Yd) at (4.1,-1.5) {$\cdots$};
\node[bcb] (Ys) at (5.2,-1.5) {$Y_{\sigma+1}$};
\node[bcb] (Yt) at (6.7,-1.5) {$Y_{\sigma+2}$};
\node[bftb] (B0) at (3.9,0) {$B_0$};
\node[bftb] (B1) at (5.6,0) {$B_1$};
\node[bftb] (B2) at (7.3,0) {$B_2$};
\foreach \c/\p in {X1/O, X2/X1, Xd/X2, Xs/Xd, Xt/Xs,
    Y1/O, Y2/Y1, Yd/Y2, Ys/Yd, Yt/Ys, B1/B0, B2/B1}
  \draw[par] (\c) -- (\p);
\draw[ctx] (Xt) -- (B1);
\draw[ctx] (Yt) -- (B2);
\draw[snap] (B1) -- (X1);
\draw[snap] (B2) -- (Y1);
\draw (2.45,1.95) -- (2.45,2.05) -- (5.75,2.05) -- (5.75,1.95);
\node[font=\small, above] at (4.1,2.05)
  {$B_1.\headersbc$, $\sigma$ headers};
\draw (2.45,-1.95) -- (2.45,-2.05) -- (5.75,-2.05) -- (5.75,-1.95);
\node[font=\small, below] at (4.1,-2.05)
  {$B_2.\headersbc$, $\sigma$ headers};
\node[font=\small, above=1mm of X1] {$\fin_i$};
\node[font=\small, below=1mm of Y1] {$\fin_j$};
\node[font=\small, above=1mm of Xt] {$H_i$};
\node[font=\small, below=1mm of Yt] {$H_j$};
\draw[{Latex[length=1.6mm]}-{Latex[length=1.6mm]}, tierspend, thick]
  (X1) -- node[right, font=\small] {conflict} (Y1);
\node[font=\small, anchor=east] at (-0.6,-3.0) {height};
\foreach \x/\h in {0/0, 1.6/1, 3.0/2, 5.2/\sigma+1, 6.7/\sigma+2}
  \node[font=\small] at (\x,-3.0) {$\h$};
\end{tikzpicture}
\caption{The execution of Proposition~\ref{prop:fa-needs-lin}, symbolic in
$\sigma$. The bft-chain $B_0 = \Obft$, $B_1$, $B_2$ (orange outline) has
Final Agreement. The best-chain branches $X$ and $Y$ fork at $\Obc$; $B_1$
carries the headers of $X_2, \ldots, X_{\sigma+1}$ and $B_2$ those of
$Y_2, \ldots, Y_{\sigma+1}$, so the dashed arrows give $\snapshot(B_1) =
X_1$ and $\snapshot(B_2) = Y_1$, and $B_2$ violates the Linearity rule. The
green arrows are the fields $\contextbft$ of the tips $H_i = X_{\sigma+2}$
and $H_j = Y_{\sigma+2}$; every other best-chain block cites $\Obft$. The
finalised chains $\fin_i = X_1$ and $\fin_j = Y_1$ (blue) conflict (red). Solid
black arrows point to parents.}
\label{fig:linearity-counterexample}
\end{figure}

\begin{remark}[Linearity in the book]\label{rem:linearity-book}
The construction chapter of the TFL book calls the Linearity rule ``key to
combining simplicity with strong security properties in Crosslink 2'' and
argues that under Final Agreement it implies Assured Finality ``without
needing to assume any safety property of $\Pibc$''. A sentence marked
``Edit'' in the security analysis, in its section ``Every rule in Crosslink 2
is individually necessary'', written of the Increasing Score rule of
Crosslink~1 and renamed for Crosslink~2, states that ``some rules, e.g.\ the
Linearity rule, only contribute heuristically to security in the analysis so
far''. It holds of the analysis as printed (see the discussion after
Theorem~\ref{thm:fa-af}) and conflicts with the construction chapter;
Theorem~\ref{thm:fa-af} with Proposition~\ref{prop:fa-needs-lin} shows the
rule's contribution proved, not heuristic. The chapter ``Questions About
Crosslink'' opens by asking why there is no bc-block-validity rule
$\snapshot(\LF(H)) \preceqbc H$, which is the Last Final Snapshot rule that
Crosslink~2 then adopted; that chapter predates the construction and is not
read against it.
\end{remark}

\subsection{Consequences of the two routes}\label{sec:saf-either}

\begin{corollary}[Assured Finality under either premise]\label{cor:af-either}
Assured Finality holds in every execution in which $\Pibc$ has Prefix
Agreement at $\sigma$ or $\Pibft$ has Final Agreement.

Status: proved (this volume), of a designed-but-unspecified construction.
\end{corollary}

\begin{proof}
Theorem~\ref{thm:pa-af} gives the first case and Theorem~\ref{thm:fa-af}
the second.
\end{proof}

The finalised chains of two honest nodes can therefore conflict only in an
execution in which both premises fail.

\begin{remark}[Role of the best-chain rules]\label{rem:bc-rules-role}
Neither safety proof uses the Extension, Last Final Snapshot or Finality
depth rule. The proof of Theorem~\ref{thm:pa-af} needs only $\updatefin$,
the Valid context rule, which makes $\candidate(H)$ defined, and
$\candidate(H) \preceqbc H\trunc{\sigma}$; the proof of
Theorem~\ref{thm:fa-af} needs the Genesis bc-block, Valid context, Parent and
Linearity rules and $\updatefin$. The three rules serve other ends. The
Extension and Last Final Snapshot rules, with the Linearity rule, make the
final snapshot of each best-chain block lie on its chain and advance along it
(Lemma~\ref{lem:lf-monotone} and~\eqref{eq:cons-snap-linear}), and the
candidate advance with it (Lemmas~\ref{lem:candidate-lfs}
and~\ref{lem:fd-step}), which serves the spending of finalised outputs
(item~(4) of Remark~\ref{rem:sac-deficiencies}); the Finality depth rule serves
bounded availability, through Stalled Mode (Definition~\ref{def:stalled}).
\end{remark}

\begin{remark}[Residual attack]\label{rem:residual}
Let the network be partitioned, and let the adversary hold more than one
third of the voting units of a BFT epoch. Voting with the honest validators
on each side, it can make each side's last final bft-block consistent with
the fork that side follows, and the finalised chains of honest nodes on the
two sides then conflict. Both premises of Corollary~\ref{cor:af-either} fail
in such an execution: the partition breaks Prefix Agreement, and Final
Agreement fails, the two last final bft-blocks being in conflict; the units
cast on both sides violate Assumption~\ref{ass:one-third}. The TFL book
calls the attack unavoidable, since a best-chain protocol has no mechanism
that maintains finality under partition, and contrasts it with
snap-and-chat, which is unsafe without any partition once $\Pibft$ is
subverted (chapter ``Security Analysis of Crosslink 2'', section ``Safety
argument'').

Status: argued (TFL book).
\end{remark}

\begin{remark}[Necessity of the other rules]\label{rem:rule-necessity}
Beyond Proposition~\ref{prop:fa-needs-lin}, no result shows a rule of
Crosslink~2 necessary for Assured Finality or for bounded availability. The
section ``Every rule in Crosslink 2 is individually necessary'' of the TFL
book's security analysis contains only its statement of intent and the edit
quoted in Remark~\ref{rem:linearity-book}. Witness runs of a model in
Z.~Manian's Quint suite for Crosslink (18 May to 4 June 2026) that is aligned
with the book restate the Last Final Snapshot and Extension rules as
invariants and show that dropping either rule violates its own restatement;
the model is only witness-tested (Remark~\ref{rem:mc-zq}), and that is not a
necessity result.

Status: open problem.
\end{remark}

\subsection{Safety of the bounded-available chain}\label{sec:saf-ba}

\begin{theorem}[Prefix Consistency implies Prefix Consistency of $\ba$]
\label{thm:pc-ba}
Let $\mu \le \sigma$ (Definition~\ref{def:mu}). In every execution in which
$\Pibc$ has Prefix Consistency at depth $\mu$ (Definition~\ref{def:pc}),
$(\ba_\mu)_i^t \preceqbc \ch_j^u$ for all times $t \le u$, every node $i$
honest at $t$ and every node $j$ honest at $u$.

Status: proved (TFL book), of a designed-but-unspecified construction.
\end{theorem}

\begin{proof}
By Lemma~\ref{lem:ba-depth} there is a time $r \le t$ with
$(\ba_\mu)_i^t \preceqbc \ch_i^r\trunc{\mu}$; node $i$ is honest at $r$, and
$r \le u$. Prefix Consistency at $\mu$, applied at $(i, r)$ and $(j, u)$,
gives $\ch_i^r\trunc{\mu} \preceqbc \ch_j^u$, and transitivity of
$\preceqbc$ ends the proof.
\end{proof}

The hypothesis is Prefix Consistency at the depth $\mu$ at which the node
reads $\ba$. Since $C\trunc{k'} \preceqbc C\trunc{k}$ for $k \le k'$, Prefix
Consistency at a depth implies it at every greater depth, so for
$\mu < \sigma$ the hypothesis is stronger than Prefix Consistency at
$\sigma$. The TFL book states the theorem in chapter ``Security Analysis of
Crosslink 2'' below its marker ``Not updated for Crosslink 2 below this
point'', as it does the two theorems on Assured Finality; all three are
written in Crosslink~2 notation.

\begin{proposition}[Retention of the finalised prefix in $\ba$]
\label{prop:ba-keeps-fin}
For times $t \le u$ and a node $i$ honest at $u$, $\fin_i^t \preceqbc
(\ba_\mu)_i^u$. In every execution with Assured Finality, for a node $i$ honest
at $t$ and a node $j$ honest at $u$, the shorter of $\fin_i^t$ and $\fin_j^u$ is
a prefix of both $(\ba_\mu)_i^t$ and $(\ba_\mu)_j^u$.

Status: proved (this volume), of a designed-but-unspecified construction.
\end{proposition}

\begin{proof}
Lemma~\ref{lem:fin-linear} gives $\fin_i^t \preceqbc \fin_i^u$, and
Proposition~\ref{prop:ledger-prefix} gives $\fin_i^u \preceqbc (\ba_\mu)_i^u$;
transitivity gives the first claim.
Under Assured Finality the chains $\fin_i^t$ and $\fin_j^u$ agree, so the
shorter of them is a prefix of both, hence of $(\ba_\mu)_i^t$ and of
$(\ba_\mu)_j^u$ by Proposition~\ref{prop:ledger-prefix} and transitivity.
\end{proof}

The proof uses facts about each node separately, so it holds with a
different depth $\mu$ at each node. The proposition does not claim that
$(\ba_\mu)_i^t$ and $(\ba_\mu)_j^u$ agree: both extend the shorter
finalised chain, and Lemma~\ref{lem:linear-prefix} orders two prefixes of one
chain, not two extensions of one prefix. Nor does it bound how far $\ba_\mu$ can
roll back above $\fin$; the TFL book notes that such rollbacks ``may be of
unbounded length''. The book's version, ``Theorem: $\LOGba$ does not roll back
past the agreed $\LOGfin$'' (section ``Proof of safety for $\LOGba$''), is
written in Crosslink~1 notation, over transaction logs, and transfers
directly to the chains of Crosslink~2.

\begin{definition}[Agreement on $\ba$]\label{def:ba-goal}
Agreement on $\ba$ holds in an execution iff Agreement on the view
$(i, t) \mapsto (\ba_\mu)_i^t$ holds (Definition~\ref{def:agree-view}).
\end{definition}

The TFL book states Agreement on $\ba$ as its security goal for $\ba_\mu$
(chapter ``The Crosslink 2 Construction'', section ``Locally
bounded-available chain'') and does not prove it; Corollary~\ref{cor:ba-agree}
derives it from Prefix Agreement at $\mu$.

\begin{corollary}[Agreement on $\ba$ from Prefix Agreement]\label{cor:ba-agree}
In every execution in which $\Pibc$ has Prefix Agreement at depth $\mu$, and
in particular in every execution in which it has Prefix Consistency at $\mu$,
Agreement on $\ba$ holds; with depths differing between nodes, Prefix
Agreement at the smallest suffices.

Status: proved (this volume), of a designed-but-unspecified construction.
\end{corollary}

\begin{proof}
Let node $i$ be honest at $t$ and node $j$ at $u$, with depths $\mu_i$ and
$\mu_j$ at least $\mu$. Lemma~\ref{lem:ba-depth} gives times $r \le t$ and
$s \le u$, at which the nodes are honest, with $(\ba_{\mu_i})_i^t \preceqbc
\ch_i^r\trunc{\mu_i} \preceqbc \ch_i^r\trunc{\mu}$ and $(\ba_{\mu_j})_j^u
\preceqbc \ch_j^s\trunc{\mu}$ (Definition~\ref{def:trunc}). Prefix Agreement at
$\mu$ makes the two truncated chains agree; by symmetry let the first be a
prefix of the second.
Both values of $\ba$ are then prefixes of $\ch_j^s\trunc{\mu}$, and
Lemma~\ref{lem:linear-prefix} makes them agree. Prefix Consistency at $\mu$
implies Prefix Agreement at $\mu$ (\S\ref{sec:iface-bc}).
\end{proof}

\begin{remark}[Missing results for $\ba$]\label{rem:ba-missing}
Agreement on $\ba$ holds under Prefix Agreement at $\mu$
(Corollary~\ref{cor:ba-agree}). It does not follow from Final Agreement,
which constrains only $\Pibft$: with $\Obft$ the only bft-block, two best chains
that fork at $\Obc$, each with more than $\mu$ blocks above $\Obc$, give
$\fin = \Obc$ at both nodes and conflicting values of $\ba$. Nor does it follow
from Prefix Agreement at $\sigma$ when $\mu < \sigma$: a fork more than $\mu$
and at most $\sigma$ blocks below both tips gives conflicting values of $\ba$.
The safety of the transaction sequence of $\ba_\mu$
(Remark~\ref{rem:ledger-names}) is therefore conditional on Prefix Agreement
at $\mu$, an assumption for Zcash (Remark~\ref{rem:pow-prefix}). The TFL
book's two results on
snapshots in the finalised and bounded-available ledgers,
under Assumption~\ref{ass:one-third} and Prefix Consistency, are written in
Crosslink~1 notation, and its argument for the bounded-available ledger
breaks off (``What we want to show is that, under some conditions on
executions, \ldots''); no Crosslink~2 form exists. No analogue for
Crosslink~2 of the dynamic-availability property P2 of
Definition~\ref{def:ntt-ebbflow} is proved, and no bound on the rollback of
$\ba_\mu$.

Status: open problem.
\end{remark}

\section{Liveness}\label{sec:live}

This section states each liveness claim of the TFL book with its status. The
claims are those of chapter ``Security Analysis of Crosslink 2'', section
``Liveness argument'', with the transfer argument of its section ``Proof of
safety for $\LOGfin$''. Four are argued in prose and none is proved, and the
section proves one necessary condition, $L \ge \sigma + 1$; the
transfer of the prefix properties from $\Piorigbc$, the advance of
finalisation and the choice of $L$ are open problems. The book defines no
liveness property of either sub-protocol: ``live'' below has the meaning of
whatever liveness property is assumed of $\Piorigbc$ or $\Piorigbft$, and for
$\Pibc$ the property argued is growth of the best chain
(Proposition~\ref{prop:bc-live}). The book argues that the dependence of each
sub-protocol's validity on the other (Remark~\ref{rem:bidirectional}) does not
obstruct liveness because the fork choice of $\Pibc$ is unchanged
(Definition~\ref{def:fork-choice}), unlike that of a protocol whose fork choice
follows its finality layer.

\subsection{Liveness of the sub-protocols}\label{sec:live-sub}

\begin{proposition}[Liveness of $\Pibc$]\label{prop:bc-live}
Under the conditions under which the best chain of $\Piorigbc$ grows, the
best chain of $\Pibc$ grows, given that coinbase-only stalled blocks are
always producible.

Status: argued (TFL book).
\end{proposition}

The book's argument is the following. Let the set of bc-blocks that are
both coinbase-only (Definition~\ref{def:coinbase-only}) and stalled
(Definition~\ref{def:params}) be such that a producer can always produce a
block in it (Remark~\ref{rem:producible}). No rule of
Definition~\ref{def:bc-validity} then prevents producing such a block on the
current best chain: the Genesis bc-block rule concerns only the genesis
block; the Valid context, Extension and Last Final Snapshot rules hold when
the new block cites the context of its parent; and the Finality depth rule
always admits a stalled block. An honest producer may produce such a block
(Construction~\ref{def:honest-production}), so the best chain of $\Pibc$
grows under the conditions under which that of $\Piorigbc$ does.

The argument is prose. It also lists a change to contextual validity that
does not affect coinbase transactions; that change belongs to Crosslink~1,
and Crosslink~2 changes no contextual rule (Definition~\ref{def:bc-validity}),
so the item is vacuous. What the argument shows is that a valid block can
always be produced, so the property it transfers is the growth of the best
chain. Inclusion of user transactions is not preserved: the book notes that
Crosslink~2 intentionally gives up availability during a long finalisation
stall, whose blocks may all be coinbase-only. For Zcash the liveness of
$\Piorigbc$ is itself an assumption: the liveness clause of Assumption
``Persistence and liveness'' (\emph{FlyClient Guide}, \S``The
variable-difficulty backbone model'') is adopted there without proof.

\begin{proposition}[Liveness of $\Pibft$]\label{prop:bft-live}
The protocol $\Pibft$ is live whenever $\Piorigbft$ is.

Status: argued (TFL book).
\end{proposition}

The book argues as follows. Since $\Pibc$ is live, a fresh header chain that
satisfies the Linearity and Tail Confirmation rules
(Definition~\ref{def:bft-validity}) eventually exists for the field
$\headersbc$ of a proposal. When none is available, an honest proposer
repeats the field of the parent bft-block
(Construction~\ref{def:honest-proposal}), which satisfies both rules, so an
honest proposer can always make a valid proposal, as the liveness of a
protocol such as adapted Streamlet (Example~\ref{ex:streamlet}) requires.
The changes to voting only require a vote to be for a proposal that ``could
have been honestly proposed'', and finality in context is unchanged from
$\Piorigbft$.

The argument rests on Proposition~\ref{prop:bc-live}, itself argued. It
does not address the Confirmed best-chain criterion of
Construction~\ref{def:honest-voting}, which each validator checks against
its own best chain: whether enough voting units can be cast for an honest
proposal depends on the validators' best chains, and the construction
chapter relies there on enough validators being synced that Prefix
Consistency holds for them.

\begin{proposition}[Transfer by erasure]\label{prop:erasure}
Every safety property of $\Piorigbft$, Final Agreement included, holds for
$\Pibft$.

Status: argued (TFL book).
\end{proposition}

The book argues that the changes from $\Piorigbft$ to $\Pibft$ only add the
field $\headersbc$ (Definition~\ref{def:fields}) and tighten validity
(Definition~\ref{def:bft-validity}). Erasing the added field therefore maps
every execution of $\Pibft$ to an execution of $\Piorigbft$ with the same
nodes honest at every time. A safety property asserts only that certain
executions do not occur and does not mention the added field, so each safety
property of $\Piorigbft$, Final Agreement included, holds for $\Pibft$;
liveness is the subject of Proposition~\ref{prop:bft-live}. The argument is
sound in outline, but two of its steps are unjustified. Erasure
does not by itself give an execution of $\Piorigbft$: the proposer's
signatures, the ballots, the notarisation proofs and the parent hashes
commit to $\headersbc$ (Definitions~\ref{def:bft-protocol}, \ref{def:units}
and~\ref{def:bft-block}), so the field must instead be read as part of the
content of a $\Piorigbft$ proposal, which the interface leaves open. And the
book asserts without argument that the same nodes are honest: honest voting
in $\Pibft$ is ``only if'' (Remark~\ref{rem:only-if}), so the image may
withhold votes that honest voting in $\Piorigbft$ would cast.

\begin{remark}[Transfer from $\Piorigbc$]\label{rem:origbc-transfer}
That Prefix Consistency and Prefix Agreement transfer from $\Piorigbc$ to
$\Pibc$ is unproved. Of Prefix Agreement the TFL book holds that the changes
to $\Pibc$ ``seem unlikely to have any significant effect on this property,
given that we proved above that they do not affect liveness'', the liveness
argument being prose (Proposition~\ref{prop:bc-live}), and marks the sentence
``TODO: that is a handwave; we should be able to do better, as we do for
$\Pibft$ below'' (section ``Proof of safety for $\LOGfin$''). Erasure does
not settle the question: an honest node of $\Pibc$ rejects blocks valid in
$\Piorigbc$ that fail a rule of Definition~\ref{def:bc-validity}, so its
best chain need not be that of an honest node of $\Piorigbc$ with the same
view, and the rules it adds depend on $\Pibft$. This is the open problem to
which the caveat of Corollary~\ref{cor:bft-subversion} refers. For Zcash the
prefix properties of $\Piorigbc$ are themselves assumptions
(Remark~\ref{rem:pow-prefix}).

Status: open problem.
\end{remark}

\subsection{Progress of finalisation}\label{sec:live-fin}

\begin{remark}[Advance of finalisation]\label{rem:fin-advances}
That finalisation advances, that is, that a later snapshot eventually
enters $\fin$ when both sub-protocols are live and enough participants are
honest, is unproved. The TFL book argues that liveness of $\Pibft$ and the
Linearity rule imply that, ``provided there are sufficient honest
bft-proposers/validators'', a bft-block with a higher-scoring snapshot is
eventually final in the context of the longest valid bft-chain; that by the
Extension rule it descends from the previous last final block in the
context that best-chain producers see; and that the new finalised chain
therefore extends the old. It marks the first step ``TODO make that more
precise''. Neither the number of honest participants nor any time bound is
stated.

Status: open problem.
\end{remark}

\begin{proposition}[Stalled Mode only on failure of the liveness
assumptions]\label{prop:stall-only}
While there are sufficient honest best-chain producers and sufficient honest
bft-proposers and validators, Stalled Mode is not entered, provided that $L$
is large enough relative to $\sigma$, the network delay, the honest hash rate
and the details of $\Pibft$.

Status: argued (TFL book).
\end{proposition}

The book argues informally that, while there are sufficient honest
best-chain producers and sufficient honest bft-proposers and validators, the
finalisation point implied by the field $\contextbft$ at the tip of each
node's best chain advances fast enough for the bound $L$ not to be reached.
It states that this depends on $L$ relative to $\sigma$, on the network
delay, on the hash rate of honest producers, on the number of honest
proposers and the proportion of voting units they hold, and on details of
the BFT protocol, and it marks the argument ``TODO more detailed argument
needed, especially for the dependence on $L$''. No hypothesis is
quantified.

Two results show hypotheses of the statement to be necessary. By
Proposition~\ref{prop:forced-entry}, producers that never advance
$\contextbft$ over $L$ consecutive heights force a stalled block whatever
$\Pibft$ does, so honest best-chain producers are needed as well as a live
$\Pibft$. By the following proposition, Stalled Mode persists while
$\Pibft$ is live if $L \le \sigma$, so $L \ge \sigma + 1$ is needed.

\begin{proposition}[Stall when $L \le \sigma$]\label{prop:stall-short-L}
If $L \le \sigma$, then in every execution in which the bc-blocks of all
views lie on one bc-chain, every participant is honest and no voting unit is
cast non-honestly (Definition~\ref{def:nonhonest}), every bc-valid block of
height greater than $L$ in a view is in Stalled Mode, whether or not $\Pibft$
is live.

Status: proved (this volume), of a designed-but-unspecified construction.
\end{proposition}

\begin{proof}
Let $H$ be such a block. If $\LF(H) = \Obft$, then $\finalitydepth(H) =
\mathsf{height}(H) > L$, since $\snapshot(\Obft) = \Obc$.
Otherwise, by the Tail Confirmation rule, $\LF(H)$ carries the
$\sigma$-block tail of a bc-valid chain $T'$. Its notarisation proof is valid,
so at least one unit was cast for the proposal of $\LF(H)$
(Definition~\ref{def:units}), honestly by hypothesis, and its holder
downloaded the blocks of the carried headers
(Construction~\ref{def:honest-voting}(i)); so $T'$ is in a view and lies with
$H$ on one bc-chain. It is neither $H$ nor a descendant of $H$,
since either would give a cycle of commitments through $H.\contextbft$ and
$\LF(H)$ (Remark~\ref{rem:validity-recursion}); so $T'$ is an ancestor of
$H$, the carried headers are headers of ancestors of $H$, and
Proposition~\ref{prop:lag}(ii) gives $\finalitydepth(H) \ge \sigma + 1 > L$.
\end{proof}

\begin{remark}[Choice of $L$]\label{rem:choose-L}
No source gives a rule for choosing $L$. The TFL book requires $L$
``significantly greater than $\sigma$'' and states that in practice it
``should be \emph{at least} $2\sigma$'' (chapter ``The Crosslink 2
Construction'', sections ``Parameters'' and ``Stalled Mode''), without
derivation, and its liveness argument marks the dependence on $L$ as
missing. The condition $L \ge \sigma + 1$, necessary for
Proposition~\ref{prop:stall-only} by Proposition~\ref{prop:stall-short-L},
is a lower bound only.

Status: open problem.
\end{remark}

\section{Bounded availability and finality overrides}\label{sec:bounded}

This section states the TFL book's informal notion of bounded availability
and the book's case for it, proves the precise statement for Crosslink~2
that the book leaves unproved, shows that the statement needs
$\sigma \le L$, and states the book's argument for finality overrides, for
which the construction has no rule. The objects of the section are designed
but unspecified (Definition~\ref{def:classes}), and each claim carries its
status.

\subsection{Bounded availability}\label{sec:bd-notion}

\begin{remark}[Bounded availability]\label{rem:bounded-avail}
Bounded availability is the TFL book's weakening of dynamic availability,
the property P2 of Definition~\ref{def:ntt-ebbflow}: the finalisation gap,
the number of best-chain blocks above the finalisation point, may grow
during a finalisation stall only up to a bound, beyond which ordinary spends
stop. The book does not define it formally: ``We do not define bounded
availability formally in this document'' (chapter ``The Arguments for
Bounded Availability and Finality Overrides'', section ``What is Bounded
Availability?''). Its only formal property, Bounded hazard-freeness, is
stated for snap-and-chat (chapter ``Notes On Snap-and-Chat''): no node $i$
at time $t$ observes in $\LOGda[i]^t$ a hazardous transaction, such as a
spend, from a block $H$ of $\ch_i^t$ with $\mathsf{height}(H) -
\mathsf{height}(\lca(\snapshot(\LF(H)), H)) > L$, the Crosslink~1 form of the
finality depth. The book gives it no form for Crosslink~2. Read with
$\LOGba[i]^t$ for $\LOGda[i]^t$ (Remark~\ref{rem:ledger-names}), it holds for
Crosslink~2 under the hypotheses of Proposition~\ref{prop:hazard-free}. It
bounds spends relative to the final snapshot of each block; the next
subsection bounds them relative to $\fin_i^t$.
\end{remark}

\begin{proposition}[Bounded hazard-freeness]\label{prop:hazard-free}
Suppose that no stalled block carries a hazardous transaction, and let the best
chain $\ch_i^t$ of node $i$ be bc-valid. Then no block $H \preceqbc \ch_i^t$
with $\mathsf{height}(H) - \mathsf{height}(\lca(\snapshot(\LF(H)), H)) > L$
contributes a hazardous transaction to $\LOGba[i]^t$.

Status: proved (this volume), of a designed-but-unspecified construction.
\end{proposition}

\begin{proof}
For $H = \Obc$ the difference of heights is $0 < L$, since
$\snapshot(\LF(\Obc)) = \Obc$ (Definitions~\ref{def:bc-validity}
and~\ref{def:params}), so $H \neq \Obc$; as a block of the bc-valid chain
$\ch_i^t$, $H$ is bc-block-valid (Definition~\ref{def:bc-validity}). By the
Last Final Snapshot rule $\snapshot(\LF(H)) \preceqbc H$, so the last common
ancestor is $\snapshot(\LF(H))$ (Definition~\ref{def:lca}) and the difference
is $\finalitydepth(H) > L$ (Definition~\ref{def:fdepth}); the Finality depth
rule makes $H$ stalled, and a stalled block carries no hazardous transaction.
\end{proof}

\begin{remark}[Rationale of bounded availability]\label{rem:bound-reasons}
The TFL book argues as follows (chapter ``The Arguments for Bounded
Availability and Finality Overrides''). A partition can stall any protocol
with finality for as long as it lasts, a consequence of the CAP theorem
that the book attributes, made precise, to Lewis-Pye and Roughgarden. Under
strict dynamic availability the best chain meanwhile keeps advancing, so the
finalisation gap grows without bound while spends continue, and the cost of
rolling back spent funds grows with it. Bounded availability caps that
exposure at $L$ blocks and then stops ordinary spends. The best chain keeps
producing stalled blocks rather than halting, because a halt gives an
attacker unbounded time to build a fork of the recent tail, which the book
calls a tail-thrashing attack; the book adds that honest hash rate is
expected to fall once a stall is known, miners expecting their blocks to be
rolled back, which raises the attacker's share.

Status: argued (TFL book).
\end{remark}

% tail_thrash.py: share 1/10: one block per 10 block times; 10 blocks in 100 block times: ok
The book's example of tail thrashing: an attacker with a tenth of the hash
rate finds in expectation one block per $10$ block times, hence, by
linearity of expectation (\emph{Math Guide}, \S``Random variables and
expectation''), $10$ blocks, a $10$-block fork, in $100$ block times.

\subsection{The bounded non-stalled run}\label{sec:bd-run}

\begin{proposition}[Bounded non-stalled run]\label{prop:bounded-run}
Let $H = \ch_i^t$, and let node $i$ have followed the update rule $\updatefin$
(Definition~\ref{def:fin}) up to $t$. Suppose that the Extension, Last Final
Snapshot and Finality depth rules hold along $H$, that the Parent and
Linearity rules hold along the bft-chain that ends at $\LF(H)$, that
$\sigma \le L$, and that $\fin_i^t \preceqbc H$. Then every block of $H$ at
height greater than $\mathsf{height}(\fin_i^t) + L$ is stalled.

Status: proved (this volume), of a designed-but-unspecified construction.
\end{proposition}

\begin{proof}
The Extension rule presupposes that every block of $H$ has a last final
context, so each context cited along $H$ is a valid bft-block
(Definition~\ref{def:blf}). Write $S = \snapshot(\LF(H))$, the final
snapshot of $H$. By the Last Final Snapshot rule $S \preceqbc H$, and
$\fin_i^t \preceqbc H$ by hypothesis, so $S$ and $\fin_i^t$ agree
(Lemma~\ref{lem:linear-prefix}).

First, $\snapshot(\LF(H')) \preceqbc S$ for every $H' \preceqbc H$.
Lemma~\ref{lem:lf-monotone} gives $\LF(H') \preceqbft \LF(H)$; by the Parent
rule every bft-block $D$ with $\LF(H') \precbft D \preceqbft \LF(H)$ is
valid, the Linearity rule gives $\snapshot(D\trunc{1}) \preceqbc
\snapshot(D)$ for each, and transitivity gives the claim.

\emph{Case $S \preceqbc \fin_i^t$.} Let $H' \preceqbc H$ with
$\mathsf{height}(H') > \mathsf{height}(\fin_i^t) + L$. Then
$\snapshot(\LF(H')) \preceqbc S \preceqbc \fin_i^t$, so
\[
  \finalitydepth(H') \ge \mathsf{height}(H') - \mathsf{height}(\fin_i^t) > L,
\]
and the Finality depth rule makes $H'$ stalled.

\emph{Case $\fin_i^t \precbc S$.} Then $H \ne \Obc$, so node $i$'s best
chain has changed to $H$, and by Lemma~\ref{lem:candidate-lfs} the chain
$\candidate(H)$ is the shorter of $S$ and $H\trunc{\sigma}$. At that change,
node $i$ following $\updatefin$, the rule either set $\fin_i^t = \candidate(H)$,
or kept a value with $\fin_i^t \mathrel{\not\preceq}_{\bc} \candidate(H)$. In
the first sub-case $\candidate(H) \precbc S$, so $\candidate(H) =
H\trunc{\sigma}$ and $\fin_i^t = H\trunc{\sigma}$. In the second, $\fin_i^t$ and
$\candidate(H)$ are prefixes of $H$ and agree, so $\candidate(H) \precbc
\fin_i^t \precbc S$; again $\candidate(H) = H\trunc{\sigma}$, now a strict
prefix of $\fin_i^t$. In both sub-cases
\[
  \mathsf{height}(H) - \mathsf{height}(\fin_i^t) \le \sigma \le L,
\]
so no block of $H$ lies at height greater than $\mathsf{height}(\fin_i^t) + L$.
\end{proof}

% ledger_checks.py: 5. sigma <= L: every block more than L above fin is stalled (3000 trials, 270980 checks of 44776 distinct blocks): ok
The proof uses $\sigma \le L$ only in the second case. A randomised check by
script, over $3000$ trials of random block trees, bft-chains and best-chain
switches, finds, in $270980$ checks of $44776$ distinct blocks, every block
more than $L$ above $\fin$ stalled.
The proposition is the precise Crosslink~2 form of the TFL book's remark
that ``the number of non-Stalled Mode bc-blocks after the consensus
finalization point is bounded by $L$, but we have also not proven that so
far'' (chapter ``Security Analysis of Crosslink 2'', section ``Proof of
safety for $\LOGba$''). The bound needs $\sigma \le L$. The book's
requirement that $L$ be ``significantly greater than $\sigma$'', and in
practice at least $2\sigma$, implies it, but the book ties the requirement to
no bound; for some best-chain protocol of the interface the constraint cannot
be dropped.

\begin{proposition}[Necessity of the constraint $\sigma \le L$]
\label{prop:short-L}
For all $L$ and $\sigma$ with $1 \le L < \sigma$ there are a best-chain
protocol $\Piorigbc$ with the interface of Definition~\ref{def:bc-protocol}
and an execution in which the bc-valid best chain $\ch_i^t$ of a node $i$
has a tip that is not stalled and lies more than $L$ blocks above
$\fin_i^t$.

Status: proved (this volume), of a designed-but-unspecified construction.
\end{proposition}

\begin{proof}
Fix $h \ge \sigma$. Let $A$ be a bc-chain with blocks $A_1, \ldots, A_h$,
the index being the height, and let a branch fork from it at $A_{h-1}$ with
blocks $A'_h, \ldots, A'_{h+\sigma-1}$. Let $B$ be a bft-block with parent
$\Obft$ whose field $\headersbc$ is the $\sigma$-block tail of the chain
$A'_{h+\sigma-1}$, so that $\snapshot(B) = A_{h-1}$, and let $C$ be a valid
bft-block with $\bftlastfinal(C) = B$ (in a protocol that decides each
block individually, $C = B$). The blocks $A_1, \ldots, A_{h-1}$ and those of
the branch cite $\Obft$, each block whose finality depth exceeds $L$ being
stalled (Remark~\ref{rem:producible}); so the chain $A'_{h+\sigma-1}$ is
bc-valid, and $B$ satisfies the Tail Confirmation rule and, since
$\snapshot(\Obft) = \Obc$, the Linearity rule.

Let the tip $A_h$ cite $C$ and not be stalled, as under the proposed Zcash
policy (Remark~\ref{rem:stalled-policy}) any block with a non-coinbase
transaction. The tip satisfies the Valid context rule; the Extension rule,
since $\LF(A_{h-1}) = \Obft \preceqbft B = \LF(A_h)$; the Last Final
Snapshot rule, since $A_{h-1} \preceqbc A_h$; and the Finality depth rule,
since $\finalitydepth(A_h) = 1 \le L$. Let node $i$'s best chain grow from
$\Obc$ through $A_1, \ldots, A_h$, node $i$ receiving the branch with $A_h$,
and let $\score(A_h) > \score(A'_{h+\sigma-1})$, which
Definition~\ref{def:bc-protocol} permits, since it requires the score to
increase only along a chain. Every candidate before $A_h$ is $\Obc$.
By Lemma~\ref{lem:candidate-lfs}, $\candidate(A_h)$ is the shorter of
$A_{h-1}$ and $A_h\trunc{\sigma} = A_{h-\sigma}$, which is $A_{h-\sigma}$
since $\sigma > L \ge 1$; $\updatefin$ sets $\fin_i^t = A_{h-\sigma}$, and the
tip lies $\sigma > L$ blocks above it.
\end{proof}

% ledger_checks.py: 6. for every 1 <= L < sigma <= 20 a non-stalled tip lies sigma above fin; its context carries a sibling branch's headers (on its own chain the finality depth exceeds L): ok
A script checks the construction for every $1 \le L < \sigma \le 20$. The
context of the tip carries the headers of a sibling branch, as it must: by
Proposition~\ref{prop:lag}, a context whose last final block carries
headers of ancestors of the tip gives the tip finality depth at least
$\sigma + 1 > L$. Every hypothesis of Proposition~\ref{prop:bounded-run}
other than $\sigma \le L$ holds in the construction. Under accumulated work
on Mainnet (Remark~\ref{rem:bc-zcash}) the instance is no best chain of a
node that holds the branch: $A_h$ and $A'_h$ have the same ancestors, hence
the same target and work, and the branch adds $\sigma - 1 \ge 1$ blocks; on
Testnet a minimum-difficulty sibling can carry less work, and the argument
fails. The
proposition is shown for the interface of Definition~\ref{def:bc-protocol},
not for the Zcash instance; whether $\sigma \le L$ is necessary for Zcash is
open.

\begin{remark}[Finalised chain off the best chain]\label{rem:fin-off-chain}
Proposition~\ref{prop:bounded-run} does not cover a node whose finalised
chain is not a prefix of its best chain, $\fin_i^t \mathrel{\not\preceq}_{\bc}
\ch_i^t$. Under unchanged fork choice, for a node honest at $t$, that case
arises only after a failure of Prefix Consistency at $\sigma$:
Lemma~\ref{lem:fin-depth} gives a time $r \le t$ with $\fin_i^t \preceqbc
\ch_i^r\trunc{\sigma}$, and Prefix Consistency at $\sigma$, with $j = i$, gives
$\ch_i^r\trunc{\sigma} \preceqbc \ch_i^t$. A node that follows $\updatefin$ but
is not honest, such as an eclipsed one (Definition~\ref{def:exec}), can reach
the case while Prefix Consistency holds. The TFL book does not define bounded
availability formally, and every theorem about it beyond
Propositions~\ref{prop:bounded-run}, \ref{prop:short-L}
and~\ref{prop:hazard-free}, this case included, is open.

Status: open problem.
\end{remark}

\subsection{Finality overrides}\label{sec:bd-override}

\begin{remark}[Override argument]\label{rem:override}
The TFL book argues that consensus protocols fail, through a design flaw of
the finality layer that causes or permits a stall, a balance-violation or
spend-authorisation flaw that is exploited or likely to be, or an
implementation bug that makes many nodes diverge from consensus, and that
overriding finality may then be the best response (chapter ``The Arguments
for Bounded Availability and Finality Overrides'', sections ``Finality
Overrides'' and ``Complementarity''). The override procedure should be
formalised, its governance process included, because that makes security
analysis of the protocol, of the governance and of their interaction
feasible, and settles in advance the procedure stakeholders agree to
follow. Modelling an overridden protocol as a new protocol and chain is
rejected as insufficient: it does not match how users regard an override
and does not help its analysis. Overrides complement Stalled Mode.
Validators asked to stop validating cause a finalisation stall; once the
stall exceeds the availability bound, the best chain enters Stalled Mode,
so that at most $L$ blocks with spends lie above a node's finalised chain
$\fin_i^t$ (Proposition~\ref{prop:bounded-run}, under $\sigma \le L$ and
$\fin_i^t \preceqbc \ch_i^t$, when stalled blocks carry no spends, as under the
coinbase-only policy of Remark~\ref{rem:stalled-policy}) and a long rollback
stays realistic. A coalition of validators that ignores the request removes only
the extra time for a governance decision. An intentional rollback should
preserve the coinbase-only mining rewards, so that producers keep an incentive
to produce blocks during the stall. The argument depends on the freedom of
validators to refuse a vote (Remark~\ref{rem:only-if}).

Status: argued (TFL book).
\end{remark}

\begin{remark}[Overrides outside the construction]\label{rem:override-outside}
Crosslink~2 has no override rule, message or state: an override is a change
of the consensus rules, a hard fork outside the construction. Assured
Finality (Definition~\ref{def:af}) binds executions of the protocol only. A
reversal out of band, by users who fork the network with a new validity
rule, remains possible, and for that reason the TFL book avoids the term
``absolute finality'' (chapter ``Terminology''). The deployed chain has the
same boundary: consensus imposes no maximum reorganisation depth, and the
requirement that a validator's chain include the activation block of the
most recent settled network upgrade is a social checkpoint, distinct from
probabilistic finality (\emph{Consensus Guide}, \S``Rollback limits,
coinbase maturity, and finality'').
\end{remark}

\begin{remark}[Open parts of the override design]\label{rem:override-open}
No source designs a governance process for finality overrides; the TFL book
argues only that one should be formalised. Its out-of-date chapter
``Potential Changes to Crosslink'' recommends that honest best-chain
producers record every proposed and notarised bft-block, not only the cited
context, which would expose double proposals and double notarisations. Its
chapter ``Questions About Crosslink'' proposes posting objective evidence of
a violation of the safety of $\Pibft$ or of its assumptions, two final
bft-blocks with the same parent or more than one third of a BFT epoch's
voting units cast for conflicting blocks, in a best-chain transaction,
optionally latching the best chain into Stalled Mode until manual
intervention. Neither is part of Crosslink~2.

Status: open problem.
\end{remark}

\section{The Shielded Labs design}\label{sec:sl}

This section states how Shielded Labs' \emph{Zcash Crosslink Design
Overview} (draft~10, 14 September 2026, last edited 5 October 2026), the
\emph{overview} below, departs from the construction of \S\ref{sec:cons}
and \S\ref{sec:ledgers}: its terms and the location of ledger state, the
rules it keeps and drops, its certificates and cross-link references, its
fork choice, its proposer behaviour, its finality vocabulary, and which
results of \S\ref{sec:safety} and \S\ref{sec:bounded} survive, each premise
read through the design's own validity notion. The overview is cited by
section title. Every component this section states is designed but
unspecified (Definition~\ref{def:classes}) unless it is marked an open
problem; no ZIP specifies any of them.

\subsection{Terms and ledger state}\label{sec:sl-terms}

The overview departs from the book chiefly in four respects; its further
departures are stated in \S\ref{sec:sl-rules}, \S\ref{sec:sl-certs} and
Remark~\ref{rem:sl-proposer}. It has no Stalled Mode, no finalisation gap bound
$L$ and no bounded-available chain. It replaces the unchanged fork choice by a
sticky fork choice, which never follows a chain that does not contain $\fin$. It
adds an Inclusion Depth rule, a lower bound on how far a bc-block lies above the
snapshot of the bft-block it cites. It punishes the stake behind finalisers by
user-configured slash forks instead of by any in-protocol slashing. It is a
living document, and the items it marks as to be confirmed are unsettled.

The following terms of the overview are used in this section and the next
two. \emph{Tenderlink} is the design's BFT protocol $\Piorigbft$, derived
from Tendermint (E.~Buchman, J.~Kwon and Z.~Milosevic, ``The latest gossip
on BFT consensus'', arXiv 1807.04938, 2018); it has no published
specification. A \emph{trailing finality certificate}, or \emph{certificate}, is
the design's bft-block: it carries the field $\headersbc$ of
Definition~\ref{def:fields} and has the snapshot of
Definition~\ref{def:snapshot}, but it contains neither its notarisation
signatures nor an outer signature. A \emph{fat pointer} is the field of a
bc-block that references a certificate by hash together with notarisation
signatures; its rule is Definition~\ref{def:fat-pointer}. A \emph{lockbox}
is a finaliser's reward account. The staking objects are glossed here and
defined in \S\ref{sec:fs}. A staker holds delegation bonds, each an amount
of ZEC, its \emph{stake}, assigned to one finaliser and reassignable to
another by redelegation (Definition~\ref{def:bond}). The voting power of a
finaliser is the stake delegated to it (Definition~\ref{def:voting-power}). The
\emph{BFT height} of a certificate is its height in the bft-chain, and the
\emph{roster} of a BFT height is the set of finalisers with voting powers
derived from staking state on the bc-chain, of which the active roster's votes
count at that height (Definition~\ref{def:roster}). A \emph{slash fork} is
a user-configured hard fork that burns the bonds delegated to named
finalisers (Definition~\ref{def:slash-fork}).

All ledger state of the design, the roster included, lives on the
bc-chain, and the bft-chain holds none (overview, \S``3. The proof-of-work
side: the fat pointer and where it lives''). Every payout, burn and slash is
computed from what the bc-chain records about the bft-chain, and the fat
pointer is the only such record.

\subsection{Retained and removed rules}\label{sec:sl-rules}

The design keeps, on the certificate side, the Tail Confirmation rule and
the Linearity rule of Definition~\ref{def:bft-validity}; on the best-chain
side, the Valid context, Extension and Last Final Snapshot rules of
Definition~\ref{def:bc-validity}; and $\updatefin$ of Definition~\ref{def:fin}
(overview, \S``4. The BFT side: TFCs and $\sigma$'', \S``9. Sticky fork
choice''). The Parent rule has no counterpart: the parent
of a certificate carries no signatures, so judging a certificate does not
judge its parent, which is judged only by a bc-block that cites it, if one
exists. The Tail
Confirmation rule is read as in the book: a validator downloads and validates
the bc-blocks under the carried headers, since the rule requires the tail of a
bc-valid chain and headers alone do not establish validity.

The design drops the Finality depth rule, Stalled Mode with its predicate
$\isstalledblock$, the parameter $L$, the bounded-available chain $\ba_\mu$
with its depth $\mu$, and the outer proposer signature of
Definition~\ref{def:bft-block}. The overview's bc-block rules are the Valid
context, Extension, Last Final Snapshot and Inclusion Depth rules
(\S``4. The BFT side: TFCs and $\sigma$''), it states that the design has no
Stalled Mode (\S``9. Sticky fork choice''), and neither its certificate nor its
fat pointer carries an outer signature (\S``3. The proof-of-work side: the fat
pointer and where it lives'').

\begin{remark}[Rationale of the Shielded Labs design]\label{rem:sl-rationale}
A blog post of J.~McGee for Shielded Labs of 29 May 2026 allocates roles:
wallets protect users; miners include transactions and produce blocks;
finalisers enforce finality; stakers ``select good finalizers''; users
``slash stakers if they fail to do their job''. The overview supplies the
mechanisms, redelegation and the slash fork, and draws the consequence for
stalls. Redelegation and slash forks are to resolve them,
``finality is a ratchet rather than an override'' (\S``2. The one-paragraph
version''), the finality layer's only job is to finalise the proof-of-work
chain, and if that layer fails the network ``decays gracefully back to plain
PoW'' (\S``1. Purpose''); a stall ``can never be proven permanent'', so the
judgement is left to users, and the threat of a slash fork is to move stake
before any burn (\S``16. Social slashing''). The book's Stalled Mode instead
bounds the run of non-stalled blocks above $\fin$
(Proposition~\ref{prop:bounded-run}; Remark~\ref{rem:bound-reasons}).

Status: argued (Shielded Labs).
\end{remark}

\subsection{Certificates and cross-link references}\label{sec:sl-certs}

In the design $\Piorigbft$ is Tenderlink (overview, \S``6. The BFT
implementation''), where the book leaves it abstract and gives adapted
Streamlet as an example (Example~\ref{ex:streamlet}). Tenderlink's peers
agree on each decision, whether a certificate is approved by two thirds of
the stake-weighted voting power, but not on the set of signatures that
made up that two thirds.

\begin{definition}[Fat pointer]\label{def:fat-pointer}
In the Shielded Labs design each bc-block $H$ other than $\Obc$ carries a
\emph{fat pointer}, which is either nil or the hash of a certificate $C$
other than $\Obft$, which is the commitment $H.\contextbft$ of
Definition~\ref{def:fields}, together with a set of notarisation signatures
attesting $C$, which lie outside $C$; every bc-block of height at most $h_2$
(Definition~\ref{def:bootstrap}) carries a nil fat pointer. The signatures
are Ed25519 signatures, a digital signature scheme in the sense of the
\emph{Crypto Guide}, \S``Syntax and security goal''. The block $H$
\emph{cites} $C$. A block with a nil fat pointer has context $\Obft$ and
cites no certificate, and~(i) and the Inclusion Depth rule do not apply to
it. Under the design:
\begin{enumerate}[label=(\roman*)]
\item the Valid context rule of Definition~\ref{def:bc-validity} reads: the
  certificate that $H$ cites satisfies the Linearity and Tail Confirmation
  rules of Definition~\ref{def:bft-validity} and Tenderlink's own rules for
  a proposal, without the Parent rule, and the signatures in the fat pointer
  of $H$ are valid signatures on it by holders of at least two thirds of the
  stake of the active roster (Definition~\ref{def:active-roster}) of the
  roster at its BFT height;
\item a certificate $C$ is \emph{validly cited in honest view} iff some
  party, at a time at which it is honest, observes a bc-block that cites
  $C$ and whose fat pointer meets~(i); this reading replaces ``in honest
  view'' of Definition~\ref{def:bft-block} for certificates.
\end{enumerate}
A copy of a certificate is not evidence of canonicity, since the
certificate carries no signature attesting it.
\end{definition}

The book's bft-block $(P, \proofsym_P)$ under an outer signature is valid
or invalid on its own (Remark~\ref{rem:objective}). In the design ``a
certificate cannot carry its own signatures'' (overview, \S``3. The
proof-of-work side: the fat pointer and where it lives''), so validity of a
cited certificate is decided from the citing bc-block: two blocks may cite
one certificate with different signature sets, and a certificate that no
block cites is judged by no block. Where a block stores its fat pointer
is an open problem: the overview weighs options from a commitment behind the
block commitment hash (\emph{FlyClient Guide}, \S``The block commitment hash'',
Remark ``Extension of the block commitment hash''), which it judges probably
insufficient, to a change of the block header, and lists the choice among its
open questions.

\begin{remark}[Finality of each decision]\label{rem:sl-final-each}
Tenderlink decides each certificate individually, so $\bftlastfinal$ is the
identity on decided certificates and $\LF(H)$ is the certificate that $H$
cites: the overview takes $\LF$ of a tip to be the certificate at which
the tip's $\contextbft$ points (\S``9. Sticky fork choice''). In the book's
interface $\bftlastfinal(C)$ may instead be a strict ancestor of $C$, as
in adapted Streamlet (Example~\ref{ex:streamlet}).
\end{remark}

\begin{remark}[Deferred validity]\label{rem:sl-deferred}
The design adds a third state, ``not yet determinable'', on both sides: a
bc-block, proposal or certificate whose referenced data a node does not yet
hold is deferred, not rejected, and the deferral must be able to expire or
be abandoned, since a malicious peer can reference a hash for which no data
exists (overview, \S``5. Mutual reference and its consequences''). The state
concerns a node's knowledge of referenced data; validity, as a function of a
block, proposal or certificate and of the blocks it commits to by hash
(Remarks~\ref{rem:validity-recursion} and~\ref{rem:objective}), stays
two-valued. The expiry rule is marked to be confirmed.

Status: open problem.
\end{remark}

% ledger_checks.py: 3. candidate below the truncated tip (47850 cases); with the carried headers on the tip's chain, lag = finality depth >= sigma + 1 and >= sigma + 2 confirmations (7330 cases); lag sigma + 1 blocks, sigma + 2 confirmations, attained, for every sigma <= 20; genesis context excluded (lag height(H) <= sigma for a low tip): ok
\begin{lemma}[Inclusion Depth as a finality-depth bound]
\label{lem:inclusion-depth}
In the Shielded Labs design let a bc-block $H$ cite a certificate $C$ (so
$C \ne \Obft$), so that $\LF(H) = C$ (Remark~\ref{rem:sl-final-each}). The
Inclusion Depth rule, that $H$ may cite $C$ only if
\[
  \mathsf{height}(H) \ge \mathsf{height}(\snapshot(C)) + \sigma + 1,
\]
holds for $H$ if and only if $\finalitydepth(H) \ge \sigma + 1$.

Status: proved (this volume), of a designed-but-unspecified construction.
\end{lemma}

\begin{proof}
By Definition~\ref{def:fdepth} and $\LF(H) = C$,
\[
  \finalitydepth(H) = \mathsf{height}(H) - \mathsf{height}(\snapshot(\LF(H)))
  = \mathsf{height}(H) - \mathsf{height}(\snapshot(C)),
\]
so the two inequalities coincide.
\end{proof}

Under the Last Final Snapshot rule the finality depth of $H$ is the number
of blocks of $H$ above its final snapshot. In the book the bound
$\finalitydepth(H) \ge \sigma + 1$ holds without a rule whenever $\LF(H)$
is not $\Obft$ and the headers that $\LF(H)$ carries are ancestors of $H$
(Proposition~\ref{prop:lag}). The same holds for a cited certificate $C$
other than $\Obft$: if the $\sigma$ headers of $C$ are ancestors of $H$, the
newest is a strict ancestor, because the header of $H$ commits to $C$, $C$
commits to its headers, and a cycle of hash commitments is excluded
(Remark~\ref{rem:validity-recursion}); the deepest carried header then lies
at height at most $\mathsf{height}(H) - \sigma$, and $\snapshot(C)$, its
parent, at height at most $\mathsf{height}(H) - \sigma - 1$. Hence, among
blocks that cite a certificate other than $\Obft$, the rule excludes
exactly those whose certificate carries a header that is not an ancestor of $H$
and whose final snapshot lies at most $\sigma$ blocks below $H$: a decision
whose tail lies on another branch. Neither Theorem~\ref{thm:pa-af} nor
Theorem~\ref{thm:fa-af} uses the rule, nor any lower bound on finality
depth.

\subsection{Sticky fork choice}\label{sec:sl-forkchoice}

The accumulated work of a bc-chain is the sum of the work of its blocks
(\emph{Consensus Guide}, \S``Work, and the best chain'', Definition ``Work
of a block''), the score of Zcash (Remark~\ref{rem:bc-zcash}).

\begin{definition}[Sticky fork choice]\label{def:sticky}
A node $i$ follows \emph{sticky fork choice} iff, whenever some bc-valid
chain $T$ of its view with $\fin_i^s \preceqbc T$ has more accumulated work
than $\ch_i^s$, or the same accumulated work and a tip hash that wins the
tie-break, it replaces $\ch_i^s$ by the heaviest such chain, ties broken by
tip hash, and it changes its best chain at no other time. The overview
states the rule as ``best = heaviest chain that contains $\fin$'': a heavier
chain that forks below $\fin_i$ is refused, whatever its work (overview, \S``9.
Sticky fork choice''). The overview does not state the direction of the
tie-break.
\end{definition}

The rule departs from Definition~\ref{def:fork-choice}, which chooses among
bc-valid chains by score alone, and its tie-break departs from the
first-received rule of the protocol specification (\emph{Consensus Guide},
\S``Work, and the best chain'', Definition ``Best valid block chain'').

\begin{proposition}[Absence of finalisation hazards under sticky fork
choice]\label{prop:sticky-no-hazard} Let node $i$ follow sticky fork choice and
$\updatefin$. Then $\fin_i^t \preceqbc \ch_i^t$ at every time $t$, and at every
change of the best chain of $i$ from $\ch_i^s$ to $\ch_i^t$ the chains
$\candidate(\ch_i^t)$ and $\fin_i^s$ agree, so the finalisation safety
hazard of $\updatefin$ is never recorded.

Status: proved (this volume), of a designed-but-unspecified construction.
\end{proposition}

\begin{proof}
By induction over the sequence of best-chain changes of $i$, on which
$\updatefin$ is defined (Definition~\ref{def:fin}). Initially
$\fin_i^0 = \Obc = \ch_i^0$. Between changes neither chain changes. At a
change from $\ch_i^s$ to $\ch_i^t$, with $\fin_i^s \preceqbc \ch_i^s$ by
the induction hypothesis, Definition~\ref{def:sticky} gives
$\fin_i^s \preceqbc \ch_i^t$. By Definitions~\ref{def:candidate}
and~\ref{def:lca}, $\candidate(\ch_i^t) \preceqbc \ch_i^t\trunc{\sigma}$,
and by Definition~\ref{def:trunc} $\ch_i^t\trunc{\sigma} \preceqbc \ch_i^t$.
Both $\fin_i^s$ and $\candidate(\ch_i^t)$ are thus prefixes of $\ch_i^t$, and
they agree by Lemma~\ref{lem:linear-prefix}; the case of $\updatefin$ in which
the candidate conflicts with $\fin_i^s$ does not arise, and no hazard is
recorded. The new value $\fin_i^t$ is one of the two, hence a prefix of
$\ch_i^t$.
\end{proof}

The overview asserts the first claim without proof (\S``10. Kinds of finality
and which ledger state is current'': local finality ``always lies on the best
chain''); the proof, and the absence of finalisation hazards, are this
volume's.

\begin{remark}[Consequences of sticky fork choice]
\label{rem:sticky-consequences}
The overview states the following consequences (\S``9. Sticky fork
choice''). The rule departs from most work only when a heavier chain
excludes $\fin$; while the BFT protocol is stalled, $\fin$ is frozen and
selection above it is by work alone. A refused switch is ``the only
observable sign of a finality conflict''. A partition in which only one side
can finalise past the fork point leaves that side ``permanently unwilling to
switch to the other''. Without Stalled Mode the gap between $\fin$ and the
tip is unbounded during a stall, and ordinary spends keep landing above
$\fin$. A miner who dominates the best chain can withhold finality progress
by never updating $\contextbft$, at no cost in validity.

Status: argued (Shielded Labs).
\end{remark}

\begin{remark}[Status of sticky fork choice]\label{rem:sticky-open}
Safety and liveness under sticky fork choice are unproved. The overview
states: ``The \emph{TFL Book} has no safety or liveness proof for this rule.
Its liveness argument relies on unmodified fork choice'' (\S``9. Sticky fork
choice''). No source proves Prefix Consistency, Prefix Agreement
(Definitions~\ref{def:pc} and~\ref{def:pa}) or liveness of the best chain
that the rule selects, and the book's arguments for $\Pibc$
(Proposition~\ref{prop:bc-live}, Remark~\ref{rem:origbc-transfer}) assume
Definition~\ref{def:fork-choice}.

Status: open problem.
\end{remark}

\subsection{Proposer behaviour}\label{sec:sl-proposer}

\begin{remark}[Proposer deviations from the book]\label{rem:sl-proposer}
The design's proposer departs from Construction~\ref{def:honest-proposal} in
two ways, both behaviour and not validity, since a validator cannot tell
whether carried headers are the proposer's best-chain tail (overview, \S``4. The
BFT side: TFCs and $\sigma$''). First, it bounds how far the snapshot of a
proposal may lie above the snapshot of the last decided certificate; when the
bound binds, the proposal carries a tail that ends below the proposer's tip,
which still meets the Tail Confirmation rule. The overview calls the bound ``a
heuristic, not part of Crosslink 2'', and its value is not given here. Second,
where the book's honest proposer repeats the headers of its parent bft-block,
the design's proposer makes no proposal; how often it should repeat instead is
not decided. The second deviation removes the proposal on which the book's
argument for Proposition~\ref{prop:bft-live} relies.
\end{remark}

The overview is internally inconsistent on $\sigma$: its section ``4. The BFT
side: TFCs and $\sigma$'' and its Part~VI, ``Open questions'', state
different values, the first described as unchecked. No value of $\sigma$ is
given here.

\subsection{Finality vocabulary}\label{sec:sl-finality}

\begin{definition}[Shielded Labs finality quantities]\label{def:sl-quantities}
The overview distinguishes three quantities (\S``10. Kinds of finality and
which ledger state is current''), of which the volume keeps two, (i)
and~(ii); the third, a ``database finalized tip'', which the overview calls
``Not finality'', concerns a node's storage, not the protocol, and is
omitted. Item~(iii) is the quantity from which local finality advances
(\S``9. Sticky fork choice''):
\begin{enumerate}[label=(\roman*)]
\item the \emph{BFT final snapshot} of node $i$ at time $t$, the snapshot
  of the newest certificate decided in the view of $i$: a function of the
  decided bft-chain, which need not be a prefix of $\ch_i^t$, for any
  length of time;
\item \emph{local finality}, the chain $\fin_i^t$ of
  Definition~\ref{def:fin}, of which the overview states ``The node will not
  reorg past this block'';
\item for each bc-block $H$, the chain $\candidate(H)$ of
  Definition~\ref{def:candidate}, a function of $H$ and of the chains it
  commits to.
\end{enumerate}
\end{definition}

Local finality is node-local, since it depends on the history of best chains of
$i$; it is monotone (Lemma~\ref{lem:fin-linear}); under sticky fork choice it
lies on the best chain (Proposition~\ref{prop:sticky-no-hazard}). It lies at or
below the BFT final snapshot whenever the decided certificates form one
bft-chain along which the Linearity rule holds: under Definition~\ref{def:fin}
alone, the chain $\fin_i^t$ is $\Obc$ or a prefix of $\snapshot(C)$ for a
decided certificate $C$ (Definitions~\ref{def:fin} and~\ref{def:candidate}), the
newest decided certificate $C^*$ extends $C$, and the Linearity rule along
the path from $C$ to $C^*$ gives $\snapshot(C) \preceqbc \snapshot(C^*)$.
The bootstrap's finalisation of the block at height $h_1$
(Definition~\ref{def:bootstrap}) lies outside $\updatefin$: from height
$h_2$ until a decided certificate's snapshot extends that block, local
finality need not lie at or below the BFT final snapshot, which before the
first decision is $\Obc$, strictly below $\fin_i$, if the genesis bft-block
points at the best-chain genesis; where the genesis bft-block points is open.
The overview requires that ``Consensus must never read $\fin$'': every quantity
that all nodes must compute identically, among them the roster and staking
rewards, is derived from chain data. The overview states that honest nodes'
values of $\fin$ are ``only prefix-compatible, not equal''; prefix-compatibility
across honest nodes is Assured Finality (Definition~\ref{def:af}), which under
the design rests on the unproved hypotheses of Proposition~\ref{prop:sl-fa}.

\subsection{Theorems under the Shielded Labs design}\label{sec:sl-theorems}

The \emph{Shielded Labs rules} below are the kept rules of
\S\ref{sec:sl-rules}, with the Valid context rule and honest view read
through Definition~\ref{def:fat-pointer}, together with the Inclusion Depth
rule and sticky fork choice. An honest node's best chain is bc-valid under
these rules, since sticky fork choice chooses among bc-valid chains. The
rules leave out the bootstrap of Definition~\ref{def:bootstrap}: under them
$\fin$ starts at $\Obc$ and changes only by $\updatefin$. With the bootstrap,
Proposition~\ref{prop:sl-fa} needs the further hypothesis that all honest
nodes finalise the same block at height $h_1$, which
Proposition~\ref{prop:bootstrap-safety} derives, for the honest nodes whose
best chains reach height $h_2$, from Prefix Agreement of $\Pibc$ at depth
$h_2 - h_1$; Proposition~\ref{prop:sl-pa} needs $h_2 - h_1 \ge \sigma$; and in
Proposition~\ref{prop:sticky-no-hazard} the bootstrap value, a prefix of the
best chain, is a third possible value of $\fin$.

\begin{proposition}[Final Agreement route under the Shielded Labs rules]
\label{prop:sl-fa}
Under the Shielded Labs rules, suppose that in an execution
\begin{enumerate}[label=(\alph*)]
\item any two certificates validly cited in honest view
  (Definition~\ref{def:fat-pointer}) lie on one bft-chain, that is, agree;
\item for any two such certificates $C \preceqbft C'$, every certificate $D$
  with $C \precbft D \preceqbft C'$ satisfies the Linearity rule relative to
  its parent: $\snapshot(D\trunc[\bft]{1}) \preceqbc \snapshot(D)$.
\end{enumerate}
Then Assured Finality (Definition~\ref{def:af}) holds in the execution.

Status: proved (this volume), of a designed-but-unspecified construction.
\end{proposition}

\begin{proof}
Let $i$ be honest at $t$ and $j$ honest at $u$. If $\fin_i^t = \Obc$, it is
a prefix of every bc-chain (Definition~\ref{def:chain}) and agrees with
$\fin_j^u$; likewise if $\fin_j^u = \Obc$. Otherwise, by $\updatefin$
(Definition~\ref{def:fin}) there is a best-chain change of $i$ at a time
$r \le t$ with $\fin_i^t = \candidate(\ch_i^r)$, namely the last change up
to $t$ at which $\updatefin$ took its first case. With $H = \ch_i^r$,
Definitions~\ref{def:candidate} and~\ref{def:lca} give $\fin_i^t \preceqbc
\snapshot(\LF(H))$. Since $\fin_i^t \ne \Obc$, the snapshot is not $\Obc$, so
$\LF(H)$ is not $\Obft$ and $H$ is not $\Obc$.
Node $i$ is honest at $r$, so $H$ is bc-valid and its fat pointer meets
the Valid context rule as read in Definition~\ref{def:fat-pointer}(i); as
$i$ observed $H$ at $r$ while honest, the certificate $C_i$ that $H$ cites
is validly cited in honest view. By Remark~\ref{rem:sl-final-each}, $\LF(H) =
C_i$, so $\fin_i^t \preceqbc \snapshot(C_i)$. The same argument for $j$ gives a
certificate $C_j$, validly cited in honest view, with $\fin_j^u \preceqbc
\snapshot(C_j)$.

By (a), $C_i$ and $C_j$ agree; by symmetry let $C_i \preceqbft C_j$, and let
$C_i = D_0, D_1, \ldots, D_k = C_j$ be the certificates of $C_j$ from $C_i$
on, each the parent of the next. By (b), $\snapshot(D_{l-1}) \preceqbc
\snapshot(D_l)$ for $1 \le l \le k$, and by transitivity of $\preceqbc$,
$\snapshot(C_i) \preceqbc \snapshot(C_j)$.
Hence $\fin_i^t$ and $\fin_j^u$ are both prefixes of $\snapshot(C_j)$, and
they agree by Lemma~\ref{lem:linear-prefix}.
\end{proof}

The proof is that of Theorem~\ref{thm:fa-af} with three substitutions:
the Valid context rule and honest view are read through the fat pointer;
Final Agreement (Definition~\ref{def:fa}) becomes hypothesis~(a), because
$\bftlastfinal$ is the identity on decided certificates; and hypothesis~(b)
takes the place of the Parent rule's propagation of validity to every
ancestor of a valid bft-block, since a certificate that no block cites is
judged by no block (Definition~\ref{def:fat-pointer}). It uses neither the
fork choice, nor Stalled Mode, nor an outer signature, nor the Inclusion
Depth rule. Hypotheses~(a) and~(b) concern Tenderlink with its
stake-weighted roster and are unproved (Remark~\ref{rem:tenderlink-fa}).

\begin{proposition}[Prefix Agreement route under the Shielded Labs rules]
\label{prop:sl-pa}
Under the Shielded Labs rules, in every execution in which the best chain
selected by sticky fork choice has Prefix Agreement at depth $\sigma$
(Definition~\ref{def:pa}, with $\ch_i^t$ that chain), Assured Finality
holds.

Status: proved (this volume), of a designed-but-unspecified construction.
\end{proposition}

\begin{proof}
Lemma~\ref{lem:fin-depth} and its proof use $\updatefin$ and
Definition~\ref{def:candidate} only, and no property of the fork choice, so
they hold with $\ch_i^t$ the chain that sticky fork choice selects. Let $i$
be honest at $t$ and $j$ honest at $u$. The lemma gives $r \le t$ and
$r' \le u$ with $\fin_i^t \preceqbc \ch_i^r\trunc{\sigma}$ and
$\fin_j^u \preceqbc \ch_j^{r'}\trunc{\sigma}$. Nodes $i$ and $j$ are honest
at $r$ and $r'$, so by Prefix Agreement at $\sigma$ the chains
$\ch_i^r\trunc{\sigma}$ and $\ch_j^{r'}\trunc{\sigma}$ agree; let the first
be a prefix of the second, the other case being symmetric. Then
$\fin_i^t$ and $\fin_j^u$ are both prefixes of $\ch_j^{r'}\trunc{\sigma}$
and agree by Lemma~\ref{lem:linear-prefix}.
\end{proof}

No source proves Prefix Agreement of the sticky-choice chain
(Remark~\ref{rem:sticky-open}); Manian's Quint suite witnesses its failure on a
hand-built instance with more than one third of the BFT protocol's voting power
subverted (Remark~\ref{rem:mc-zq-findings}). In the
overview's partition example (Remark~\ref{rem:sticky-consequences}), once
both branches extend more than $\sigma$ blocks above the fork point, Prefix
Agreement at $\sigma$ fails in the execution, as under any fork choice. What
the sticky rule adds is that the conflict persists after the partition
heals: the nodes of the side that finalised past the fork point never switch
to the other branch, and the nodes of the other side, whose $\fin$ is at or
below the fork point, keep their branch while it is the heavier. Assured
Finality can still hold in that example, through Proposition~\ref{prop:sl-fa},
since the BFT protocol decided on one side only.

\begin{proposition}[Unbounded run under the Shielded Labs rules]
\label{prop:sl-unbounded}
Under the Shielded Labs rules, which have no Finality depth rule, the
number of blocks above $\fin_i^t$ on the best chain $\ch_i^t$ of an honest
node $i$, none of which is stalled since the design has no Stalled Mode, is
unbounded over executions: for every $d \in \N$ there is an execution in
which $\mathsf{height}(\ch_i^t) - \mathsf{height}(\fin_i^t) \ge d$.

Status: proved (this volume), of a designed-but-unspecified construction.
\end{proposition}

\begin{proof}
Let $H_0 \ne \Obc$ be a bc-block valid under the Shielded Labs rules, so that
it carries a fat pointer (Definition~\ref{def:fat-pointer}) and meets the
Inclusion Depth rule if it cites a certificate, and let
$S = \snapshot(\LF(H_0))$. Consider an execution in which $H_0$ is the best
chain of an honest node $i$ at some time, and then blocks $H_1, \ldots, H_d$ are
produced, each $H_l$ a child of $H_{l-1}$ that meets the Inherited origbc
rules and repeats the fat pointer of $H_{l-1}$, as when no new certificate is
decided or when producers never advance $\contextbft$; node $i$ adopts each
$H_l$ in turn and learns of no other chain. Each $H_l$ is valid: the Valid
context rule holds as for $H_{l-1}$, since the certificate, the signatures and
the roster at its BFT height are the same; by Remark~\ref{rem:sl-final-each},
$\LF(H_l) = \LF(H_{l-1})$, so the Extension rule holds with equality; the Last
Final Snapshot rule holds because $S \preceqbc H_0 \preceqbc H_l$; and the
Inclusion Depth rule holds because the height grows and the snapshot does not
change. Each $H_l$ extends the previous best chain, which contains $\fin_i$
(Proposition~\ref{prop:sticky-no-hazard}), and has more accumulated work, so
sticky fork choice adopts it.

When $H_0$ is adopted, $\fin_i \preceqbc H_0$
(Proposition~\ref{prop:sticky-no-hazard}), so $\mathsf{height}(\fin_i) \le
\mathsf{height}(H_0)$. At each later change, $\updatefin$ either keeps $\fin_i$
or sets it to $\candidate(H_l)$, which is a prefix of $S$
(Definitions~\ref{def:candidate} and~\ref{def:lca}) and hence
of $H_0$. Thus $\mathsf{height}(\fin_i) \le \mathsf{height}(H_0)$ throughout,
and at the adoption of $H_d$, $\mathsf{height}(H_d) - \mathsf{height}(\fin_i)
\ge d$.
\end{proof}

The proposition matches the overview's own statement that the gap between
$\fin$ and the tip is unbounded during a stall
(Remark~\ref{rem:sticky-consequences}); it is the failure, in the design,
of Proposition~\ref{prop:bounded-run}, whose bound rests on the Finality
depth rule.

\begin{remark}[Final Agreement of Tenderlink]\label{rem:tenderlink-fa}
No analysis shows that Tenderlink, with a stake-weighted roster derived
from bc-chain state, has Final Agreement, that the certificates it decides
are valid, or that its notarisations are sound in the sense of
Definition~\ref{def:units}. Hypotheses~(a) and~(b) of
Proposition~\ref{prop:sl-fa} are therefore open, and with them the premise
of the Final Agreement route for the design. A BFT height of Tendermint is
decided in one or more rounds, and a validator that has voted to commit a
value in one round is locked on it in later rounds of that height until a
later round's quorum releases it. The book's argument for
Proposition~\ref{prop:bft-live}, that a proposer can always repeat its
parent's snapshot, is not reconciled with these locks, and the design's
proposer declines to repeat (Remark~\ref{rem:sl-proposer}).

Status: open problem.
\end{remark}

\section{The finaliser set}\label{sec:fs}

This section states the book's abstract finaliser model and proves the power
of a finaliser coalition in it, and then describes the finaliser set as
Shielded Labs designs it, in the vocabulary of \S\ref{sec:sl-terms}: the
finalisers and their roster, bootstrap, delegation bonds, staking windows and
the privacy of staking, rewards, slashing by user-configured hard forks,
finaliser keys, and the mechanism-design review that shaped parts of the
reward rule. Each component carries its class (Definition~\ref{def:classes}).
The design documents are Shielded Labs' \emph{Zcash Crosslink Design Overview},
draft 10 of 14 September 2026, last edited 5 October 2026 (``the overview''
below), cited by section title, and Shielded Labs' \emph{Rewards Distribution}
(22 September 2026). Calendar values and bootstrap heights are unsettled, two
current texts of the overview differing, and are given without numbers.

\subsection{The abstract finaliser model}\label{sec:fs-abstract}

In the book's model the finaliser set of a BFT epoch is the set of holders
of its $n$ voting units (Definition~\ref{def:units}); notarisation requires
$q = \lceil 2n/3 \rceil$ units cast for a proposal
(Definition~\ref{def:quorum}), and fewer than one third of the units of each
BFT epoch are ever cast non-honestly (Assumption~\ref{ass:one-third}). The
model says nothing about how voting units are assigned. Among its ``Major
Missing Components'' the book lists ``PoS subprotocol selection'' and
``Issuance and supply mechanics, such as how much ZEC stakers may earn''
(chapter ``Status and Next Steps''). Every concrete decision about the
finaliser set in the rest of this section therefore comes from Shielded
Labs' documents, and each is designed but unspecified.

\begin{remark}[Delegation goal]\label{rem:delegation-goal}
The book's chapter ``Design Goals'' states, as an ideal goal that binds no
rule, that users of shielded mobile wallets can delegate ZEC to ``PoS
consensus providers'' and earn a return ``via ZEC issuance or fees'', that
delegation may expose them to a loss of delegated ZEC ``(such as through
`slashing fees')'', and that the providers, holders of voting units, ``have
no discretionary control over such delegated funds''. The goal fixes no
mechanism; the book defines neither delegation nor slashing.
\end{remark}

\begin{remark}[Units and threshold in the two designs]\label{rem:units-designs}
In the book each BFT epoch has a fixed number $n$ of voting units, and the
notarisation rule requires at least two thirds of them, in the integer form
$q = \lceil 2n/3 \rceil$ (Definition~\ref{def:quorum}). In Shielded Labs'
design voting power is stake, and a certificate is decided when it is
``approved by two thirds of stake-weighted power'' of the finalisers whose
votes count at its BFT height (overview, \S``6. The BFT implementation'';
\S\ref{sec:sl-terms}). Read in the book's model, with the least amount of
stake as the unit, the second is an instance of the first in which the units
of a BFT epoch are the stake of the active roster at its BFT height; the
overview does not state whether two thirds is attained with equality. Both are
designed but unspecified.
\end{remark}

A set of holders \emph{withholds} in a BFT epoch when none of them casts any
of its units of that BFT epoch for any proposal. The honest $\fin$ values are
the chains $\fin_i^t$ of nodes $i$ honest at $t$ (Definition~\ref{def:fin}).

\begin{proposition}[Power of a finaliser coalition]\label{prop:coalition}
In an execution that satisfies Assumption~\ref{ass:ballots} and in which
notarisation proofs are sound (Definition~\ref{def:units}), holders of more
than one third of a BFT epoch's units can prevent every notarisation of that
BFT epoch by withholding. In every execution, whatever the holders of voting
units do, they cannot make honest $\fin$ values conflict while $\Pibc$ has
Prefix Agreement at $\sigma$ (Definition~\ref{def:pa}).

Status: proved (this volume), of a designed-but-unspecified construction.
\end{proposition}

\begin{proof}
Let the holders hold $c > n/3$ of the $n$ units of the BFT epoch and
withhold. By Assumption~\ref{ass:ballots} no party other than its holder
casts a unit, so for every proposal $P$ of that BFT epoch the units cast for
$P$ are among the $n - c$ units of the other holders, and
\[
  n - c < 2n/3 \le \lceil 2n/3 \rceil = q .
\]
No set of ballots for $P$ therefore meets the notarisation rule
(Definition~\ref{def:quorum}), and by soundness of notarisation proofs
(Definition~\ref{def:units}) no valid notarisation proof $\proofsym_P$ is
obtainable, so no bft-block $(P, \proofsym_P)$ is valid
(Definition~\ref{def:bft-block}): no proposal of the BFT epoch is notarised.

If they hold more than one third of the units of every BFT epoch from some
BFT epoch $e$ on and withhold in each, every valid bft-block is $\Obft$ or
belongs to a BFT epoch before $e$. The context of every
block $H$ of a valid bc-chain is valid (Valid context rule,
Definition~\ref{def:bc-validity}), so $\LF(H)$, a prefix of it, is valid
(Definition~\ref{def:blf}) and is such a block. For $i$ honest at $t$,
either $\fin_i^t = \Obc$ or $\fin_i^t = \candidate(\ch_i^r)$ for some
$r \le t$, a prefix of $\snapshot(\LF(\ch_i^r))$ (Definitions~\ref{def:fin}
and~\ref{def:candidate}). Hence $\fin_i^t$ is a prefix of $\snapshot(B)$ for $B
= \Obft$ or a valid bft-block $B$ of a BFT epoch before $e$. The abstract model
does not exclude that a proposal of such a BFT epoch is notarised after the
withholding begins, so this bounds finalisation by the earlier BFT epochs
without proving a stall.

The second clause is Corollary~\ref{cor:bft-subversion}: under Prefix
Agreement of $\Pibc$ at $\sigma$, Assured Finality (Definition~\ref{def:af})
holds whatever the holders of voting units do.
\end{proof}

\subsection{Finalisers and the roster}\label{sec:fs-roster}

The definitions of this subsection follow the overview, sections ``10. Kinds
of finality and which ledger state is current'', ``11. Finalizers and
rosters'' and ``8. Bootstrap and activation''.

\begin{definition}[Voting power]\label{def:voting-power}
In the Shielded Labs design the finalisers are the participants of
Tenderlink (\S\ref{sec:sl-terms}); the term widens the meaning of
Definition~\ref{def:units}: the holders of the voting units of a BFT height
are the finalisers of its active roster (Definition~\ref{def:active-roster}),
and a finaliser of the inactive roster holds none. The \emph{voting
power} of a finaliser at a best-chain state is the stake delegated to it in
that state, the total value of the bonds that target it
(\S\ref{sec:sl-terms}). Every fraction of finalisers is a fraction of stake,
never of their number, and the notarisation threshold is two thirds of the
stake of the active roster, the finalisers whose votes count at a BFT
height. The definition is designed but unspecified (overview, \S``11. Finalizers
and rosters'').
\end{definition}

\begin{remark}[Finaliser consent]\label{rem:opt-in}
A finaliser publishes an address that ``acts as a capability carrying its
own signature consenting to act as a finalizer'', and a delegation to it
requires that consent (overview, \S``11. Finalizers and rosters''). The design
has no registration transaction, no minimum self-stake and no web of trust
inside the protocol: that a finaliser is run by a given operator is left to
attestations outside it. The consent rule is designed but unspecified. No
current design fixes a minimum stake for a finaliser.

Status: open problem.
\end{remark}

\begin{definition}[Active roster]\label{def:active-roster}
Let $N$ be a fixed positive integer, the \emph{active roster size}. At a
best-chain state the \emph{active roster} consists of the $N$ finalisers of
largest voting power at that state, or of all finalisers when there are at
most $N$, and the \emph{inactive roster} of the others. The split bounds the
number of parties in a BFT round; it is neither a reward nor a penalty. The
value of $N$ and the selection rule, the order among finalisers of equal
voting power included, are open: the overview lists ``How the active roster
is selected'' among its open questions (Part VI). The definition is designed but
unspecified (overview, \S``11. Finalizers and rosters''), with $N$ and the
selection rule open.
\end{definition}

\begin{definition}[Roster derivation]\label{def:roster}
For a BFT height (\S\ref{sec:sl-terms}) greater than $1$, let $C$ be the
certificate decided at the preceding BFT height. The \emph{roster} for the
BFT height is the set of finalisers with the voting powers that the bonds in
the best-chain state at $\snapshot(C)$ determine (Definitions~\ref{def:snapshot}
and~\ref{def:voting-power}), together with the active roster of that state
(Definition~\ref{def:active-roster}); only the active roster's votes count. The
roster for BFT height $1$ is that of Definition~\ref{def:bootstrap}, which this
rule also gives when the genesis certificate's snapshot is the block of height
$h_1$. Agreement of the BFT protocol on $C$ fixes $\snapshot(C)$, the final
snapshot of the newest decided certificate, which generally lies above $\fin$
and may lie on a chain the node does not take as best
(Definition~\ref{def:sl-quantities}). The roster is required to be a pure
function of the best-chain state at that snapshot and of the slashing
configuration, the set of slash forks a node applies (\S\ref{sec:sl-terms});
under a fixed slashing configuration it is frozen while no new certificate is
decided, and a slash fork adopted during a stall changes it by jailing the
finalisers it names (Definition~\ref{def:slash-fork}). The derivation is
designed but unspecified (overview, \S``10. Kinds of finality and which
ledger state is current'', \S``11. Finalizers and rosters'').
\end{definition}

Validators that hold the same certificate at the preceding BFT height and
the same slashing configuration derive the same roster. Agreement on that
certificate is a property of Tenderlink, which no analysis establishes
(Remark~\ref{rem:tenderlink-fa}).

\begin{definition}[Bootstrap]\label{def:bootstrap}
The bft-chain has no genesis received from outside the best chain: every
node starts it from best-chain state. A \emph{roster height} $h_1$ and a
later \emph{start height} $h_2$ of the best chain fix the initial roster and
the start of BFT operation. In the numbering of \S\ref{sec:sl-terms}, BFT
height $0$ is the genesis certificate, on which no roster votes. The roster
for BFT height $1$, the first voted height (the overview's BFT height $0$),
is derived, as in Definition~\ref{def:roster}, from the best-chain state at
the block of height $h_1$. On reaching height $h_2$ a node finalises the
block of its best chain at height $h_1$ and starts the BFT protocol; every
best-chain block of height at most $h_2$ carries a nil fat pointer, which
cites no certificate (Definition~\ref{def:fat-pointer}). The difference $h_2 -
h_1$ is required to exceed the node's rollback limit, the depth beyond which it
does not reverse blocks (\emph{Consensus Guide}, \S``Rollback limits,
coinbase maturity, and finality''). The heights are unsettled: the two
current texts of the overview differ, the second copy adding a third height,
the first at which a staking action is valid. Whether the genesis bft-block
points at the best-chain genesis or at the block of height $h_1$ is open
(Part VI). The bootstrap is designed but unspecified (overview, \S``8. Bootstrap
and activation''), with the heights open.
\end{definition}

\begin{proposition}[Bootstrap safety]\label{prop:bootstrap-safety}
If $\Pibc$ has Prefix Agreement at depth $h_2 - h_1$ (Definition~\ref{def:pa}),
any two honest nodes whose best chains reach height $h_2$ hold the same block at
height $h_1$, finalise it, and, with the same slashing configuration, derive the
same initial roster.

Status: proved (this volume), of a designed-but-unspecified construction.
\end{proposition}

\begin{proof}
Let nodes $i$ and $j$ reach height $h_2$ at times $t$ and $u$ at which they
are honest, their best chains then having tips at height $h_2$ or above.
Prefix Agreement at $h_2 - h_1$ makes $\ch_i^t\trunc{h_2 - h_1}$ and
$\ch_j^u\trunc{h_2 - h_1}$ agree; both have tips at height at least $h_1$, so
the shorter is a prefix of the longer and both hold the same block at height
$h_1$. Each node finalises the block of its best chain at height $h_1$ at
that time (Definition~\ref{def:bootstrap}), and the initial roster is
a function of the best-chain state at it and of the slashing configuration
(Definition~\ref{def:roster}).
\end{proof}

\begin{remark}[Rollback-limit argument]\label{rem:bootstrap-rollback}
The overview argues instead from rollback limits (\S``8. Bootstrap and
activation'', ``Safety argument''): when a node reaches $h_2$, its block at
height $h_1$ lies below its rollback limit, ``and therefore identical on every
chain a node could be following'', so every node computes the same roster from
it. The argument rests on the rollback limit, which is node policy: consensus
imposes no maximum reorganisation depth (\emph{Consensus Guide}, \S``Rollback
limits, coinbase maturity, and finality''), and the rollback window that
ZIP 218 recommends is a SHOULD on implementations (ZIP 218,
``Block-count-based constants''). It is the only safety argument the design
has for activation, and the overview calls the finality it relies on
``reorg-depth finality rather than BFT finality''. The clause ``identical on
every chain a node could be following'' is a hypothesis on the best chains
of distinct honest nodes, not a consequence of a node-local limit: a limit
stops a node from abandoning its own block at height $h_1$, and does not
stop two nodes whose best chains forked below $h_1$ from holding different
ones. The argument thus shows only that a node does not abandon its own
block at height $h_1$; Prefix Agreement of the best chain at depth
$h_2 - h_1$ supplies the missing hypothesis
(Proposition~\ref{prop:bootstrap-safety}).

Status: argued (Shielded Labs).
\end{remark}

\subsection{Delegation bonds}\label{sec:fs-bonds}

\begin{definition}[Bond]\label{def:bond}
A \emph{bond} is created by a transaction that spends shielded funds and
carries a staking action, the creation of a delegation bond, which states an
amount, required to be a power of ten ZEC, and a target finaliser, named
through its consent (Remark~\ref{rem:opt-in}). The amount is quantised at
creation only. The bond is a free-floating note, in the overview's term: a
ledger object, not a note of a shielded pool, controlled by a signing key
derived from the creating action and a salt, so that identical delegations
yield distinct bonds. Holding the key is holding the bond, which is
decoupled from the funds that created it. Bonds are atomic: no bond is
split, and every later action applies to a whole bond. The definition is
designed but unspecified (overview, \S``12. Delegation bonds'').
\end{definition}

\begin{definition}[Staking actions]\label{def:staking-actions}
Besides its creation, a bond admits three staking actions. A
\emph{retarget} changes the bond's target finaliser; it may occur at any
time, and its transaction records both the old and the new target, so that
from the bond state at any time the targets at every earlier time are
recovered by replaying the retargets backwards, which slashing requires. An
\emph{exit} is an \emph{unbond} followed by a \emph{withdraw}, in two
different staking windows, the fixed height ranges to which creation,
unbonding and withdrawal are confined. The unbond fixes the bond's numeric
value, which depends on the block it lands in; the withdraw must state that
value exactly, or it is invalid. The overview gives two reasons for the
separation: the amount is an explicit number before the funds re-enter
ordinary transactions, and the delay puts stake at risk, since ``if an
attack is discovered there is time to slash before the funds escape into the
shielded pool''. The actions are designed but unspecified (overview, \S``12.
Delegation bonds'').
\end{definition}

\begin{remark}[Transaction contents]\label{rem:bond-tx}
\emph{Rewards Distribution}, ``Staking Mechanics'', requires that a
transaction acting on bonds contain no inputs or outputs other than its bond
actions, actions of the ``latest/greatest shielded pool'' and the payment of
its fee. The document does not name the pool: whether it is the Ironwood
pool (\emph{Ironwood Guide}, \S``The Orchard protocol and its two
pools'') or any shielded pool is open. In-protocol liquid staking
derivatives are not a goal of the initial design: the book lists them among
its out-of-scope goals, which it says ``may be goals for future protocol
improvements'' (chapter ``Design Goals''), and the overview defines none.
ZIP 248 (Draft), the extensible transaction format, lists hypothetical bundle
types ``Staking'' and ``Unstaking (if it can't be combined with the Staking
bundle)'' with no defining ZIP, and no ZIP that ZIP 259 lists for NU7 provides
the transaction format of staking actions or the header change that the
cross-link fields require. The restriction is designed but unspecified; the
ZIP 248 entries are specified text that defines no bundle.
\end{remark}

\begin{remark}[Open bond questions]\label{rem:bond-open}
Two questions about bonds are undecided. First, whether a withdraw may
complete before the unbond it follows is finalised, and whether an exit may
complete during a finalisation stall. Second, how a wallet derives bond keys
from its seed: no design states a derivation, by ZIP 32 (\emph{Wallet
Guide}, \S``Hierarchical key derivation (ZIP 32)'') or otherwise, so
the recovery of bonds from a seed is undesigned, and the overview lists as
open the wallet's handling of the calendar: completing a withdraw without a
live wallet and scheduling an action from outside a staking window
(Part VI).

Status: open problem.
\end{remark}

\begin{remark}[Redelegation as censure]\label{rem:redelegation}
Short of a slash fork, a staker censures a finaliser by retargeting its
bonds away from it, at any height, or by exiting them, in staking windows
(Definitions~\ref{def:staking-actions} and~\ref{def:staking-window}). Shielded
Labs' role allocation gives stakers the selection of finalisers
(Remark~\ref{rem:sl-rationale}), and the overview describes the threat of a
slash fork as what moves stake, ``the burn is the fallback'' (\S``16. Social
slashing'', ``A big stick''). Two limits bound the means. A retarget or exit
changes voting power from the block that contains it, but the votes that count
only through a later roster, and, short of a slash fork, the roster changes
only through decided certificates (Definition~\ref{def:roster}), so during a
stall it is frozen and no redelegation reaches it. A party with more than
half of the hash
rate can exclude transactions at will (\emph{Consensus Guide}, Remark ``Majority
is a different problem'', in \S``Reward-normalised exposure''), redelegations
among them, and so disable the means. The two limits are this volume's
consequences of Definition~\ref{def:roster} and the \emph{Consensus Guide}'s
remark; the role allocation is Shielded Labs' argument.

Status: argued (Shielded Labs).
\end{remark}

\subsection{Staking windows and privacy}\label{sec:fs-windows}

\begin{definition}[Staking window]\label{def:staking-window}
A fixed calendar divides the best-chain heights into recurring \emph{staking
windows}, the overview's staking days, and the heights between them. A
creation, unbond or withdraw takes effect only in a block whose height lies
in a staking window (overview, \S``13. Staking days''); a retarget takes
effect at any height (\S``12. Delegation bonds''). The overview gives the
windows three roles: batching removes timing information
that could link actions to users; the separation of unbond and withdraw into
different windows gives the exit delay; and the slash window of a slash fork
(Definition~\ref{def:slash-fork}) is measured in staking windows. The period of
the calendar and the length of a window are unsettled. The definition is
designed but unspecified (overview, \S``13. Staking days''), with the calendar
open.
\end{definition}

\begin{remark}[Calendar divergence]\label{rem:calendar}
The two current copies of the overview, its main text and a second copy
whose staking-calendar and bootstrap text, revised on 29 September 2026, is
newer than the main text's of 22 September, give different periods.
\emph{Rewards Distribution} gates a different set of actions: it confines
creation and its own withdraw, which takes the place of the overview's unbond,
to staking windows, and leaves unconstrained its transfer, which follows a
``sufficient delay'' that it leaves to be defined and which takes the place of
the overview's withdraw (the volume's reading of the two lifecycles). Its text
is inconsistent: item 2 of ``Staking Mechanics'' confines every transaction with
bond actions to staking-day heights, while item 3 leaves redelegation and
transfer unconstrained. The public \emph{Crosslink FAQ} (modified 30 July 2026)
describes an older scheme of fixed periods.
The volume gives no calendar value.

Status: open problem.
\end{remark}

\begin{remark}[Privacy design]\label{rem:staking-privacy}
Public are each finaliser's total stake, since ``that total is what voting
weight is computed from, so it has to be'', and, for each bond, its
power-of-ten amount, its target, its retargets with old and new target, its
withdrawal value and the staking window of each action. Hidden is the link
from a bond to the shielded notes that funded it and to its holder's other
activity: ``bonds are anonymous''. Two devices serve this: quantisation of
amounts, so that bonds cannot be fingerprinted by size, and batching into
staking windows (overview, \S``12. Delegation bonds'', \S``13. Staking days'').
The exit delay of Definition~\ref{def:staking-actions} serves security. No
design shields the delegation itself, which would hide a bond's amount and
target. The two devices are designed but unspecified; shielded delegation
is an open problem.
\end{remark}

\begin{remark}[Open privacy and expiry questions]\label{rem:staking-open}
Quantisation protects creation only: once a bond has earned rewards, which
are added to it, its value at withdrawal is not a power of ten, and no source
discusses the linkage that this value permits. No source couples
transaction expiry to the calendar. A transaction carrying a creation,
unbond or withdraw is minable only at heights inside a staking window
(Definition~\ref{def:staking-window}) and, if its expiry height
$h_{\mathrm{exp}}$ is not $0$, at no height above $h_{\mathrm{exp}}$
(\emph{Wallet Guide}, \S``Expiry (ZIP 203)'').
If it is not mined before its window ends and $h_{\mathrm{exp}}$ lies beyond
that end, it cannot be mined until the next window opens although it has not
expired, and it expires unmined if no later window opens at a height at most
$h_{\mathrm{exp}}$; with $h_{\mathrm{exp}} = 0$ it never expires and waits for
the next window (the volume's synthesis).

Status: open problem.
\end{remark}

\subsection{Rewards}\label{sec:fs-rewards}

\begin{definition}[Reward source]\label{def:reward-source}
Staking rewards are paid from the block subsidy (\emph{Consensus Guide},
\S``Scheduled issuance with shorter blocks''), with no issuance beyond what
the existing rules provide: in the words of Shielded Labs' \emph{Crosslink
FAQ}, ``The 21 million coin cap remains unchanged''. The distribution of
transaction fees is unchanged: ``Crosslink does not alter the distribution of
transaction fees in any way'' (\emph{Rewards Distribution}, ``Crosslink
Rewards Distribution Rules''). The block subsidy remaining after the share of
Zcash's development fund, the overview's post-dev-fund issuance, is divided
between miners and staking (\emph{Rewards Distribution} calls the remainder
the Operational Consensus Rewards); the part paid for staking is the
\emph{staking share}. No source states whether the additional block subsidy
by which ZIP 237 (Draft) reissues the reserve of the Network Sustainability
Mechanism, the NSM (\emph{Consensus Guide}, \S``The NSM reserve and delayed
reissuance''), is divided in the same way. The definition is designed but
unspecified.
\end{definition}

% issuance.py: development fund 20% of subsidy (ZIP 214 streams 8% + 12%); miners 40% and stakers 40% of subsidy = 50% each of post-development-fund issuance (0.8 x 0.5 = 0.4); bonds 36%, commission 4% of subsidy: ok
\begin{remark}[The published split]\label{rem:split}
The published split gives 40 per cent of the block subsidy to miners, 40 per
cent to stakers and 20 per cent to Zcash's development fund
(\emph{Crosslink FAQ}). The 20 per cent is not a Crosslink share: it is
ZIP 214's two funding streams, of 8 and 12 per cent of the block subsidy,
which end at the third halving (\emph{Consensus Guide}, \S``Scheduled
issuance with shorter blocks''). The 40 and 40 per cent of the subsidy are
the 50 and 50 per cent of post-development-fund issuance of \emph{Rewards
Distribution}, since $0.8 \times 0.5 = 0.4$. Inside the staking share,
90 per cent goes to bonds pro rata and 10 per cent to finalisers as a
commission fixed by the protocol; as fractions of the block subsidy these
are 36 and 4 per cent. The split is designed but unspecified.
\end{remark}

% issuance.py: overview rule: miner 48%, staking 48%, first-inclusion bounty 4% of post-development-fund issuance; 52% unpaid without a qualifying certificate: ok
\begin{remark}[Unsettled payment rule]\label{rem:payment-open}
The two texts disagree. The overview (\S``14. Rewards'', ``Proposed issuance
split'' and ``Proposed issuance rule'') divides post-development-fund
issuance into 48 per cent miner subsidy, 48 per cent staking share, divided
90 to 10 as above, and a 4 per cent bounty to the miner that includes a new
finality. It pays the staking share and the bounty only if a best-chain
block ``points at a new certificate'' that finalises a best-chain block
``close to the tip''; otherwise the 52 per cent is burned through the NSM.
The overview measures that distance on the best chain, as the height of the
citing block minus the height of the snapshot of the certificate it cites,
which is the citing block's finality depth (Definition~\ref{def:fdepth};
Remark~\ref{rem:sl-final-each}), and leaves undefined the bound that
``close'' denotes (Part VI). \emph{Rewards Distribution} pays 50 and 50 per
cent, then 90 and 10 per cent, in every block. The texts also disagree on
the commission: the overview pays it to active finalisers ``uniformly'',
\emph{Rewards Distribution} in proportion to the total stake weight of each
active finaliser.

Status: open problem.
\end{remark}

\begin{definition}[Commission and bond earnings]\label{def:commission}
The \emph{commission} is 10 per cent of the staking share, fixed by the
protocol and not set by each finaliser. Whether it is divided uniformly
among active finalisers or in proportion to their stake, and whether its
denominator is the active roster's stake or the total stake, are open
(overview, \S``14. Rewards''; \emph{Rewards Distribution}, ``Crosslink Rewards
Distribution Rules''); the text of the \emph{Crosslink FAQ} on ``the
commission rate charged by their selected finalizer'' is stale. A bond earns
from the other 90 per cent whether or not its finaliser is in the active
roster: ``Backing an up-and-coming finalizer costs nothing''. A finaliser's
commission accrues to its lockbox (\S\ref{sec:sl-terms}), which a staking
action converts into a bond and which counts, in the 90 per cent pool, as a
virtual bond delegated to its finaliser. Rewards increase a bond's
principal, and no yield is guaranteed: ``Rewards come from a fixed pool''.
Quantisation applies at creation only (Definition~\ref{def:bond}), so this
growth does not conflict with it. The definition is designed but unspecified
(overview, \S``14. Rewards''), with the division and the denominator open.
\end{definition}

\begin{remark}[Conditional rewards]\label{rem:conditional-rewards}
The overview's single-shot rule (Remark~\ref{rem:payment-open}) answers
three recommendations of N.~Della Penna's mechanism-design review
(\S\ref{sec:intro-sources}; \emph{the review} below): a
first-inclusion bounty for miners (R2); an anti-jackpot constraint (R4), in
that nothing accrues while finalisation is stalled, each payment being made
once, in its block; and, in the overview's words, certificate aging ``by
another route'' (R6): where the review pays a bounty for any certificate
that advances finality, optionally weighted by the blocks advanced, the
overview withholds rewards from finalisations far behind the tip, to
encourage catch-up on a shared prefix. The review's anti-jackpot results,
its Theorem 3.3 and Proposition 3.5, belong to its own model
(Remark~\ref{rem:review-scope}). Any reward trigger must be
objective, computable from chain data: consensus never reads the node-local
$\fin$ (overview, \S``10. Kinds of finality and which ledger state is current'';
Definition~\ref{def:sl-quantities}). Rounding must be deterministic, every
party agreeing exactly on the value paid to each recipient, and its rule is
unfinished (overview, \S``15. Accounting and the acceleration structure'';
\emph{Rewards Distribution}). These constraints are designed but unspecified;
the rounding rule is open.
\end{remark}

\begin{remark}[Withheld rewards and the NSM]\label{rem:nsm}
ZIP 237 (Draft; listed by ZIP 259 among the NU7 consensus changes) defines
the NSM Value Balance in its Terminology as funds removed via ZIP 235 or
ZIP 233, plus block subsidy and fees left unclaimed by coinbase transactions
before NU6. Its recurrence for that balance, in ``Changes to the Zcash
Protocol Specification'', adds $\removed(\mathsf{height})$, ``the total value
removed from circulation in the block at $\mathsf{height}$ by any deployed
mechanism (such as ZIP 233 or ZIP 235)''. A Crosslink ZIP that removed
withheld staking rewards would therefore enter the recurrence through
$\removed(\mathsf{height})$, with an amendment of the Terminology only.
Under NU7 as specified no such mechanism exists: ZIP 259 states that ``NU7
does not deploy ZIP 233 or ZIP 234'', and a coinbase cannot leave value
unclaimed since ZIP 236 (\emph{Consensus Guide}, \S``The NSM reserve
and delayed reissuance''). These statements about ZIP 237 are specified. The
route by which withheld staking rewards would reach the NSM is designed by
no text.

Status: open problem.
\end{remark}

\subsection{Slashing}\label{sec:fs-slashing}

The book defines no slashing. Its out-of-date chapter ``Potential Changes to
Crosslink'' notes that a stake requirement for proposals ``potentially
allows parties that make invalid proposals to be slashed'', and its chapter
``Questions About Crosslink'' sketches evidence of a safety violation, two
bft-blocks final in some context with the same parent, or more than one
third of stake voting for conflicting blocks of one BFT epoch, posted in a
transaction to the best chain. Neither is adopted, and the ``slashing fees''
of the delegation goal (Remark~\ref{rem:delegation-goal}) have no mechanism.

\begin{remark}[No in-protocol slashing]\label{rem:no-slashing}
The Shielded Labs design has no automatic slashing in the protocol. The
stated reason is a property of Tenderlink: peers agree on the decision for
a certificate, but different peers may hold different subsets of the votes
that made up the two thirds, so ``the protocol therefore cannot use `who
voted for what' as evidence for anything --- not slashing, not per-vote
payouts'' (overview, \S``6. The BFT implementation'', \S``11. Finalizers and
rosters''). The property supports a narrower claim: the exact signature set
is not agreed, so the absence of a finaliser's signature from a fat pointer
is no evidence that it did not vote, whereas a signature that is present is
attributable to its signer under Assumption~\ref{ass:ballots}. The
absence of in-protocol slashing is designed but unspecified.
\end{remark}

\begin{definition}[Slash fork]\label{def:slash-fork}
A \emph{slash fork} is a user-configured hard fork. Node operators set a
configuration that names one or more finalisers and a best-chain activation
height; every slash is a hard fork, with no default switch. On the
best-chain side every bond delegated to a named finaliser within a window
reaching back two staking days, that is two staking windows
(Definition~\ref{def:staking-window}), from the activation height is burned,
bonds unbonding from it or retargeted away from it within the window
included; the retarget records of Definition~\ref{def:staking-actions} make
that set recoverable, and there is no partial slashing and no tunable
amount. On the BFT side the named finalisers are jailed: they vote in no
current or later round, and their voting power leaves the pool, so that the
others can reach two thirds. Each certificate carries a best-chain height,
``do not include by'', below which no best-chain block may cite it, so that
certificates of the reduced set reach the best chain no earlier than the
activation height (the exact comparison and encoding are open, Part VI). The
height is non-decreasing along the bft-chain and raised only by the first
certificate of a forked bft-chain, which embeds the configuration; the
forked finalisers thus resume at once while the best chain changes at the
one activation height. Votes are cryptographically namespaced by configuration.
The configuration obeys hard constraints: the roster stays a pure function of
the best-chain state at the final snapshot of the BFT protocol and of the
configuration (Definition~\ref{def:roster}); staking actions that have landed in
the best chain are respected; and application is robust to users, miners and
finalisers adopting the configuration at different times. The definition is
designed but unspecified (overview, Part V, \S``16. Social slashing'').
\end{definition}

\begin{remark}[Open slashing details]\label{rem:slash-open}
Several interactions are marked as unsettled: slashing for misbehaviour other
than stalls (``Other malicious behaviour''), the recognition of peers on
another configuration as well behaved, and certificates produced between the
creation of a configuration and its application (overview, \S``16. Social
slashing''; Part VI). A withdrawn bond cannot be burned. The exit and the
slash window are both counted in staking windows: an exit needs two different
windows (Definition~\ref{def:staking-actions}), and the slash window reaches
back two. A bond that unbonds from a named finaliser in the earlier window of
the slash window may therefore withdraw in the later one, below the activation
height, whatever the calendar. Such a bond belongs to the set of
Definition~\ref{def:slash-fork} but escapes the burn, and no source states how
the burn treats it.

Status: open problem.
\end{remark}

\subsection{Finaliser keys}\label{sec:fs-keys}

\begin{remark}[Keys and signatures]\label{rem:keys}
Finaliser keys, the signatures in fat pointers and bond keys are Ed25519
(overview, \S``3. The proof-of-work side: the fat pointer and where it
lives'', \S``12. Delegation bonds''). The signature set of a fat pointer has
variable size. A fixed-size signature set for the fat pointer, by aggregate
signatures or by a zero-knowledge proof over the signatures, is open, and the
post-quantum story for the fat pointer, bond keys and finaliser keys is ``not
yet considered'' (Part VI). In this design Assumption~\ref{ass:ballots}, on
which Proposition~\ref{prop:coalition} rests, is a hypothesis on the Ed25519
signatures of finalisers and on the custody of their keys. The use of
Ed25519 is designed but unspecified; a fixed-size signature set and
post-quantum keys are open problems.
\end{remark}

\begin{remark}[Key rotation]\label{rem:key-rotation}
Assumption~\ref{ass:one-third} counts every unit ever cast non-honestly, a
cast with a compromised key included (Definition~\ref{def:nonhonest}). The
book therefore requires that ``validator keys of honest nodes must remain
secret indefinitely'' and that, ``whenever a key is rotated, the old key must
be securely deleted'' (chapter ``The Crosslink 2 Construction'', section
``Model for BFT protocols''), and the book's issue on the point is open.
Neither the book nor the overview has a mechanism to rotate or revoke a
finaliser key, or a threshold-custody design for finaliser keys.

Status: open problem.
\end{remark}

\subsection{The mechanism-design review}\label{sec:fs-review}

\begin{remark}[Scope of the review]\label{rem:review-scope}
The review examined a Shielded Labs design text of 11 December 2025, which
predates the overview. Its results are stated in its own abstract model,
with proofs of varying completeness: rational agents with a discount factor, a
validator set frozen at the last finalised block, a finality function from
chains to heights, and certificates valid when their signers hold stake above a
fraction $\tau \ge 2/3$ of the total, $\tau$ the review's threshold fraction.
They are not theorems about a Crosslink rule set, and the volume uses them as
framing only. Its Theorem 6.4 (``Fork Choice Dichotomy'') concerns a
work-preferring and a finality-preferring rule. Sticky fork choice
(Definition~\ref{def:sticky}) selects by work among the chains that contain
the node-local $\fin$, not among chains extending the highest known
finalised block, and does not halt the best chain during a stall
(Remark~\ref{rem:sticky-consequences}); it is neither case as stated, and the
theorem is not read as constraining it. Its Theorem 2.3 (``Accountable
Safety'') is quorum intersection, the argument of Lemma~\ref{lem:uniqueness} in
units of stake.
\end{remark}

\begin{remark}[Results and failure modes]\label{rem:review-results}
As framing, the review's results that bear on the finaliser set are these.
Conflicting finality requires that stake of at least $2\tau - 1$ times the
total have signed conflicting certificates (Theorem 2.3). Without slashing,
the largest penalty for a deviation is at most the withheld rewards plus the
time value of a delayed exit (Theorem 4.2); the overview's slash forks lie
outside that assumption, and short of a slash fork the design is its case.
A frozen validator set whose participating stake falls below the threshold
during a stall cannot recover without intervention from outside the protocol
(Theorem 5.4); the jailing of named finalisers by a slash fork
(Definition~\ref{def:slash-fork}), which removes their voting power so that
the others can reach two thirds, is an intervention of that kind. If exits
can complete during a stall, the cost of corrupting the frozen set decays
while its voting weights stay fixed (Proposition 5.8), the first question of
Remark~\ref{rem:bond-open}. The review names three failure modes: the
``Hostage Holdout'', a blocking coalition delaying finality to increase
future payouts; the ``Zombie Set'', a frozen set made permanently inoperable
by attrition; and the ``Liquid Exit Paradox'', exits during a stall that
leave the frozen set with ``phantom security''. Of its recommendations, the
overview adopts the first-inclusion bounty (R2) and the anti-jackpot
constraint (R4) and states that it meets certificate aging (R6) by another
route (Remark~\ref{rem:conditional-rewards}); it does not adopt
signer-conditioned commission (R1), since it pays active finalisers without
regard to vote sets (Remark~\ref{rem:no-slashing}). Stated in the review's
model, these results are not statements about Crosslink~2 or the overview's
rules, and carry no status word of this volume.
\end{remark}

\section{The client contract}\label{sec:contract}

This section states what a client of a Crosslink~2 node may rely on. It
defines the word \emph{final} for a best-chain block and for a transaction,
the only definition of the word for these objects in the volume, and proves
what finality at an honest node guarantees, how the finality status of a
block is decided, and that the finalised chain is valid. It then separates
finality from confirmation depth, states which chain an application reads in
each design and what a finalisation stall does to it, records what changes
for the confirmation policy, anchors and expiry of the \emph{Wallet Guide},
and states why no current design lets a client that checks headers only
verify finality. The constructions the section uses are designed but
unspecified (Definition~\ref{def:classes}); the deployed facts it cites are
specified.

\subsection{Finality as a status}\label{sec:k-final}

\begin{definition}[Final]\label{def:final}
A bc-block $X$ is \emph{final at node $i$ at time $t$} iff
$X \preceqbc \fin_i^t$, where $\fin_i^t$ is the locally finalised chain of
Definition~\ref{def:fin}. A transaction is \emph{final at node $i$ at time
$t$} iff a block that contains it is final at $i$ at $t$, that is, iff it is
a transaction of a block of $\fin_i^t$.
\end{definition}

The notion is the book's assured finality: a final transaction cannot be
reverted by the protocol, while a fork made outside the protocol can still
reverse it (the book, chapter ``Terminology'', entries ``Final'' and
``Assured Finality''); a finality override is such a fork
(Remark~\ref{rem:override-outside}). It is designed but unspecified, and it
is distinct from four other uses of the word.
\begin{itemize}
\item Finality in the context of a bft-block
  (Definition~\ref{def:final-context}) is a relation between two
  bft-blocks; finality of a bc-block is relative to one node's $\fin_i^t$.
\item The final treestate of a block is the output treestate of its last
  transaction (\emph{Ironwood Guide}, \S``Anchors'', Definition ``Treestate
  and anchor''), an object and not a status.
\item The final chain of the \emph{Wallet Guide}, \S``Reorganisation
  detection and recovery'', Theorem ``Correctness of synchronisation'', is
  the server's best chain from the time at which it stops changing, a
  hypothesis on an execution; its Lemma ``Expired unmined is final on its
  chain'' (\S``Status tracking and terminal states'') concludes that an
  expired transaction lies in no chain that agrees with a given chain up to
  the expiry height. Neither is a status that a node computes.
\item Confirmation by depth gives probabilistic assurance only: no finite
  confirmation count makes the modelled probability of reversal zero
  (\emph{Consensus Guide}, \S``Confirmation requirements'', Proposition ``No
  finite count is final'').
\end{itemize}

\begin{proposition}[Path to finality]\label{prop:path-final}
Let node $i$ be honest at time $t$ and let $X \neq \Obc$ be a bc-block. Then
$X \preceqbc \fin_i^t$ iff there are a proposal $P$, bft-blocks $B$ and $C$
and a bc-block $H$ such that the following hold, in order.
\begin{enumerate}[label=(\arabic*)]
\item $X \preceqbc \snapshot(P)$.
\item $B = (P, \proofsym_P)$ is a bft-block formed from $P$
  (Definition~\ref{def:bft-block}).
\item $B$ is final in the context of the valid bft-block $C$,
  $B \preceqbft \bftlastfinal(C)$ (Definition~\ref{def:final-context}).
\item $H$ is a valid bc-block whose context has a last final block that
  extends $B$, $B \preceqbft \LF(H)$ (Definition~\ref{def:lf}).
\item $X \preceqbc H\trunc{\sigma}$.
\item At an update of node $i$ at a time at most $t$, $H$ becomes its best
  tip with $\fin_i^s \preceqbc \candidate(H)$ for the value $\fin_i^s$ before
  the update, so that $\updatefin$ sets $\fin_i$ to $\candidate(H)$
  (Definition~\ref{def:fin}).
\end{enumerate}

Status: proved (this volume), of a designed-but-unspecified construction.
\end{proposition}

\begin{proof}
Suppose (1) to (6). By the Valid context rule and Definition~\ref{def:blf},
$\LF(H)$ is valid; every block of the bft path from $B$ to $\LF(H)$ is then
valid by the Parent rule, and the Linearity rule applied along the path gives
$X \preceqbc \snapshot(B) \preceqbc \snapshot(\LF(H))$
(Definitions~\ref{def:bc-validity} and~\ref{def:bft-validity}). With~(5), $X$ is
a common prefix of $\snapshot(\LF(H))$ and $H\trunc{\sigma}$, so $X \preceqbc
\candidate(H)$ (Definitions~\ref{def:candidate} and~\ref{def:lca}). By~(6) the
update sets $\fin_i$ to $\candidate(H)$, and Lemma~\ref{lem:fin-linear}, node
$i$ having followed the update rule up to $t$ (Definition~\ref{def:exec}), gives
$X \preceqbc \fin_i^t$.

Conversely, if $X \preceqbc \fin_i^t$, then $\fin_i^t \neq \Obc$, so the
last update up to $t$ that changed $\fin_i$ set it, in the first case of
$\updatefin$, to $\candidate(H)$ for its new best tip $H$, which is~(6).
Since $\candidate(H) \neq \Obc$, $\LF(H) \neq \Obft$. With $B = \LF(H)$,
$C = H.\contextbft$ and $P$ the proposal of $B$, conditions~(2) to~(4) hold,
$H$ being valid as the best chain of an honest node, and
$X \preceqbc \candidate(H)$ gives~(5) and, through $\candidate(H) \preceqbc
\snapshot(\LF(H)) = \snapshot(P)$, also~(1).
\end{proof}

The steps are those of Construction~\ref{def:node-update}, in which a BFT
decision alone changes nothing locally, and Proposition~\ref{prop:lag} gives
the distance from the tip to $\fin_i$ at such an update. By the Tail
Confirmation rule (Definition~\ref{def:bft-validity}), which $B$ satisfies
since it is valid by~(4), and the paragraph after Definition~\ref{def:snapshot},
the block whose header is the newest that $P.\headersbc$ carries lies $\sigma$
blocks above $\snapshot(P)$, hence at least $\sigma$ blocks above $X$. The valid
notarisation proof of $B$ requires, under soundness
(Definition~\ref{def:units}), ballots for $P$ authenticated by the holders
of at least two thirds of the voting units of $P$'s BFT epoch, an honest
holder casting only if $\snapshot(P)$ is a prefix of its own best chain
truncated by $\sigma$ (Construction~\ref{def:honest-voting}). In adapted
Streamlet (Example~\ref{ex:streamlet}), condition~(3) is a run in $C$ of
three adjacent blocks of consecutive BFT epochs whose middle block has $B$
as a prefix; in a protocol that decides each block individually, $C = B$. A
candidate that conflicts with $\fin_i^s$ leaves $\fin_i$ unchanged and is
recorded as a finalisation safety hazard.

\begin{proposition}[Guarantee of finality]\label{prop:final-guarantee}
Let a bc-block $X$ be final at node $i$ at time $t$, with $i$ honest at $t$.
Then $X$ is final at $i$ at every time $u \ge t$ at which $i$ is honest. In
every execution in which $\Pibc$ has Prefix Agreement at depth $\sigma$
(Definition~\ref{def:pa}) or $\Pibft$ has Final Agreement
(Definition~\ref{def:fa}), moreover, $X$ and $\fin_j^u$ agree for every time
$u$ and every node $j$ honest at $u$.

Status: proved (this volume), of a designed-but-unspecified construction.
\end{proposition}

\begin{proof}
A node honest at $u$ has followed the protocol, $\updatefin$ included, up to
and including $u$ (Definition~\ref{def:exec}). For $u \ge t$,
Lemma~\ref{lem:fin-linear} gives $\fin_i^t \preceqbc \fin_i^u$, so $X
\preceqbc \fin_i^u$ by transitivity of the prefix order. Under either
premise, Corollary~\ref{cor:af-either} gives Assured Finality
(Definition~\ref{def:af}), so $\fin_i^t$ and $\fin_j^u$ agree. If $\fin_i^t
\preceqbc \fin_j^u$, then $X \preceqbc \fin_j^u$. Otherwise $\fin_j^u
\preceqbc \fin_i^t$, and $X$ and $\fin_j^u$ are both prefixes of
$\fin_i^t$, so they agree by Lemma~\ref{lem:linear-prefix}.
\end{proof}

Agreement of $X$ with $\fin_j^u$ means that $X$ is final at $j$ at $u$ or
that $\fin_j^u$ is a strict prefix of $X$: no honest node's finalised chain
conflicts with $X$. The proposition guarantees nothing to a node that is not
honest in the model, which includes a node that the adversary has cut off
from the rest of the network, for example by an eclipse
(Remark~\ref{rem:exec-props}), and nothing in it binds a fork made outside
the protocol (Remark~\ref{rem:override-outside}). Its premises are
hypotheses on the sub-protocols: Prefix Agreement of $\Pibc$ at $\sigma$ rests
on an assumption for Zcash proof of work (Remark~\ref{rem:pow-prefix}) and on a
transfer from $\Piorigbc$ to $\Pibc$ that is unproved
(Remark~\ref{rem:origbc-transfer}); Final Agreement of $\Pibft$ is a hypothesis
on $\Piorigbft$ carried over by an argued step (Proposition~\ref{prop:erasure}),
and for the BFT protocol of Shielded Labs' design it is open
(Remark~\ref{rem:tenderlink-fa}).

\begin{proposition}[Finality status]\label{prop:status-tri}
Let node $i$ follow $\updatefin$, let $B$ be a bc-block and $t$ a time.
Exactly one of the following holds:
\begin{enumerate}[label=(\roman*)]
\item \emph{final}: $B \preceqbc \fin_i^t$;
\item \emph{not yet final}: $\fin_i^t \precbc B$;
\item \emph{never final}: $B$ and $\fin_i^t$ conflict, and then $B$ is final
  at $i$ at no time $u \ge t$.
\end{enumerate}

Status: proved (this volume), of a designed-but-unspecified construction.
\end{proposition}

\begin{proof}
By Definition~\ref{def:prefix}, $B$ and $\fin_i^t$ either agree or conflict.
If they agree, then $B \preceqbc \fin_i^t$, case~(i), or $\fin_i^t
\preceqbc B$ with $\fin_i^t \neq B$, case~(ii); the two exclude each other,
since $B \preceqbc \fin_i^t$ and $\fin_i^t \preceqbc B$ together give $B =
\fin_i^t$. Otherwise they conflict, case~(iii), which excludes both. In
case~(iii), suppose $B \preceqbc \fin_i^u$ for some $u \ge t$.
Lemma~\ref{lem:fin-linear} gives $\fin_i^t \preceqbc \fin_i^u$, so $B$ and
$\fin_i^t$ are both prefixes of $\fin_i^u$ and agree by
Lemma~\ref{lem:linear-prefix}, a contradiction.
\end{proof}

The status is decided by the prefix relation, not by height. A block of
height greater than $\mathsf{height}(\fin_i^t)$ is in case~(ii) only if
$\fin_i^t$ is its ancestor; a block on a branch that forks from $\fin_i^t$
below its tip is in case~(iii) whatever its height, while a classification by
height alone would place it in case~(ii) if its height exceeds
$\mathsf{height}(\fin_i^t)$ and in case~(i) otherwise. A block of height at most
$\mathsf{height}(\fin_i^t)$ is in case~(i) iff it is the block of
$\fin_i^t$ at its height, and in case~(iii) otherwise. Case~(ii) promises
nothing: a later update may move $\fin_i$ to a chain that conflicts with
$B$. A transaction contained in $B$ has, relative to $B$, the status of $B$;
a transaction contained in no block has none of the three statuses and is
not final.

\begin{proposition}[Validity of the finalised chain]\label{prop:fin-valid}
For node $i$ honest at $t$, the chain $\fin_i^t$ is bc-valid
(Definition~\ref{def:bc-validity}), and every transaction of $\fin_i^t$, its
anchors included, satisfies the transaction consensus rules in the context of
its own chain, the prefix of $\fin_i^t$ that its block tips.

Status: proved (this volume), of a designed-but-unspecified construction.
\end{proposition}

\begin{proof}
Lemma~\ref{lem:fin-depth} gives a time $r \le t$ with $\fin_i^t \preceqbc
\ch_i^r\trunc{\sigma} \preceqbc \ch_i^r$. Node $i$ is honest at $r$, having
followed the protocol up to $t \ge r$ (Definition~\ref{def:exec}), so its
best chain $\ch_i^r$ is a bc-valid chain, a chain of bc-valid blocks
(Definitions~\ref{def:bc-protocol} and~\ref{def:bc-validity}), and every
prefix of $\ch_i^r$, $\fin_i^t$ among them, is a bc-valid chain.
Validity in $\Pibc$ includes the Inherited origbc rules,
under which the transaction consensus rules are those of $\Piorigbc$
unchanged, with no reordering or sanitisation and heights preserved
(Definition~\ref{def:bc-protocol}). Each transaction of $\fin_i^t$ therefore
satisfies them in the context of the ancestry of its block, which is the
same in $\fin_i^t$ and in $\ch_i^r$. Among them is the anchor rule: an
anchor of a bundle equals the anchor, in its pool, of a block that precedes
the transaction's block (\emph{Ironwood Guide}, \S``Anchors''), and that
block is a block of $\fin_i^t$.
\end{proof}

Since bc-validity is a function of a block and of the blocks it commits to by
hash (Remark~\ref{rem:validity-recursion}), never of a chain that contains
it, the conclusion holds of every block of $\fin_i^t$ on every chain that
extends $\fin_i^t$. Snap-and-chat's finalised ledger is a sanitised sequence
of transactions, whose two readings coincide when double spends and missing
inputs are the only contextual invalidities and can differ in Zcash, where a
transaction invalid for a missing anchor can become valid at a later position
(Remark~\ref{rem:sac-deficiencies}, item~(3)).

\begin{remark}[Local finality and global consistency]%
\label{rem:local-not-global}
The finality status of a block is a function of one node's $\fin_i^t$, and
its value at one node is not evidence of global consistency. Two honest
nodes, or one node at two times, can hold finalised chains of different
lengths. The values of one honest node at two times agree in every execution
(Proposition~\ref{prop:af-single}). Under the premises of
Corollary~\ref{cor:af-either} the values of any two honest nodes agree, one a
prefix of the other (Definition~\ref{def:af}); when both premises fail, no
result of this volume excludes conflicting values at two distinct honest
nodes, and a node that observes a candidate conflicting with its finalised
chain records a finalisation safety hazard (Definition~\ref{def:fin}).
Shielded Labs' \emph{Zcash Crosslink Design Overview} (draft 10) states the same
of its design: honest nodes reach one chain through different histories, so
their values of $\fin$ are only prefix-compatible, not equal, and consensus
never reads $\fin$ (the overview, \S``10. Kinds of finality and which ledger
state is current''; Definition~\ref{def:sl-quantities}).
\end{remark}

\subsection{Finality and confirmation depth}\label{sec:k-depth}

\begin{remark}[Confirmation depth under Crosslink]\label{rem:not-depth}
The deployed facts are specified. No finite confirmation count is final: in
the model of the \emph{Consensus Guide}, for every attacker share of the hash
rate that is positive and below the honest share, the probability that the
attacker reverses a transaction is positive at every number of confirmations
(\S``Confirmation requirements'', Proposition ``No finite count is final'').
Consensus imposes no maximum reorganisation depth; a node may limit how far
it rolls back, and such a limit is node policy, under which different nodes
may stop on different sides of a long split, not global agreement
(\S``Rollback limits, coinbase maturity, and finality'', Remark ``Effect of a
rollback limit''). This section uses no rollback window. The
\emph{Consensus Guide} leaves to the layers above it the policy of when to
treat a transaction as final (\S``Layer boundary''). Under Crosslink~2,
$\fin_i^t$ is a separate state, computed by Construction~\ref{def:node-update}
from the BFT decisions that the best chain cites, while the best chain and its
confirmation counts remain those of the best-chain protocol. A block deep in
$\ch_i^t$ that is not a prefix of $\fin_i^t$ is not final, whatever its number
of confirmations, and a block of $\fin_i^t$ is final whatever its depth. A
displayed count of confirmations remains proof-of-work depth, as the overview
also states of its design (\S``10. Kinds of finality and which ledger state is
current'').
\end{remark}

% ledger_checks.py: 3. candidate below the truncated tip (47850 cases); with the carried headers on the tip's chain, lag = finality depth >= sigma + 1 and >= sigma + 2 confirmations (7330 cases); lag sigma + 1 blocks, sigma + 2 confirmations, attained, for every sigma <= 20; genesis context excluded (lag height(H) <= sigma for a low tip): ok
\begin{corollary}[Confirmations of a newly final block]%
\label{cor:final-confirmations}
Let an update of node $i$ set $\fin_i$ to $\candidate(H)$ for its new best
tip $H$, where $\LF(H) \neq \Obft$ and the headers carried by $\LF(H)$ are
ancestors of $H$. Every block that becomes final at $i$ at that update has
at least $\sigma + 2$ confirmations on the best chain $H$ at that time.

Status: proved (this volume), of a designed-but-unspecified construction.
\end{corollary}

\begin{proof}
A block $X$ that becomes final at the update satisfies $X \preceqbc
\candidate(H)$ (Definition~\ref{def:final}). Under the hypotheses,
Proposition~\ref{prop:lag} gives $\candidate(H) \preceqbc H\trunc{\sigma}
\preceqbc H$ and $\mathsf{height}(H) - \mathsf{height}(\candidate(H)) \ge
\sigma + 1$, so $X$ lies on $H$ at depth $\mathsf{height}(H) -
\mathsf{height}(X) \ge \sigma + 1$ (Definitions~\ref{def:chain}
and~\ref{def:trunc}). In the inclusive count of Remark~\ref{rem:depth-conf}, a
block at depth $d$ has $d + 1$ confirmations, so $X$ has at least $\sigma + 2$.
\end{proof}

The bound is attained by $\candidate(H)$ itself when the attainment
condition of Proposition~\ref{prop:lag} holds and $\candidate(H)$ was not
final at $i$ before the update. The corollary is the only derivation of the
count in the volume, and no argument rests on a value of $\sigma$
(Definition~\ref{def:params}).

\begin{remark}[Terminal states]\label{rem:terminal}
In the \emph{Wallet Guide}, a mined transaction is terminal only relative to
the wallet's rollback window, since consensus bounds no reorganisation depth
(\S``Status tracking and terminal states'', Definition ``Terminal states'';
\S``Reorganisation detection and recovery'', Remark ``The rollback
window''); the absence of a bound is specified, and the states are designed
but unspecified there. Finality would strengthen this status: by
Proposition~\ref{prop:final-guarantee}, a transaction of a block final at an
honest node stays final there whatever the depth of a later reorganisation
of its best chain, and under the premises of that proposition no honest
node's finalised chain conflicts with the block. No specification defines a
terminal state by finality. The \emph{Wallet Guide}'s states are taken on the
server's best chain, which under unchanged fork choice need not contain
$\fin_i^t$ (Remark~\ref{rem:fin-off-chain}).
\end{remark}

\subsection{Ledger views and stalls}\label{sec:k-views}

\begin{definition}[Application views]\label{def:views}
An \emph{application view} of node $i$ at time $t$ is one of the two chains
$\fin_i^t$ (Definition~\ref{def:fin}) and $(\ba_\mu)_i^t$
(Definition~\ref{def:ba}), for the node's bounded-available depth $\mu$
(Definition~\ref{def:mu}). Either view is to be exposed only once the node
is synced in the sense of Remark~\ref{rem:sync}, a recommendation and not a
consensus rule. The best chain $\ch_i^t$ is not an application view.
\end{definition}

The definition is the book's (chapter ``The Crosslink 2 Construction'',
sections ``Stalled Mode'', ``Locally bounded-available chain'' and
``Syncing and checkpoints''), and it is designed but unspecified. By
Proposition~\ref{prop:ledger-prefix}, at a node $i$ honest at $t$,
$\fin_i^t \preceqbc (\ba_\mu)_i^t$, so the second view extends the first. In
the book, for $\mu = \sigma$, the main factor in an application's choice
between them is not latency but its desired behaviour under a finalisation
stall: ``stall immediately, or keep processing user transactions until $L$
blocks have passed'' (section ``Stalled Mode'');
Remark~\ref{rem:stall-behaviour} states the second behaviour precisely. The
checkpoint and the freshness threshold of the syncing gate are open
(Remark~\ref{rem:sync}).

\begin{remark}[Application views in the Shielded Labs design]%
\label{rem:sl-views}
Shielded Labs' design has no Stalled Mode, no bound $L$ and no
bounded-available chain (\S\ref{sec:sl-rules}). An application there either
stops with $\fin_i^t$ during a finalisation stall or follows the best chain,
which under sticky fork choice contains $\fin_i^t$
(Proposition~\ref{prop:sticky-no-hazard}), with no bound on how far it runs
past $\fin_i^t$ (Proposition~\ref{prop:sl-unbounded}). The design offers an
application no ledger view between the best tip and $\fin_i^t$, although the
BFT final snapshot of Definition~\ref{def:sl-quantities} can lie on the best
chain above $\fin_i^t$; a displayed confirmation count is proof-of-work depth,
never finality (the overview, \S``10. Kinds of finality and which ledger state
is current''). These views are designed but unspecified.
\end{remark}

\begin{remark}[Behaviour during a stall]\label{rem:stall-behaviour}
In the book, the run of non-stalled blocks above the finalised chain is
bounded: under the hypotheses of Proposition~\ref{prop:bounded-run}, which
include $\sigma \le L$ and $\fin_i^t \preceqbc \ch_i^t$, every block of
$\ch_i^t$ at height greater than $\mathsf{height}(\fin_i^t) + L$ is stalled.
Under the proposed Zcash policy, a stalled block is coinbase-only
(Remark~\ref{rem:stalled-policy}), so no ordinary spend enters a block more
than $L$ above $\fin_i^t$, and the user transactions that $(\ba_\mu)_i^t$
carries above $\fin_i^t$ lie in blocks at most $L$ above it. The case in
which $\fin_i^t$ is not on the best chain is not covered
(Remark~\ref{rem:fin-off-chain}). In Shielded Labs' design, the number of
non-stalled blocks above $\fin$ on an honest node's best chain is unbounded
over executions (Proposition~\ref{prop:sl-unbounded}), and ordinary spends
keep landing above $\fin$ during a stall (the overview, \S``9. Sticky fork
choice''). Each behaviour is designed but unspecified in its design; which
behaviour a specification adopts is an open problem.
\end{remark}

\begin{remark}[Design goals]\label{rem:goals}
The book's design goals bind nothing: it calls them ideal goals, some of
which trade-offs may preclude (the book, chapter ``Design Goals''). Three
bear on clients: that users can rely on assured finality with an expected
time to finality below $30$ minutes; that the cost of attack for a one-hour
rollback is not reduced; and that the cost of attack to halt the chain
exceeds the $24$-hour revenue of proof-of-work mining rewards. With no value of
$\sigma$ or $L$ chosen (Definition~\ref{def:params}), the time goal cannot be
checked: under the hypotheses of Proposition~\ref{prop:lag}, $\fin_i$ trails the
tip by at least $\sigma + 1$ blocks at each update that sets $\fin_i$ to
$\candidate(H)$, the latency of $\Pibft$ adding to it, and no validated latency
bound exists (Remark~\ref{rem:latency}). A time goal fixes a block count only
through the block target spacing, which NU7 changes (\S\ref{sec:intro-scope}).
\end{remark}

\subsection{Confirmation policy, anchors and expiry}\label{sec:k-wallet}

\begin{remark}[Confirmation policy under finality]\label{rem:policy}
A confirmation policy is a triple $(\ctr, \cuntr, z)$ of integers $1 \le
\ctr \le \cuntr$, the trusted and the untrusted confirmation depth, and a
flag $z$ that permits zero-confirmation shielding (\emph{Wallet Guide},
\S``The confirmation policy'', Definition ``Confirmation policy''). ZIP 315
recommends $\ctr = 3$ and $\cuntr = 10$ (``Trusted and untrusted TXOs'') and
permits, without recommending, $z = \mathsf{true}$ (``Transparent UTXOs'');
it says that wallets SHOULD NOT permit the spending of outputs with fewer
than $3$ confirmations, with exceptions for certain transparent outputs,
among them those spent in shielding transactions (``Trusted and untrusted
TXOs''). The floor that ZIP 315 sets is therefore $3$ confirmations; the
bound $1 \le \ctr$ is a constraint of the definition. These rules are
specified. ZIP 315 does not mention finality, and no source maps the policy
onto the statuses of Proposition~\ref{prop:status-tri}: a wallet could spend
only final outputs, could combine a depth with the three statuses, or could
keep its depth policy and display finality separately. Under the hypotheses
of Corollary~\ref{cor:final-confirmations}, an output of a block that
becomes final has at least $\sigma + 2$ confirmations at that time. The
choice of a policy under finality is an open problem.
\end{remark}

\begin{remark}[Anchors under finality]\label{rem:anchors}
Transaction validity is unchanged by Crosslink~2
(Definition~\ref{def:bc-protocol}), so the anchor rule is inherited: the
anchor of a bundle must be the anchor, in its pool, of a block that precedes
the transaction's block in the chain, with no bound on how far back that
block lies (\emph{Ironwood Guide}, \S``Anchors''), and the rule is evaluated
against the ancestry of the block that contains the transaction. The recommended
anchor depth is a recommendation to the builder of a transaction, not a
consensus rule (\emph{Ironwood Guide}, \S``Anchors'', Remark ``Recommended
anchor depth''; \emph{Wallet Guide}, \S``Target and anchor heights''). These
rules are specified. No source restricts anchors to the finalised chain. The
anchors of the transactions of $\fin_i^t$ at an honest node are valid in their
own chains by Proposition~\ref{prop:fin-valid}.
\end{remark}

\begin{remark}[Anchor policy relative to $\fin$]\label{rem:anchor-fin}
Whether consensus should restrict anchors relative to finality, and whether a
wallet's default anchor should lie in $\fin_i^t$, is an open problem; the
trade-off below is this volume's analysis. A consensus rule can refer only to
objective data of a block, such as its final snapshot $\snapshot(\LF(H))$
(Definition~\ref{def:lf}), since $\fin_i^t$ is node state that consensus does
not read (Remark~\ref{rem:local-not-global}). An anchor at a block $X
\preceqbc \fin_i^t$ satisfies the anchor rule in every chain that extends $X$
and contains the transaction in a later block, so no reorganisation that
keeps $X$ invalidates the anchor, and $X$ stays final at the node by
Proposition~\ref{prop:final-guarantee}. Against this, at each update that
sets $\fin_i$ to $\candidate(H)$ under the hypotheses of
Proposition~\ref{prop:lag}, the distance from the tip $H$
to $\fin_i$ is at least $\sigma + 1$ blocks and equals $\finalitydepth(H)$,
which varies from block to block; an anchor in $\fin_i^t$ is therefore
deeper than the recommended anchor whenever $\sigma + 1$ exceeds the
recommended depth, and it has no fixed depth, unlike both conventions of the
\emph{Wallet Guide}, \S``Target and anchor heights'', Remark ``The anchor
convention''. A note mined above $X$ is not a leaf of the tree that the
anchor of $X$ roots (\emph{Ironwood Guide}, \S``Anchors'', Remark ``Old
anchors''), so it cannot be spent against that anchor; the candidate set of
the anchor is no larger than that of a later anchor; and an anchor at a depth
that other builders do not use distinguishes its builder (\emph{Ironwood
Guide}, \S``Anchors'', Remark ``Anchor choice as public data''). Since the
finalised chains of honest nodes can differ in length, wallets that anchor
at the finalised chains of different nodes can choose different anchors.
\end{remark}

\begin{remark}[Expiry under finality]\label{rem:expiry}
Transaction expiry is a consensus rule: a non-coinbase transaction with
expiry height $h_{\mathrm{exp}} \neq 0$ must not be mined at a height greater
than $h_{\mathrm{exp}}$, and $h_{\mathrm{exp}} = 0$ disables expiry
(\emph{Wallet Guide}, \S``Expiry (ZIP 203)'', Definition ``Transaction
expiry''). The default expiry delta at NU7 is that of the same section read
together with ZIP 218, ``Default expiry delta''; no argument of this volume
rests on its value. The rule is specified and unchanged by Crosslink~2.
Lemma~\ref{lem:expired-fin} below is the analogue, for the finalised chain, of
the \emph{Wallet Guide}'s Lemma ``Expired unmined is final on its chain''.
During Stalled Mode the best chain still grows, by stalled blocks, so an unmined
transaction still expires by height; under the coinbase-only policy of
Remark~\ref{rem:stalled-policy} no stalled block includes it, so a transaction
whose expiry height is reached while every new block is stalled expires unmined.
No source relates the expiry delta to $L$; the coupling is an open problem.
\end{remark}

\begin{lemma}[Expired unmined below $\fin$]\label{lem:expired-fin}
Let node $i$ be honest at $t$ and let a non-coinbase transaction have
$h_{\mathrm{exp}} \neq 0$. If $\mathsf{height}(\fin_i^t) \ge h_{\mathrm{exp}}$
and no block of $\fin_i^t$ contains the transaction, then no bc-valid chain
that extends $\fin_i^t$ contains it, and neither does $\fin_i^u$ for any time
$u \ge t$ at which $i$ is honest.

Status: proved (this volume), of a designed-but-unspecified construction.
\end{lemma}

\begin{proof}
The blocks of a chain that extends $\fin_i^t$ at heights at most
$h_{\mathrm{exp}}$ are those of $\fin_i^t$, and the expiry rule, a consensus
rule, excludes the transaction from every higher block of a bc-valid chain.
For $u \ge t$ at which $i$ is honest, Lemma~\ref{lem:fin-linear} gives
$\fin_i^t \preceqbc \fin_i^u$ and Proposition~\ref{prop:fin-valid} makes
$\fin_i^u$ bc-valid, so $\fin_i^u$ is such a chain.
\end{proof}

\begin{remark}[Downstream services]\label{rem:downstream}
For a service that credits deposits, three questions are open problems: the
meaning of a deposit received during a finalisation stall, when $\fin_i$
does not advance and, in the book, ordinary spends stop at a bounded
distance above it; the treatment of a transaction already credited whose
block falls into case~(iii) of Proposition~\ref{prop:status-tri} at the
service's node, and so can no longer become final there; and the protocol by
which a client that runs no node learns finality status. A service that
credits by finality status relies on the guarantee of
Proposition~\ref{prop:final-guarantee} instead of a delay of its own
choosing; no design states the obligations of such a service.
\end{remark}

\subsection{Light clients}\label{sec:k-light}

\begin{remark}[Commitment hook]\label{rem:light-hook}
Under ZIP 244, ``Block Header Changes'', the block commitment hash is a
linked list of hashes terminated by arbitrary data, and a later network
upgrade may replace its all-zero terminator by the hash of a new structure;
no upgrade through NU7 does (\emph{FlyClient Guide}, \S``The block commitment
hash'', Remark ``Extension of the block commitment hash''). This is the
specified place in the header at which a commitment to Crosslink data could
enter. No ZIP or specification section specifies a consumer of the
chain-history tree to which the header commits, a proof object or a verifier
(\emph{FlyClient Guide}, \S``The specified commitment and its consumers'').
These facts are specified.
\end{remark}

\begin{remark}[The light-client service]\label{rem:light-server}
The light-client service of the \emph{Wallet Guide} extracts its data,
compact blocks included, from the node's view of the best chain, and carries
no finality status (\S``Architecture and authority'', Definition
``Light-client architecture''). Its Assumption ``Honest-but-curious server''
concerns answers about the current best chain and has no clause for finality
status, and no light-client protocol for finality status is specified. A
light wallet that learns the finality status of a block from a server
therefore relies on the server for a status that the assumption does not
cover. These facts concern specified objects.
\end{remark}

\begin{remark}[Verification of finality by light clients]\label{rem:light-open}
A \emph{header-chain client} is a light client that verifies block headers
only, as the simplified-payment-verification client of the \emph{FlyClient
Guide}, \S``Header-only verification'', Definition ``Simplified payment
verification''. Verification of finality by such a client is an open
problem: no current design places in the header chain both the notarisation
signatures and a commitment to the set of finalisers against which they are
counted. In the book, each bc-header carries $\contextbft$, a commitment to a
bft-block (Definition~\ref{def:fields}), so a header-chain client sees the
cross-link reference; validity of the referenced block depends on its
notarisation proof, to which the header commits by hash through
$\contextbft$ but which a header-chain client does not hold, and on the
holders of the voting units of its BFT epoch (Definitions~\ref{def:units}
and~\ref{def:bft-block}), whose determination the book does not fix and to
which no header carries a commitment of its own. In Shielded Labs' design,
the roster for a BFT height is derived
from the bonds in the best-chain state at a snapshot and from the node's
slashing configuration (Definition~\ref{def:roster}); no header or certificate
carries a commitment to the roster itself, and a header-chain client, which
holds no transactions, cannot evaluate that derivation, so it cannot check the
stake weight behind a notarisation. The overview judges
its storage option ``commit only'', which uses the hook of
Remark~\ref{rem:light-hook}, probably insufficient (\S``3. The
proof-of-work side: the fat pointer and where it lives''), and lists the
choice of placement as open (Remark~\ref{rem:encoding}). Since a
certificate does not contain its notarisation signatures, which lie in the
fat pointer of a citing block, a copy of a certificate is not evidence of
canonicity (Definition~\ref{def:fat-pointer}), which bears on every light-client
design for that construction.
\end{remark}

\section{Formal models}\label{sec:formal}

This section states what each of the four formal models of Crosslink covers.
The status word ``machine-checked fragment'' is that of
Definition~\ref{def:classes}, which extends the \emph{FlyClient Guide}'s
Definition ``Epistemic classes'' (\emph{FlyClient Guide}, \S``Scope and
status''); it records that a tool has checked a stated fragment, and claims
nothing beyond that fragment. Two models are written in Quint, a specification
language whose specifications translate into Lamport's specification language
TLA+, so that the model checkers TLC, which enumerates the states of a finite
instance, and Apalache, which checks executions up to a bounded number of
steps symbolically, apply to them. The third is written for Lean, a proof
assistant that checks proofs of stated theorems. The fourth is written in TLA+
directly.

\subsection{The Lean formalisation}\label{sec:mc-lean}

\begin{remark}[Lean formalisation]\label{rem:mc-lean}
D.-E.~Hopwood, ``Add work-in-progress Lean formalization'', an unmerged pull
request to the TFL book (last commit 9 March 2026), states in Lean
Theorem~\ref{thm:pa-af}, with an honesty predicate that need not hold at
earlier times, and a form of Theorem~\ref{thm:fa-af} whose last-final
function is a placeholder, the identity, so that its premise is Agreement on
the bft-chains and nothing ties those chains to the bc-chains. Both proofs
are left unfinished, and both statements hold vacuously, because the
parameter structure over which they are stated includes the hypothesis on
tips stated below, which nothing satisfies. Lean accepts a proof left
unfinished with a warning, so a successful build does not show that every
stated theorem is proved. Machine-checked, per the pull request's continuous
integration only and not rebuilt for this volume, is the bc-chain case of
Lemma~\ref{lem:linear-prefix}. The lemmas on the last common ancestor of
bc-chains, that $\lca(A, B)$ is a prefix of both chains and that every
common prefix is a prefix of it (Definition~\ref{def:lca}), and
$\candidate(H) \preceqbc H\trunc{\sigma}$, the first clause of
Proposition~\ref{prop:lag}, are checked under the hypothesis that two
bc-chains with one tip are equal. The formalisation represents a bc-chain as
a list of block hashes, and two different lists can share a tip, so that
hypothesis fails and these lemmas hold vacuously. The last is moreover
stated for an arbitrary $\LF$, whose definition is left unfinished.
Lemma~\ref{lem:fin-depth} and both theorems are unproved in the
formalisation. Its parameter structure encodes $L \ge 2\sigma$, which the
book gives as advice for practice, as a hard constraint.

Status: machine-checked fragment.
\end{remark}

\subsection{The Informal Systems Quint model}\label{sec:mc-isq}

\begin{remark}[Informal Systems model]\label{rem:mc-isq}
Informal Systems' Quint model of Crosslink (23 July 2025) is cited by
author, kind and date, its title naming a function of the model. It models
the proposal-validity check of $\Pibft$ alone, for one node, with no
adversary, voting, signatures or network. Its invariant states that the
snapshots of consecutive accepted bft-blocks are ancestor-linked; its
description reports that random simulation finds no violation and that TLC
checks the invariant on an instance whose bc-block and bft-block stores hold
up to eight blocks each, at confirmation depth one; neither run was repeated
for this volume. The invariant follows in one step from the model's own
linearity guard, since its store of bc-blocks only grows. That guard
forbids the equal snapshot that the Linearity rule of
Definition~\ref{def:bft-validity} permits and that the book's argument for
Proposition~\ref{prop:bft-live} needs, and the model's tail is an unordered
set of at least $\sigma$ blocks rather than a sequence of exactly $\sigma$
headers (Definition~\ref{def:fields}).

Status: machine-checked fragment.
\end{remark}

\subsection{The Manian Quint suite}\label{sec:mc-zq}

\begin{remark}[Manian suite]\label{rem:mc-zq}
Z.~Manian's Quint suite for Crosslink (18 May to 4 June 2026) is cited by
author, kind and date, its only title being the name of its repository. Per
its continuous integration it runs bounded Apalache checks and TLC checks, on
small instances and at bounded depth, of models in which every finaliser has
equal voting power: a round machine shaped after Tendermint with a fixed
roster, and separate models whose roster changes across BFT heights. ``The
induction itself is not closed by Apalache''. It adds witness runs on hand-built
fixtures. Its models aligned with the book are only witness-tested. The
best-chain rules of Definition~\ref{def:bc-validity} and Stalled Mode
(Definition~\ref{def:stalled}) are absent from its core models, so the suite
is not a check of the construction.

Status: machine-checked fragment.
\end{remark}

\begin{remark}[Witnessed findings]\label{rem:mc-zq-findings}
The suite witnesses the following on hand-built fixtures; the witnesses are
instances, not general proofs.
\begin{enumerate}[label=(\arabic*)]
\item Under a partition, with more than one third of the BFT protocol's
  voting power subverted, a fork choice that rejects every chain
  conflicting with the node's finalised chain, whatever its work, as sticky
  fork choice does (Definition~\ref{def:sticky}), turns a transient failure
  of Final Agreement into a split of the best chain that persists to the end
  of the witness run, after the partition heals and after a further block on
  the heavier branch; permanence is the model's argument that the rule
  ignores work, not a checked property. The unchanged fork choice of
  Definition~\ref{def:fork-choice} reconverges. The model records no hazard;
  under Definition~\ref{def:fin} the node whose finalised block lies on the
  abandoned branch would record a finalisation safety hazard. When only one
  side finalises legitimately the comparison reverses: the rejecting fork choice
  keeps the finalised branch, and the unchanged one may abandon it.
\item Verification of each certificate against the roster in force at its
  BFT height is well-founded, and dropping that check allows a takeover of
  the roster in one step.
\item Leaked keys of former roster members forge an early certificate for a
  node that syncs without a checkpoint, and a checkpoint confines the
  forgery (Remark~\ref{rem:sync}).
\item Results for rosters of unit voting power transfer to stake-weighted
  thresholds only if the fault bound is read over stake: an instance of four
  finalisers whose one faulty finaliser, within the fault bound by count,
  holds more than one third of the stake yields conflicting certificates.
\end{enumerate}

Status: machine-checked fragment.
\end{remark}

\subsection{The Garcia TLA+ specification}\label{sec:mc-tla}

\begin{remark}[Garcia specification]\label{rem:mc-tla}
A.~Garcia's ``Crosslink2 TLA+ Specification'' (18 April 2025 to 7 April
2026), which Informal Systems' model cites as a proposal for verification
work, describes itself as a proof of concept. Its finalised chain is the
longest bc-chain of the model, truncated, computed with no bft-chain, no
context and no update rule $\updatefin$, and its last-final function is the
identity on a node's bft-chain. It submits invariants to TLC, one of them
named Assured Finality but stated over that chain, and reports no completed
run; it is therefore no machine-checked fragment.
\end{remark}

\subsection{Coverage of the formal models}\label{sec:mc-coverage}

% safety_checks.py: 1. Final Agreement and Linearity: all fin values agree (1500 trials, 4500 node runs, 1949 with forks deeper than sigma): ok
% ledger_checks.py: 5. sigma <= L: every block more than L above fin is stalled (3000 trials, 270980 checks of 44776 distinct blocks): ok
\begin{remark}[Scope of machine checking]\label{rem:mc-coverage}
No machine-checked artefact covers the book's Assured Finality
(Definition~\ref{def:af}) or either of its theorems, nor the ebb-and-flow
paper's properties P1 and P2 and its prefix condition
(Definition~\ref{def:ntt-ebbflow}). Four formal models exist: the Lean
formalisation (Remark~\ref{rem:mc-lean}), Informal Systems' model
(Remark~\ref{rem:mc-isq}), Manian's suite (Remark~\ref{rem:mc-zq}) and
Garcia's specification (Remark~\ref{rem:mc-tla}). The
proofs of this volume are on paper; the scripts whose output the text quotes
check their statements on random instances, by exhaustive enumeration over
small parameter ranges, and on constructed counterexamples.
\end{remark}

\section{Status and open problems}\label{sec:open}

This section states the components that a specification of Crosslink~2 for
Zcash must fix and the one open problem that no earlier section states,
activation on an existing chain. It then collects the status of the volume in
three tables: every open problem, with the environment that states it and its
source (Table~\ref{tab:open}); every goal of the construction, with the result
that discharges it and the open problems on which its standing for Zcash
depends (Table~\ref{tab:goal-status}); and every rule, with the party that
checks or follows it and the results that use it (Table~\ref{tab:rules}). The
tables cite environments by number and restate no statement.

\subsection{Components a specification must fix}\label{sec:open-spec}

\begin{remark}[Components left unfixed]\label{rem:spec-ledger}
A specification of Crosslink~2 for Zcash must fix each of the following
components; no source fixes any of them.
\begin{enumerate}[label=(\arabic*)]
\item The concrete BFT protocol $\Piorigbft$, the assignment of voting units
  to its BFT epochs and the format of its notarisation proofs
  (Definitions~\ref{def:bft-protocol} and~\ref{def:units}). The book uses
  adapted Streamlet as an example (Example~\ref{ex:streamlet}); the Shielded
  Labs design uses Tenderlink (\S\ref{sec:sl-terms}), which has no
  specification.
\item The values of $\sigma$ and $L$ (Definition~\ref{def:params}), with
  $\sigma + 1 \le L$: $\sigma \le L$ is needed, for some best-chain protocol
  of the interface, for the bounded non-stalled run
  (Proposition~\ref{prop:short-L}), and $L \ge \sigma + 1$ for Stalled Mode
  to end while $\Pibft$ is live (Proposition~\ref{prop:stall-short-L}); the
  depth $\mu$ is chosen by each node, by default $\sigma$
  (Definition~\ref{def:mu}). No source computes an adequate value of either
  (Remarks~\ref{rem:pow-prefix} and~\ref{rem:choose-L}).
\item The stalled-block predicate $\isstalledblock$
  (Definition~\ref{def:params}); the book proposes coinbase-only blocks
  (Remark~\ref{rem:stalled-policy}).
\item The encoding of the cross-link fields and their placement in the block
  header or in the block (Definition~\ref{def:fields},
  Remark~\ref{rem:encoding}).
\item The mapping from stake to voting units and the selection of the roster
  (Definitions~\ref{def:voting-power}, \ref{def:active-roster}
  and~\ref{def:roster}).
\item The staking calendar (Definition~\ref{def:staking-window},
  Remark~\ref{rem:calendar}).
\item The format of bond transactions, the range of quantised bond amounts
  and a minimum stake (Definition~\ref{def:bond}, Remarks~\ref{rem:bond-tx}
  and~\ref{rem:opt-in}).
\item The reward rule (Remark~\ref{rem:payment-open}), and the division of
  the commission among finalisers and its denominator
  (Definition~\ref{def:commission}).
\item The fees of staking transactions and the number of staking actions that
  one transaction may carry (Remark~\ref{rem:bond-tx}).
\item The rotation of finaliser keys and the initial roster
  (Remark~\ref{rem:key-rotation}, Definition~\ref{def:bootstrap}).
\item Activation on an existing chain (Remark~\ref{rem:activation}).
\end{enumerate}
No formal model checks the safety and liveness statements of the
construction; the existing models check fragments only
(Remark~\ref{rem:mc-coverage}).

Status: open problem.
\end{remark}

\begin{remark}[Activation]\label{rem:activation}
No design activates Crosslink~2 on the existing Zcash chain with a safety
argument that rests on consensus rules alone. The book models both chains
from genesis: every $x$-chain starts at $\Origin_x$
(Definition~\ref{def:chain}), and the Genesis bc-block rule sets
$\Obc.\contextbft = \Obft$ (Definition~\ref{def:bc-validity}). Shielded Labs'
design starts the BFT protocol from best-chain state by a deterministic
bootstrap (Definition~\ref{def:bootstrap}), whose roster height $h_1$ and
start height $h_2$ are unsettled, the two current texts of the overview
differing. The overview's safety argument (Remark~\ref{rem:bootstrap-rollback})
rests on rollback limits, which are node policy: consensus imposes no maximum
reorganisation depth (\emph{Consensus Guide}, \S``Rollback limits, coinbase
maturity, and finality''), and a node-local limit does not give its conclusion.
Proposition~\ref{prop:bootstrap-safety} proves the conclusion under Prefix
Agreement of the best chain at depth $h_2 - h_1$ instead, a prefix property
that is an assumption for Zcash (Remark~\ref{rem:pow-prefix}). The safety of
the bootstrap therefore rests on a parameter that consensus does not fix or
on that assumption.

Status: open problem.
\end{remark}

% [sec:open-spec] no worked numbers: the subsection cites no compute/ output.
\begin{table}[p]
\centering
\footnotesize \renewcommand{\arraystretch}{1.1}
\begin{tabular}{@{}p{0.42\linewidth}p{0.26\linewidth}p{0.25\linewidth}@{}}
\toprule
Open problem & Stated in & Source \\
\midrule
The open question of the paper
  & Remark~\ref{rem:ntt-open}, \S\ref{sec:model-pow} & paper \\
Status of the prefix properties for Zcash
  & Remark~\ref{rem:pow-prefix}, \S\ref{sec:iface-bc}
  & \emph{FlyClient Guide}; \emph{Consensus Guide}; ZIP 218 \\
Honestly submitted blocks
  & Remark~\ref{rem:honestly-submitted}, \S\ref{sec:iface-bft} & book \\
Parameters & \S\ref{sec:cons-params} & book \\
Encoding & Remark~\ref{rem:encoding}, \S\ref{sec:cons-struct} & overview \\
Proposed Zcash policy
  & Remark~\ref{rem:stalled-policy}, \S\ref{sec:cons-stalled} & book \\
Candidate or final snapshot
  & Remark~\ref{rem:candidate-choice}, \S\ref{sec:led-fin} & book \\
Syncing and checkpoints & \S\ref{sec:led-sync} & book \\
Necessity of the other rules
  & Remark~\ref{rem:rule-necessity}, \S\ref{sec:saf-either} & book \\
Missing results for $\ba$
  & Remark~\ref{rem:ba-missing}, \S\ref{sec:saf-ba} & book \\
Transfer from $\Piorigbc$
  & Remark~\ref{rem:origbc-transfer}, \S\ref{sec:live-sub} & book \\
Advance of finalisation
  & Remark~\ref{rem:fin-advances}, \S\ref{sec:live-fin} & book \\
Choice of $L$
  & Remark~\ref{rem:choose-L}, \S\ref{sec:live-fin} & book \\
Necessity of $\sigma \le L$ for Zcash & \S\ref{sec:bd-run} & this volume \\
Finalised chain off the best chain
  & Remark~\ref{rem:fin-off-chain}, \S\ref{sec:bd-run} & book \\
Open parts of the override design
  & Remark~\ref{rem:override-open}, \S\ref{sec:bd-override} & book \\
Deferred validity
  & Remark~\ref{rem:sl-deferred}, \S\ref{sec:sl-certs} & overview \\
Status of sticky fork choice
  & Remark~\ref{rem:sticky-open}, \S\ref{sec:sl-forkchoice} & overview \\
Final Agreement of Tenderlink
  & Remark~\ref{rem:tenderlink-fa}, \S\ref{sec:sl-theorems} & this volume \\
\bottomrule
\end{tabular}
\caption{Open problems of the volume, in the order of their statements, each
with the environment or subsection that states it and the source in which it
arises. The book, the paper and the overview are the sources so named in
``Sources and conventions'' (\S\ref{sec:intro-sources}); the stating
environment cites the chapter or section. ``This volume'' marks a problem
that no source states.}
\label{tab:open}
\end{table}

\begin{table}[p]
\ContinuedFloat
\centering
\footnotesize \renewcommand{\arraystretch}{1.1}
\begin{tabular}{@{}p{0.42\linewidth}p{0.26\linewidth}p{0.25\linewidth}@{}}
\toprule
Open problem & Stated in & Source \\
\midrule
Finaliser consent
  & Remark~\ref{rem:opt-in}, \S\ref{sec:fs-roster} & this volume \\
Active roster
  & Definition~\ref{def:active-roster}, \S\ref{sec:fs-roster} & overview \\
Bootstrap & Definition~\ref{def:bootstrap}, \S\ref{sec:fs-roster}
  & overview \\
Transaction contents
  & Remark~\ref{rem:bond-tx}, \S\ref{sec:fs-bonds} & \emph{Rewards Distribution}
  \\
Open bond questions
  & Remark~\ref{rem:bond-open}, \S\ref{sec:fs-bonds} & this volume; overview \\
Staking window
  & Definition~\ref{def:staking-window}, \S\ref{sec:fs-windows}
  & overview \\
Calendar divergence
  & Remark~\ref{rem:calendar}, \S\ref{sec:fs-windows}
  & Shielded Labs texts \\
Privacy design
  & Remark~\ref{rem:staking-privacy}, \S\ref{sec:fs-windows} & this volume \\
Open privacy and expiry questions
  & Remark~\ref{rem:staking-open}, \S\ref{sec:fs-windows}
  & this volume; \emph{Wallet Guide} \\
Unsettled payment rule
  & Remark~\ref{rem:payment-open}, \S\ref{sec:fs-rewards}
  & book; overview; \emph{Rewards Distribution} \\
Commission and bond earnings
  & Definition~\ref{def:commission}, \S\ref{sec:fs-rewards}
  & overview; \emph{Rewards Distribution} \\
Conditional rewards
  & Remark~\ref{rem:conditional-rewards}, \S\ref{sec:fs-rewards}
  & overview; \emph{Rewards Distribution} \\
Withheld rewards and the NSM
  & Remark~\ref{rem:nsm}, \S\ref{sec:fs-rewards} & this volume \\
Slash fork & Definition~\ref{def:slash-fork}, \S\ref{sec:fs-slashing}
  & overview \\
Open slashing details
  & Remark~\ref{rem:slash-open}, \S\ref{sec:fs-slashing} & overview \\
Keys and signatures & Remark~\ref{rem:keys}, \S\ref{sec:fs-keys} & overview \\
Key rotation & Remark~\ref{rem:key-rotation}, \S\ref{sec:fs-keys} & book \\
Behaviour during a stall
  & Remark~\ref{rem:stall-behaviour}, \S\ref{sec:k-views} & this volume \\
Confirmation policy under finality
  & Remark~\ref{rem:policy}, \S\ref{sec:k-wallet}
  & \emph{Wallet Guide}; ZIP 315 \\
Anchor policy relative to $\fin$
  & Remark~\ref{rem:anchor-fin}, \S\ref{sec:k-wallet}
  & \emph{Wallet Guide}; \emph{Ironwood Guide} \\
Expiry under finality
  & Remark~\ref{rem:expiry}, \S\ref{sec:k-wallet} & this volume \\
Downstream services
  & Remark~\ref{rem:downstream}, \S\ref{sec:k-wallet} & this volume \\
Verification of finality by light clients
  & Remark~\ref{rem:light-open}, \S\ref{sec:k-light}
  & overview; \emph{FlyClient Guide} \\
Components left unfixed
  & Remark~\ref{rem:spec-ledger}, \S\ref{sec:open-spec} & this volume \\
Activation & Remark~\ref{rem:activation}, \S\ref{sec:open-spec}
  & book; overview \\
\bottomrule
\end{tabular}
\caption[]{Open problems, continued.}
\end{table}

\subsection{Status by goal and by rule}\label{sec:open-status}

Table~\ref{tab:goal-status} gives, for each goal of the construction, the
result that discharges it, the status of that result, and the open problems
on which its standing for Zcash depends. Table~\ref{tab:rules} lists every
validity rule and honest-behaviour rule of the two sub-protocols and the
update rule of the node state, with the party that checks or follows it and
the results whose statement or proof uses it; it extends
Remark~\ref{rem:bc-rules-role} from the best-chain rules to the whole
mechanism. A rule that no proved result uses is marked so.

% [sec:open-status] no worked numbers: the subsection cites no compute/ output.
\begin{table}[tbp]
\centering
\footnotesize \renewcommand{\arraystretch}{1.15}
\begin{tabular}{@{}p{0.17\linewidth}p{0.29\linewidth}p{0.23\linewidth}
  p{0.21\linewidth}@{}}
\toprule
Goal & Result & Status & Depends on \\
\midrule
Assured Finality
  & Corollary~\ref{cor:af-either}, Assured Finality under either premise;
  Final Agreement of $\Pibft$ from that of $\Piorigbft$ rests on
  Proposition~\ref{prop:erasure}, Transfer by erasure, argued (TFL book)
  & proved (this volume), of a designed-but-unspecified construction
  & Remarks~\ref{rem:pow-prefix} and~\ref{rem:origbc-transfer} (Prefix
  Agreement of $\Pibc$ for Zcash); Remark~\ref{rem:tenderlink-fa} (a
  concrete BFT protocol) \\
\addlinespace
Safety under subversion of the BFT protocol
  & Corollary~\ref{cor:bft-subversion}, Safety under subversion of the BFT
  protocol
  & proved (this volume), of a designed-but-unspecified construction
  & Remarks~\ref{rem:pow-prefix} and~\ref{rem:origbc-transfer} \\
\addlinespace
Ledger prefix property
  & Proposition~\ref{prop:ledger-prefix}, Ledger prefix property
  & proved (TFL book), of a designed-but-unspecified construction
  & none \\
\addlinespace
Bounded availability
  & Propositions~\ref{prop:bounded-run}, Bounded non-stalled run;
  \ref{prop:hazard-free}, Bounded hazard-freeness
  & proved (this volume), of a designed-but-unspecified construction
  & Remarks~\ref{rem:fin-off-chain} and~\ref{rem:choose-L} \\
\addlinespace
Liveness
  & Propositions~\ref{prop:bc-live}, Liveness of $\Pibc$;
  \ref{prop:bft-live}, Liveness of $\Pibft$; \ref{prop:stall-only}, Stalled
  Mode only on failure of the liveness assumptions
  & argued (TFL book)
  & Remarks~\ref{rem:fin-advances}, \ref{rem:origbc-transfer}
  and~\ref{rem:choose-L} \\
\addlinespace
Agreement on $\ba$
  & Corollary~\ref{cor:ba-agree}, Agreement on $\ba$ from Prefix Agreement
  & proved (this volume), of a designed-but-unspecified construction
  & Remarks~\ref{rem:pow-prefix} and~\ref{rem:origbc-transfer} (Prefix
  Agreement of $\Pibc$ at $\mu$ for Zcash); Remark~\ref{rem:ba-missing} \\
\addlinespace
Client contract
  & Propositions~\ref{prop:final-guarantee}, Guarantee of finality;
  \ref{prop:status-tri}, Finality status; \ref{prop:fin-valid}, Validity of
  the finalised chain
  & proved (this volume), of a designed-but-unspecified construction
  & Remarks~\ref{rem:light-open} and~\ref{rem:policy}; for the second
  clause of Proposition~\ref{prop:final-guarantee}, those of Assured
  Finality \\
\bottomrule
\end{tabular}
\caption{Status by goal: the result that discharges each goal, its status
in the words of Definition~\ref{def:classes}, and the open problems on which
its standing for Zcash depends. The liveness results are prose arguments of
the book.}
\label{tab:goal-status}
\end{table}

% [sec:open-status] no worked numbers: the table cites no compute/ output.
\begin{table}[p]
\centering
\footnotesize \renewcommand{\arraystretch}{1.1}
\begin{tabular}{@{}p{0.33\linewidth}p{0.32\linewidth}p{0.27\linewidth}@{}}
\toprule
Rule & Defined in & Checked or followed by \\
\midrule
Genesis bft-block rule
  & Definition~\ref{def:bft-validity} & every node \\
\multicolumn{3}{@{\quad}p{0.95\linewidth}@{}}{
  Used by: Definition~\ref{def:snapshot}, Snapshot (genesis case);
  Proposition~\ref{prop:context-exists}, Existence of a valid context;
  Proposition~\ref{prop:forced-entry}, Forced entry into Stalled Mode;
  Proposition~\ref{prop:fa-needs-lin}, Necessity of the Linearity rule;
  Proposition~\ref{prop:stall-short-L}, Stall when $L \le \sigma$;
  Proposition~\ref{prop:short-L}, Necessity of the constraint $\sigma \le L$;
  Proposition~\ref{prop:sl-fa}, Final Agreement route under the Shielded Labs
  rules; Proposition~\ref{prop:path-final}, Path to finality.} \\
\addlinespace
Inherited origbft rules, with the Parent rule
  & Definition~\ref{def:bft-validity} & every node \\
\multicolumn{3}{@{\quad}p{0.95\linewidth}@{}}{
  Used by: Lemma~\ref{lem:uniqueness}, Per-BFT-epoch uniqueness;
  Proposition~\ref{prop:context-exists}, Existence of a valid context;
  Lemma~\ref{lem:fd-step}, Growth of finality depth;
  Proposition~\ref{prop:exit-branch}, Branch dependence of exit;
  Proposition~\ref{prop:forced-entry}, Forced entry into Stalled Mode;
  Theorem~\ref{thm:fa-af}, Final Agreement implies Assured Finality;
  Proposition~\ref{prop:bounded-run}, Bounded non-stalled run;
  Proposition~\ref{prop:coalition}, Power of a finaliser coalition;
  Proposition~\ref{prop:path-final}, Path to finality.} \\
\addlinespace
Linearity rule
  & Definition~\ref{def:bft-validity} & every node \\
\multicolumn{3}{@{\quad}p{0.95\linewidth}@{}}{
  Used by: Proposition~\ref{prop:context-exists}, Existence of a valid
  context; Lemma~\ref{lem:fd-step}, Growth of finality depth;
  Proposition~\ref{prop:exit-branch}, Branch dependence of exit;
  Theorem~\ref{thm:fa-af}, Final Agreement implies Assured Finality;
  Proposition~\ref{prop:fa-needs-lin}, Necessity of the Linearity rule;
  Proposition~\ref{prop:bft-live}, Liveness of $\Pibft$;
  Proposition~\ref{prop:bounded-run}, Bounded non-stalled run;
  Proposition~\ref{prop:short-L}, Necessity of the constraint $\sigma \le L$;
  Proposition~\ref{prop:sl-fa}, Final Agreement route under the Shielded Labs
  rules; Proposition~\ref{prop:path-final}, Path to finality.} \\
\addlinespace
Tail Confirmation rule
  & Definition~\ref{def:bft-validity}
  & every node; a validator also downloads and checks the blocks before
  voting \\
\multicolumn{3}{@{\quad}p{0.95\linewidth}@{}}{
  Used by: Remark~\ref{rem:struct-purpose}, Purpose of the additions, states
  its purpose; Proposition~\ref{prop:context-exists}, Existence of a valid
  context; Proposition~\ref{prop:lag}, Finalisation lag;
  Proposition~\ref{prop:fa-needs-lin}, Necessity of the Linearity rule;
  Proposition~\ref{prop:bft-live}, Liveness of $\Pibft$;
  Proposition~\ref{prop:stall-short-L}, Stall when $L \le \sigma$;
  Proposition~\ref{prop:short-L}, Necessity of the constraint
  $\sigma \le L$.} \\
\addlinespace
Honest proposal
  & Construction~\ref{def:honest-proposal} & proposer \\
\multicolumn{3}{@{\quad}p{0.95\linewidth}@{}}{Used by:
  Proposition~\ref{prop:bft-live}, Liveness of $\Pibft$;
  Proposition~\ref{prop:lag}, Finalisation lag (attainment of the
  bound).} \\
\addlinespace
Honest voting: Valid proposal criterion, Confirmed best-chain criterion
  & Construction~\ref{def:honest-voting} & validator \\
\multicolumn{3}{@{\quad}p{0.95\linewidth}@{}}{Used by:
  Definition~\ref{def:nonhonest}, Non-honest casting;
  Proposition~\ref{prop:fa-needs-lin}, Necessity of the Linearity rule;
  Proposition~\ref{prop:stall-short-L}, Stall when $L \le \sigma$;
  Remark~\ref{rem:only-if}, Refusal to vote;
  Remark~\ref{rem:override}, Override argument.} \\
\bottomrule
\end{tabular}
\caption{Rules: every validity rule and honest-behaviour rule of $\Pibft$
and $\Pibc$ and the update rule of the node state, with its defining
environment, the party that checks or follows it, and the environments
whose statement or proof uses it.}
\label{tab:rules}
\end{table}

\begin{table}[p]
\ContinuedFloat
\centering
\footnotesize \renewcommand{\arraystretch}{1.1}
\begin{tabular}{@{}p{0.33\linewidth}p{0.32\linewidth}p{0.27\linewidth}@{}}
\toprule
Rule & Defined in & Checked or followed by \\
\midrule
Genesis bc-block rule
  & Definition~\ref{def:bc-validity} & every node \\
\multicolumn{3}{@{\quad}p{0.95\linewidth}@{}}{
  Used by: Definition~\ref{def:lf}, Last final context ($\LF(\Obc) = \Obft$);
  Proposition~\ref{prop:context-exists}, Existence of a valid context;
  Theorem~\ref{thm:fa-af}, Final Agreement implies Assured Finality;
  Proposition~\ref{prop:bc-live}, Liveness of $\Pibc$;
  Proposition~\ref{prop:hazard-free}, Bounded hazard-freeness;
  Proposition~\ref{prop:bounded-run}, Bounded non-stalled run.} \\
\addlinespace
Inherited origbc rules
  & Definition~\ref{def:bc-validity} & every node \\
\multicolumn{3}{@{\quad}p{0.95\linewidth}@{}}{
  Used by: Lemma~\ref{lem:lf-monotone}, Prefix Lemma;
  Proposition~\ref{prop:context-exists}, Existence of a valid context;
  Proposition~\ref{prop:sl-unbounded}, Unbounded run under the Shielded Labs
  rules; Proposition~\ref{prop:fin-valid}, Validity of the finalised chain.} \\
\addlinespace
Valid context rule
  & Definition~\ref{def:bc-validity} & every node \\
\multicolumn{3}{@{\quad}p{0.95\linewidth}@{}}{
  Used by: Proposition~\ref{prop:context-exists}, Existence of a valid
  context; Lemma~\ref{lem:fd-step}, Growth of finality depth;
  Proposition~\ref{prop:exit-branch}, Branch dependence of exit;
  Lemma~\ref{lem:fin-depth}, Local fin-depth; Proposition~\ref{prop:lag},
  Finalisation lag; Theorem~\ref{thm:pa-af}, Prefix Agreement implies Assured
  Finality; Corollary~\ref{cor:bft-subversion}, Safety under subversion of
  the BFT protocol; Theorem~\ref{thm:fa-af}, Final Agreement implies Assured
  Finality; Proposition~\ref{prop:fa-needs-lin}, Necessity of the Linearity
  rule; Proposition~\ref{prop:bc-live}, Liveness of $\Pibc$;
  Proposition~\ref{prop:stall-short-L}, Stall when $L \le \sigma$;
  Proposition~\ref{prop:short-L}, Necessity of the constraint $\sigma \le L$;
  Proposition~\ref{prop:sl-fa}, Final
  Agreement route under the Shielded Labs rules;
  Proposition~\ref{prop:sl-unbounded}, Unbounded run under the Shielded Labs
  rules; Proposition~\ref{prop:coalition}, Power of a finaliser coalition;
  Proposition~\ref{prop:path-final}, Path to finality.} \\
\addlinespace
Extension rule
  & Definition~\ref{def:bc-validity} & every node \\
\multicolumn{3}{@{\quad}p{0.95\linewidth}@{}}{
  Used by: Lemma~\ref{lem:lf-monotone}, Prefix Lemma;
  Proposition~\ref{prop:context-exists}, Existence of a valid context;
  Lemma~\ref{lem:fd-step}, Growth of finality depth;
  Proposition~\ref{prop:fa-needs-lin}, Necessity of the Linearity rule;
  Proposition~\ref{prop:bc-live}, Liveness of $\Pibc$;
  Proposition~\ref{prop:bounded-run}, Bounded non-stalled run;
  Proposition~\ref{prop:short-L}, Necessity of the constraint $\sigma \le L$;
  Proposition~\ref{prop:sl-unbounded}, Unbounded run under the Shielded Labs
  rules.} \\
\addlinespace
Last Final Snapshot rule
  & Definition~\ref{def:bc-validity} & every node \\
\multicolumn{3}{@{\quad}p{0.95\linewidth}@{}}{
  Used by: Lemma~\ref{lem:candidate-lfs}, Candidate under the Last Final
  Snapshot rule; Proposition~\ref{prop:context-exists}, Existence of a valid
  context; Proposition~\ref{prop:exit-branch}, Branch dependence of exit;
  Proposition~\ref{prop:forced-entry}, Forced entry into Stalled Mode;
  Proposition~\ref{prop:fa-needs-lin}, Necessity of the Linearity rule;
  Proposition~\ref{prop:bc-live}, Liveness of $\Pibc$;
  Proposition~\ref{prop:hazard-free}, Bounded hazard-freeness;
  Proposition~\ref{prop:bounded-run}, Bounded non-stalled run;
  Proposition~\ref{prop:short-L}, Necessity of the constraint $\sigma \le L$;
  Proposition~\ref{prop:sl-unbounded}, Unbounded run under the Shielded Labs
  rules.} \\
\addlinespace
Finality depth rule
  & Definition~\ref{def:bc-validity} & every node \\
\multicolumn{3}{@{\quad}p{0.95\linewidth}@{}}{
  Used by: Proposition~\ref{prop:context-exists}, Existence of a valid
  context; Proposition~\ref{prop:exit-branch}, Branch dependence of exit;
  Proposition~\ref{prop:forced-entry}, Forced entry into Stalled Mode;
  Proposition~\ref{prop:fa-needs-lin}, Necessity of the Linearity rule;
  Proposition~\ref{prop:bc-live}, Liveness of $\Pibc$;
  Proposition~\ref{prop:hazard-free}, Bounded hazard-freeness;
  Proposition~\ref{prop:bounded-run}, Bounded non-stalled run;
  Proposition~\ref{prop:short-L}, Necessity of the constraint $\sigma \le L$;
  Proposition~\ref{prop:sl-unbounded}, Unbounded run under the Shielded Labs
  rules (by its absence).} \\
\bottomrule
\end{tabular}
\caption[]{Rules, continued.}
\end{table}

\begin{table}[p]
\ContinuedFloat
\centering
\footnotesize \renewcommand{\arraystretch}{1.1}
\begin{tabular}{@{}p{0.33\linewidth}p{0.32\linewidth}p{0.27\linewidth}@{}}
\toprule
Rule & Defined in & Checked or followed by \\
\midrule
Honest block production (choice of context)
  & Construction~\ref{def:honest-production} & block producer \\
\multicolumn{3}{@{\quad}p{0.95\linewidth}@{}}{Used by:
  Proposition~\ref{prop:context-exists}, Existence of a valid context;
  Proposition~\ref{prop:bc-live}, Liveness of $\Pibc$.} \\
\addlinespace
Unchanged fork choice
  & Definition~\ref{def:fork-choice} & every node \\
\multicolumn{3}{@{\quad}p{0.95\linewidth}@{}}{
  Used by: Proposition~\ref{prop:bc-live}, Liveness of $\Pibc$; no proved
  result.} \\
\addlinespace
Update rule $\updatefin$
  & Definition~\ref{def:fin} & every node, as local state \\
\multicolumn{3}{@{\quad}p{0.95\linewidth}@{}}{
  Used by: Lemma~\ref{lem:fin-linear}, Local finalization linearity;
  Lemma~\ref{lem:fin-depth}, Local fin-depth; Lemma~\ref{lem:ba-depth}, Local
  ba-depth; Proposition~\ref{prop:lag}, Finalisation lag;
  Proposition~\ref{prop:af-single}, The single-node case;
  Theorem~\ref{thm:pa-af}, Prefix Agreement implies Assured Finality;
  Corollary~\ref{cor:bft-subversion}, Safety under subversion of the BFT
  protocol; Theorem~\ref{thm:fa-af}, Final Agreement implies Assured Finality;
  Proposition~\ref{prop:fa-needs-lin}, Necessity of the Linearity rule;
  Theorem~\ref{thm:pc-ba}, Prefix Consistency implies Prefix Consistency of
  $\ba$; Proposition~\ref{prop:ba-keeps-fin}, Retention of the finalised prefix
  in $\ba$; Corollary~\ref{cor:ba-agree}, Agreement on $\ba$ from Prefix
  Agreement;
  Proposition~\ref{prop:bounded-run}, Bounded non-stalled run;
  Proposition~\ref{prop:short-L}, Necessity of the constraint $\sigma \le L$;
  Proposition~\ref{prop:sticky-no-hazard}, Absence of finalisation hazards
  under sticky fork choice; Proposition~\ref{prop:sl-fa}, Final Agreement
  route under the Shielded Labs rules; Proposition~\ref{prop:sl-pa}, Prefix
  Agreement route under the Shielded Labs rules;
  Proposition~\ref{prop:sl-unbounded}, Unbounded run under the Shielded Labs
  rules; Proposition~\ref{prop:coalition}, Power of a finaliser coalition;
  Proposition~\ref{prop:path-final}, Path to finality;
  Proposition~\ref{prop:final-guarantee}, Guarantee of finality;
  Proposition~\ref{prop:status-tri}, Finality status;
  Proposition~\ref{prop:fin-valid}, Validity of the finalised chain;
  Corollary~\ref{cor:final-confirmations}, Confirmations of a newly final
  block; Lemma~\ref{lem:expired-fin}, Expired unmined below $\fin$.} \\
\addlinespace
Syncing gate, a recommendation
  & Remark~\ref{rem:sync} & every node \\
\multicolumn{3}{@{\quad}p{0.95\linewidth}@{}}{Used by:
  Definition~\ref{def:views}, Application views; no proved result.} \\
\bottomrule
\end{tabular}
\caption[]{Rules, continued.}
\end{table}

\end{document}
